% Organisation des données en informatique — Informatique 3ème — Cours signé OnBuch+
% Programme officiel MINESEC. Compiler avec : tectonic N04-organisation-donnees-informatique.tex
\newcommand{\DOCMATIERE}{Informatique}
\newcommand{\DOCNIVEAU}{3ème}
\newcommand{\DOCMODULE}{Informatique – classe de 3ème}
\newcommand{\DOCLECON}{N04}
\newcommand{\DOCTITRE}{Organisation des données en informatique}
% OnBuch+ — préambule « cours signés » ENRICHI (compilé avec tectonic / XeLaTeX)
% Système visuel « Pop » d'OnBuch+ : fond crème (jamais blanc), bordures encre
% épaisses, ombres « dures » décalées, titres Archivo Black en capitales.
% Version MATHÉMATIQUES : graphes (pgfplots/tikz), boîtes pédagogiques
% (définition, propriété, méthode, attention, le savais-tu, exercice résolu,
% exercice, TP, bilan), page de garde et signature OnBuch+ ancrée en bas.
%
% Macros attendues dans main.tex AVANT \input{preamble} :
%   \DOCMATIERE  \DOCNIVEAU  \DOCMODULE  \DOCLECON (numéro)  \DOCTITRE
\documentclass[11pt]{article}

\usepackage{fontspec}
\usepackage[french]{babel}
\usepackage[a4paper,margin=2cm,top=3.4cm,bottom=2.8cm,headheight=1.4cm,footskip=1.3cm]{geometry}
\usepackage{amsmath,amssymb,amsfonts}
\usepackage{mathtools} % \xmapsto, \coloneqq... souvent utilisés par les modèles
\usepackage{cancel} % simplifications barrées (bilans de forces)
% \abs{z} (module, valeur absolue) et \norm{u} (norme), utilisés par les modèles
\providecommand{\abs}{}\let\abs\relax\DeclarePairedDelimiter{\abs}{\lvert}{\rvert}
\providecommand{\norm}{}\let\norm\relax\DeclarePairedDelimiter{\norm}{\lVert}{\rVert}
\usepackage[table]{xcolor} % [table] : fournit aussi \rowcolors (alternance de lignes)
\usepackage{pagecolor}
\usepackage{tikz}
\usetikzlibrary{arrows.meta,positioning,calc,shapes.geometric,shapes.misc,shapes.arrows,shapes.symbols,decorations.pathmorphing,decorations.pathreplacing,decorations.markings,patterns,shadows,mindmap,trees,fit,backgrounds,matrix,intersections,angles,quotes}
\usepackage{pgfplots}
\usepackage[version=4]{mhchem} % équations nucléaires \ce{^{235}_{92}U}
\usepackage[european,siunitx]{circuitikz} % schémas électriques
\pgfplotsset{compat=1.18}
\usepgfplotslibrary{fillbetween,groupplots} % aires entre courbes (calcul intégral)
\usepackage{enumitem}
\usepackage{fancyhdr}
% [most] charge toutes les bibliothèques tcolorbox (listings, minted, raster,
% documentation...) et gonfle inutilement la mémoire TeX sur un document long
% (~300 boîtes) : on ne charge que ce qui est réellement utilisé.
\usepackage[skins,breakable]{tcolorbox}
\usepackage{graphicx}
\usepackage{colortbl}
\usepackage{booktabs}
\usepackage{tabularx}
\usepackage{diagbox} % case d'en-tête diagonale des tableaux à double entrée
% Colonnes extensibles centrée / gauche / droite (C, L, R), souvent utilisées par les modèles
\newcolumntype{C}{>{\centering\arraybackslash}X}
\newcolumntype{L}{>{\raggedright\arraybackslash}X}
\newcolumntype{R}{>{\raggedleft\arraybackslash}X}
\usepackage{array}
\usepackage{multirow}
\usepackage{multicol}
\usepackage{siunitx}
% parse-numbers=false : accepte aussi \SI{2,5 \times 10^{-3}}{...}
\sisetup{locale=FR,output-decimal-marker={,},group-minimum-digits=4,parse-numbers=false}
\DeclareSIUnit\ohm{\ensuremath{\symup{\Omega}}} % U+2126 absent de la police
\DeclareSIUnit\division{div} % réglages d'oscilloscope (ms/div, V/div)
\DeclareSIUnit\tr{tr} % tours (vitesse de rotation, ex. \SI{1500}{\tr\per\minute})
\DeclareSIUnit\bit{bit}
\DeclareSIUnit\byte{o} % octet
\DeclareSIUnit\inch{po} % pouce
\DeclareSIUnit\ppp{ppp} % points par pouce (résolution d'image)
\usepackage{qrcode}
% Blocs de code source (PHP, C, HTML, JavaScript, SQL) — spécifique à la
% matière Informatique : listings + habillage aux couleurs OnBuch+ (voir le
% style \lstset ci-dessous, à côté des autres boîtes pédagogiques).
\usepackage{listings}
% \degree : commande fréquemment utilisée par les modèles pour un angle ou une
% température en dehors de \SI{}{} (non fournie par les paquets chargés).
\providecommand{\degree}{^{\circ}}
% \qed (fin de démonstration), utilisé par les modèles sans amsthm
\providecommand{\qed}{\hfill$\square$}
% \barre : double barre des valeurs interdites dans un tableau de variations (tkz-tab)
\providecommand{\barre}{\ensuremath{\|}}
% environnement proof (amsthm non chargé) utilisé par les modèles
\ifdefined\proof\else\newenvironment{proof}[1][Démonstration]{\par\noindent\textit{#1.}\ }{\qed\par}\fi
\usepackage{lastpage}
\usepackage{hyperref}
\hypersetup{hidelinks}

% ============================================================
% Palette « Pop » OnBuch+ — mêmes hex que la classe P (plus_ui.dart)
% ============================================================
\definecolor{poporange}{HTML}{F59321}
\definecolor{popdark}{HTML}{E07A0C}
\definecolor{popcream}{HTML}{FFF6E8}
\definecolor{popink}{HTML}{1C1714}
\definecolor{poppurple}{HTML}{7A5AE0}
\definecolor{popgreen}{HTML}{2E9E6A}
\definecolor{poppink}{HTML}{E0457B}
\definecolor{popblue}{HTML}{2F7DE1}
\definecolor{popgold}{HTML}{C98A12}
\definecolor{popmuted}{HTML}{8A7F76}
\definecolor{popline}{HTML}{EADFD2}
\definecolor{popgray}{HTML}{8A7F76}
\colorlet{popgrey}{popgray}
% Alias tolérants (le modèle nomme parfois les couleurs différemment)
\colorlet{popwhite}{white}
\colorlet{popblack}{popink}
\colorlet{popred}{poppink}
\colorlet{popyellow}{popgold}
% Teintes claires pour fonds de tableaux / zones
\colorlet{popblueL}{popblue!10!white}
\colorlet{poporangeL}{poporange!14!white}
\colorlet{popgreenL}{popgreen!12!white}
\colorlet{poppurpleL}{poppurple!12!white}
\colorlet{poppinkL}{poppink!12!white}
\colorlet{popgoldL}{popgold!14!white}

\pagecolor{popcream}

% ============================================================
% Typographie
% ============================================================
\defaultfontfeatures{Ligatures=TeX}
\setmainfont{PlusJakartaSans-Medium.ttf}[Path=../../../fonts/,BoldFont=PlusJakartaSans-Bold.ttf]
% Maths en sans-serif (Fira Math) avec chiffres et lettres droites de Plus
% Jakarta, pour que les formules se fondent dans le texte.
\usepackage[math-style=ISO,bold-style=ISO]{unicode-math}
\setmathfont{FiraMath-Regular.otf}[Path=../../../fonts/]
\setmathfont{PlusJakartaSans-Medium.ttf}[Path=../../../fonts/,range={up/num,up/{Latin,latin},\mathup}]
\usepackage{icomma} % virgule décimale sans espace en mode math
% Caractères mathématiques Unicode absents des polices de texte (Plus Jakarta,
% Archivo Black) : rendus par leur équivalent mathématique, en texte comme en
% formule — sinon ils s'affichent en carré vide (ex. « Arithmétique dans ℤ »).
\usepackage{newunicodechar}
\newunicodechar{ℝ}{\ensuremath{\mathbb{R}}}
\newunicodechar{ℤ}{\ensuremath{\mathbb{Z}}}
\newunicodechar{ℂ}{\ensuremath{\mathbb{C}}}
\newunicodechar{ℕ}{\ensuremath{\mathbb{N}}}
\newunicodechar{ℚ}{\ensuremath{\mathbb{Q}}}
\newunicodechar{𝔻}{\ensuremath{\mathbb{D}}}
\newunicodechar{α}{\ensuremath{\alpha}}
\newunicodechar{β}{\ensuremath{\beta}}
\newunicodechar{γ}{\ensuremath{\gamma}}
\newunicodechar{δ}{\ensuremath{\delta}}
\newunicodechar{ε}{\ensuremath{\varepsilon}}
\newunicodechar{θ}{\ensuremath{\theta}}
\newunicodechar{λ}{\ensuremath{\lambda}}
\newunicodechar{μ}{\ensuremath{\mu}}
\newunicodechar{σ}{\ensuremath{\sigma}}
\newunicodechar{τ}{\ensuremath{\tau}}
\newunicodechar{φ}{\ensuremath{\varphi}}
\newunicodechar{ψ}{\ensuremath{\psi}}
\newunicodechar{ω}{\ensuremath{\omega}}
\newunicodechar{Σ}{\ensuremath{\Sigma}}
% majuscules produites par \MakeUppercase dans les titres (α-aminés -> Α)
\newunicodechar{Α}{\ensuremath{\alpha}}
\newunicodechar{Β}{\ensuremath{\beta}}
\newunicodechar{Γ}{\ensuremath{\Gamma}}
\newunicodechar{Δ}{\ensuremath{\Delta}}
\newunicodechar{Ε}{\ensuremath{\varepsilon}}
\newunicodechar{Θ}{\ensuremath{\Theta}}
\newunicodechar{Λ}{\ensuremath{\Lambda}}
\newunicodechar{Μ}{\ensuremath{\mu}}
\newunicodechar{Τ}{\ensuremath{\tau}}
\newunicodechar{Φ}{\ensuremath{\Phi}}
\newunicodechar{Ψ}{\ensuremath{\Psi}}
\newunicodechar{Ω}{\ensuremath{\Omega}}
\newunicodechar{↦}{\ensuremath{\mapsto}}
\newunicodechar{∈}{\ensuremath{\in}}
\newunicodechar{∉}{\ensuremath{\notin}}
\newunicodechar{∧}{\ensuremath{\wedge}}
\newunicodechar{⇔}{\ensuremath{\Leftrightarrow}}
\newunicodechar{⟺}{\ensuremath{\Longleftrightarrow}}
\newunicodechar{⇒}{\ensuremath{\Rightarrow}}
\newunicodechar{⟹}{\ensuremath{\Longrightarrow}}
\newunicodechar{∘}{\ensuremath{\circ}}
\newunicodechar{∪}{\ensuremath{\cup}}
\newunicodechar{∩}{\ensuremath{\cap}}
\newunicodechar{≡}{\ensuremath{\equiv}}
\newunicodechar{⊂}{\ensuremath{\subset}}
\newunicodechar{⊥}{\ensuremath{\perp}}
\newunicodechar{∃}{\ensuremath{\exists}}
\newunicodechar{∀}{\ensuremath{\forall}}
\newunicodechar{∛}{\ensuremath{\sqrt[3]{\,}}}
\newunicodechar{∜}{\ensuremath{\sqrt[4]{\,}}}
\newunicodechar{⃗}{\ensuremath{{}^{\to}}}
\newunicodechar{ⁿ}{\ifmmode^{n}\else\textsuperscript{\ensuremath{n}}\fi}
\newunicodechar{ˣ}{\ifmmode^{x}\else\textsuperscript{\ensuremath{x}}\fi}
\newunicodechar{ᵇ}{\ifmmode^{b}\else\textsuperscript{\ensuremath{b}}\fi}
\newunicodechar{ʳ}{\ifmmode^{r}\else\textsuperscript{\ensuremath{r}}\fi}
\newunicodechar{ᵘ}{\ifmmode^{u}\else\textsuperscript{\ensuremath{u}}\fi}
\newunicodechar{ʸ}{\ifmmode^{y}\else\textsuperscript{\ensuremath{y}}\fi}
\newunicodechar{ᵃ}{\ifmmode^{a}\else\textsuperscript{\ensuremath{a}}\fi}
\newunicodechar{ᵉ}{\ifmmode^{e}\else\textsuperscript{\ensuremath{e}}\fi}
\newunicodechar{ᵏ}{\ifmmode^{k}\else\textsuperscript{\ensuremath{k}}\fi}
\newunicodechar{ᵖ}{\ifmmode^{p}\else\textsuperscript{\ensuremath{p}}\fi}
\newunicodechar{ᴺ}{\ifmmode^{N}\else\textsuperscript{\ensuremath{N}}\fi}
\newunicodechar{ᶜ}{\ifmmode^{c}\else\textsuperscript{\ensuremath{c}}\fi}
\newunicodechar{ʰ}{\ifmmode^{h}\else\textsuperscript{\ensuremath{h}}\fi}
\newunicodechar{ᵠ}{\ifmmode^{q}\else\textsuperscript{\ensuremath{q}}\fi}
\newunicodechar{ₙ}{\ifmmode_{n}\else\textsubscript{\ensuremath{n}}\fi}
\newunicodechar{ₐ}{\ifmmode_{a}\else\textsubscript{\ensuremath{a}}\fi}
\newunicodechar{ᵢ}{\ifmmode_{i}\else\textsubscript{\ensuremath{i}}\fi}
\newunicodechar{ᵣ}{\ifmmode_{r}\else\textsubscript{\ensuremath{r}}\fi}
\newunicodechar{ₖ}{\ifmmode_{k}\else\textsubscript{\ensuremath{k}}\fi}
\newunicodechar{ᵦ}{\ifmmode_{\beta}\else\textsubscript{\ensuremath{\beta}}\fi}
\newunicodechar{ₓ}{\ifmmode_{x}\else\textsubscript{\ensuremath{x}}\fi}
\newfontfamily\popdisplay{ArchivoBlack-Regular.ttf}[Path=../../../fonts/]
\newfontfamily\poplogo{SpaceGrotesk-Bold.ttf}[Path=../../../fonts/]
\newfontfamily\popsemi{PlusJakartaSans-SemiBold.ttf}[Path=../../../fonts/]
\newfontfamily\popxbold{PlusJakartaSans-ExtraBold.ttf}[Path=../../../fonts/]
\newfontfamily\popxboldls{PlusJakartaSans-ExtraBold.ttf}[Path=../../../fonts/,LetterSpace=3.0]

\setlist[enumerate]{leftmargin=1.7em,itemsep=5pt,topsep=4pt}
\setlist[itemize]{leftmargin=1.7em,itemsep=4pt,topsep=4pt,label={\color{poporange}\textbullet}}
\setlength{\parindent}{0pt}
\setlength{\parskip}{5pt}
\renewcommand{\arraystretch}{1.25}

% pgfplots : style Pop par défaut
\pgfplotsset{
  every axis/.append style={axis line style={popink,line width=1pt},tick style={popink},
    grid style={popline,line width=0.5pt},label style={font=\small},tick label style={font=\footnotesize}},
  popcurve/.style={poporange,line width=1.8pt,no markers,smooth},
}
% Même style disponible dans un \draw TikZ simple (hors pgfplots)
\tikzset{popcurve/.style={poporange,line width=1.8pt}}
\tikzset{popgrid/.style={popline,very thin}}
\colorlet{popcurve}{poporange} % le nom du style est parfois utilisé comme couleur

% ============================================================
% Pill / squiggle
% ============================================================
\newcommand{\poppill}[3]{%
  \tikz[baseline=-0.55ex]\node[draw=popink,line width=1.2pt,rounded corners=2.6pt,%
    fill=#2,inner xsep=6.5pt,inner ysep=3pt]{%
    \popxboldls\fontsize{7}{7}\selectfont\color{#3}\MakeUppercase{#1}};%
}
\newcommand{\popsquiggle}[1]{%
  \tikz[baseline=-0.4ex]{
    \draw[#1,line width=2.6pt,line cap=round]
      plot [smooth] coordinates {(0,0.09) (0.34,0.01) (0.68,0.21) (1.02,0.01) (1.36,0.10)};
  }%
}
% Mot-clé surligné (marqueur orange sous le texte)
\newcommand{\cle}[1]{{\popxbold\color{popdark}#1}}

% ============================================================
% En-tête / pied
% ============================================================
\pagestyle{fancy}
\fancyhf{}
\renewcommand{\headrulewidth}{0pt}
\renewcommand{\footrulewidth}{0pt}
\lhead{\raisebox{-0.15ex}{\poplogo\fontsize{13}{13}\selectfont\color{popink}OnBuch\color{poporange}+}}
\rhead{\poppill{\DOCMATIERE\ \textbullet\ \DOCNIVEAU\ \textbullet\ Leçon \DOCLECON}{white}{popink}}
\lfoot{\popxbold\fontsize{7.6}{7.6}\selectfont\color{popmuted}\MakeUppercase{Cours signé OnBuch+}\ \textbullet\ {\popsemi\DOCTITRE}}
\rfoot{\tikz[baseline=-0.6ex]\node[rounded corners=8.5pt,fill=popink,minimum height=17pt,minimum width=17pt,inner xsep=5pt,inner sep=0]{\popxbold\fontsize{8}{8}\selectfont\color{popcream}\thepage\,/\,\pageref*{LastPage}};}

% ============================================================
% Titres
% ============================================================
\newcommand{\chaptitre}[2]{%
  \begin{tcolorbox}[enhanced,colback=poporange,colframe=popink,boxrule=2.6pt,arc=11pt,
      drop shadow=popink,left=15pt,right=15pt,top=12pt,bottom=11pt,before skip=3pt,after skip=18pt]
    \poppill{#2}{popink}{popcream}\par\vspace{9pt}
    {\raggedright\hyphenpenalty=10000\exhyphenpenalty=10000\popdisplay\fontsize{22}{24}\selectfont\color{popcream}\MakeUppercase{#1}\par}
  \end{tcolorbox}
}
% Section numérotée : pastille orange avec le numéro + titre Archivo
\newcounter{popsec}
\newcommand{\coursec}[1]{%
  \stepcounter{popsec}\setcounter{popsub}{0}%
  \par\vspace{16pt}\noindent
  \begin{minipage}{\linewidth}
  \tikz[baseline=-0.7ex]\node[draw=popink,line width=1.6pt,fill=poporange,rounded corners=4pt,minimum size=22pt,inner sep=0]{\popdisplay\fontsize{12}{12}\selectfont\color{popcream}\thepopsec};%
  \hspace{8pt}{\popdisplay\fontsize{15}{17}\selectfont\color{popink}\MakeUppercase{#1}}\par
  \vspace{4pt}\noindent{\color{poporange}\rule{36pt}{3.6pt}}
  \end{minipage}\par\nopagebreak\vspace{8pt}
}
\newcounter{popsub}
\newcommand{\courssub}[1]{%
  \stepcounter{popsub}%
  \par\vspace{9pt}\noindent{\popxbold\fontsize{11.5}{13}\selectfont\color{popink}{\color{poporange}\thepopsec.\thepopsub}\hspace{6pt}#1}\par\nopagebreak\vspace{4pt}
}

% ============================================================
% Boîtes pédagogiques — même recette : fond blanc, bordure épaisse colorée,
% ombre dure encre, badge plein. Titre optionnel [#1].
% ============================================================
\tcbset{popbase/.style={breakable,enhanced,colback=white,boxrule=2.3pt,arc=9pt,
  shadow={3pt}{-3pt}{0pt}{popink},left=14pt,right=14pt,top=11pt,bottom=11pt,before skip=11pt,after skip=15pt,
  fontupper=\popsemi\fontsize{10.6}{14.5}\selectfont\color{popink},
  fontlower=\popsemi\fontsize{10.6}{14.5}\selectfont\color{popink}}}
% Coche dessinée (le glyphe ✓ n'existe pas dans les polices du document)
\newcommand{\popcheck}[1]{\tikz[baseline=-0.5ex]\draw[#1,line width=1.6pt,line cap=round,line join=round](-0.13,0.0)--(-0.04,-0.09)--(0.13,0.1);}
\providecommand{\coche}{\popcheck{popgreen}}
\providecommand{\resultat}[1]{\par\smallskip\noindent\fcolorbox{popgreen}{popgreenL}{\parbox{\dimexpr\linewidth-2\fboxsep-2\fboxrule\relax}{\textbf{Résultat.} #1}}\par\smallskip}
\newcommand{\popboxhead}[3]{\poppill{#1}{#2}{white}\ifx\relax#3\relax\else\hspace{6pt}{\popxbold\fontsize{10.6}{12}\selectfont\color{popink}#3}\fi\par\vspace{7pt}}

\newtcolorbox{aretenir}[1][]{popbase,colframe=popgreen,before upper={\popboxhead{À retenir}{popgreen}{#1}}}
\newtcolorbox{exemplebox}[1][]{popbase,colframe=poporange,before upper={\popboxhead{Exemple}{poporange}{#1}}}
\newtcolorbox{definition}[1][]{popbase,colframe=popblue,before upper={\popboxhead{Définition}{popblue}{#1}}}
\newtcolorbox{propriete}[1][]{popbase,colframe=poppurple,before upper={\popboxhead{Propriété}{poppurple}{#1}}}
\newtcolorbox{methode}[1][]{popbase,colframe=popgold,before upper={\popboxhead{Méthode}{popgold}{#1}}}
\newtcolorbox{attention}[1][]{popbase,colframe=poppink,before upper={\popboxhead{Attention}{poppink}{#1}}}
\newtcolorbox{experience}[1][]{popbase,colframe=popblue,colback=popblueL,before upper={\popboxhead{Expérience}{popblue}{#1}}}
% « Le savais-tu ? » : carte violette pleine (contexte, culture, Cameroun)
\newtcolorbox{savaistu}[1][]{popbase,colback=poppurple,colframe=popink,
  fontupper=\popsemi\fontsize{10.4}{14.2}\selectfont\color{white},
  before upper={\poppill{Le savais-tu\,?}{white}{poppurple}\ifx\relax#1\relax\else\hspace{6pt}{\popxbold\color{white}#1}\fi\par\vspace{7pt}}}
% Exercice résolu : énoncé en haut, \tcblower puis la solution sur fond crème
\newcounter{popexo}\newcounter{popexr}
\newtcolorbox{exoresolu}[1][]{popbase,colframe=poporange,colbacklower=popcream,
  segmentation style={popink,line width=1.2pt,dashed},
  before upper={\stepcounter{popexr}\popboxhead{Exercice résolu \thepopexr}{poporange}{#1}},
  before lower={\poppill{Solution}{popink}{popcream}\par\vspace{6pt}}}
% Exercice (non résolu en place) — la correction est dans la partie Corrigés
\newtcolorbox{exercice}[1][]{popbase,colframe=popink,
  before upper={\stepcounter{popexo}\popboxhead{Exercice \thepopexo}{popink}{#1}}}
% Corrigé
\newtcolorbox{corrige}[1][]{popbase,colframe=popgreen,colback=popgreenL,
  before upper={\popboxhead{Corrigé}{popgreen}{#1}}}
% Objectifs / prérequis / bilan
\newtcolorbox{objectifs}{popbase,colframe=popink,colback=poporangeL,
  before upper={\popboxhead{Objectifs de la leçon}{popink}{}}}
\newtcolorbox{prerequis}{popbase,colframe=popink,colback=popblueL,
  before upper={\popboxhead{Prérequis}{popblue}{}}}
\newtcolorbox{bilan}{popbase,colframe=popink,colback=white,boxrule=2.8pt,
  before upper={\popboxhead{Fiche bilan}{poporange}{L'essentiel à connaître pour l'examen}}}
% Figure encadrée avec légende
\newtcolorbox{popfigure}[1][]{enhanced,breakable=false,colback=white,colframe=popline,boxrule=1.4pt,arc=8pt,
  left=10pt,right=10pt,top=10pt,bottom=8pt,before skip=10pt,after skip=12pt,halign=center,
  lower separated=false,halign lower=center,
  fontlower=\popsemi\fontsize{9}{11.5}\selectfont\color{popmuted},
  #1}
\newcommand{\legende}[1]{\par\vspace{6pt}{\popsemi\fontsize{9}{11.5}\selectfont\color{popmuted}#1}}

% ============================================================
% Bloc de code source — \begin{lstlisting}[language=Php]...\end{lstlisting}
% (ou language=C, HTML, JavaScript, SQL ; option omise = pseudo-code brut).
% Même recette visuelle que les figures (\popfigure) : cadre encre épais,
% fond clair (popcream), légende possible juste après via \legende{...}
% comme pour une figure (listings ne sait pas arrondir ses coins : on garde
% un cadre rectangulaire, cohérent avec les autres tableaux du document).
% CONTRAT : les modèles ne doivent utiliser QUE \begin{lstlisting}[...] tel
% quel (pas de \begin{tcblisting}, pas de \verb pour un bloc multi-lignes).
% ============================================================
\lstdefinestyle{poplst}{
  backgroundcolor=\color{popcream},
  basicstyle=\popsemi\fontsize{8.6}{11}\selectfont\color{popink},
  keywordstyle=\bfseries\color{popblue},
  commentstyle=\itshape\color{popmuted},
  stringstyle=\color{popgreen},
  numberstyle=\tiny\color{popmuted},
  identifierstyle=\color{popink},
  frame=l,
  rulecolor=\color{popink},
  framerule=1.6pt,
  framesep=7pt,
  framexleftmargin=7pt,
  framexrightmargin=7pt,
  xleftmargin=2pt,
  xrightmargin=2pt,
  breaklines=true,
  breakatwhitespace=false,
  showstringspaces=false,
  showspaces=false,
  showtabs=false,
  tabsize=2,
  columns=flexible,
  keepspaces=true,
  numbers=none,
  aboveskip=10pt,
  belowskip=10pt,
  captionpos=b,
}
\lstset{style=poplst, language=PHP}
% La distribution TeX embarquée ne fournit pas le fichier de langue
% JavaScript de listings (lstlang*.dat) : on la redéfinit nous-mêmes
% pour que \begin{lstlisting}[language=JavaScript] fonctionne.
\lstdefinelanguage{JavaScript}{
  keywords={break,case,catch,class,const,continue,debugger,default,delete,do,
    else,export,extends,finally,for,function,if,import,in,instanceof,let,new,
    return,static,super,switch,this,throw,try,typeof,var,void,while,with,yield,
    async,await,of,get,set},
  keywordstyle=\bfseries\color{popblue},
  ndkeywords={true,false,null,undefined,NaN,Infinity},
  ndkeywordstyle=\bfseries\color{poporange},
  identifierstyle=\color{popink},
  sensitive=true,
  comment=[l]{//},
  morecomment=[s]{/*}{*/},
  string=[b]",
  morestring=[b]',
  morestring=[b]`,
}
% Même limitation pour CSS et les fichiers de configuration (apacheconf,
% nginx, ini...) : ils ne sont pas fournis. CSS et un style générique de
% type "config" (commentaires # ou ;, pas de mots-clés) couvrent les besoins.
\lstdefinelanguage{CSS}{
  keywords={color,background,background-color,margin,padding,border,
    border-radius,font,font-size,font-weight,font-family,width,height,
    display,position,top,left,right,bottom,float,clear,text-align,
    line-height,list-style,cursor,overflow,box-shadow,transition,
    flex,grid,align-items,justify-content,z-index},
  keywordstyle=\bfseries\color{popblue},
  identifierstyle=\color{popink},
  sensitive=false,
  comment=[s]{/*}{*/},
  string=[b]",
  morestring=[b]',
}
\lstdefinelanguage{apacheconf}{
  sensitive=false,
  morecomment=[l]{\#},
}
\lstdefinelanguage{Apache}{
  sensitive=false,
  morecomment=[l]{\#},
}
\lstdefinelanguage{text}{sensitive=false}
\lstdefinelanguage{powershell}{
  keywords={New-Object,New-SmbShare,Get-Acl,Set-Acl,Get-ChildItem,Set-Content,
    Get-Content,Write-Host,ForEach-Object,Where-Object,Select-Object,
    Test-Path,Remove-Item,Copy-Item,Move-Item,Start-Process,Stop-Process,
    If,Else,ElseIf,ForEach,While,Function,Return,Try,Catch},
  keywordstyle=\bfseries\color{popblue},
  identifierstyle=\color{popink},
  sensitive=false,
  comment=[l]{\#},
  string=[b]",
  morestring=[b]',
}
\lstdefinelanguage{ini}{
  sensitive=false,
  morecomment=[l]{;},
  morecomment=[l]{\#},
}

% ============================================================
% Signature OnBuch+
% ============================================================
\newcommand{\poplogoinline}[1]{{\poplogo\fontsize{#1}{#1}\selectfont\color{popink}OnBuch\color{poporange}+}}
\newcommand{\signatureonbuch}{%
  \begin{tcolorbox}[enhanced,colback=popink,colframe=popink,boxrule=2.2pt,arc=10pt,
      drop shadow=poporange,left=14pt,right=14pt,top=10pt,bottom=10pt,before skip=0pt,after skip=0pt]
    \tikz[baseline=-0.7ex]\node[circle,fill=popgreen,draw=popcream,line width=1.3pt,minimum size=22pt,inner sep=0]{\popcheck{white}};%
    \hspace{9pt}\begin{minipage}[c]{0.62\linewidth}
      {\popxbold\fontsize{12}{13}\selectfont\color{popcream}Signé\ \textbullet\ {\poplogo OnBuch\color{poporange}+}}\par\vspace{2pt}
      {\raggedright\popsemi\fontsize{8}{10.5}\selectfont\color{popline}Ludovic A.\\\MakeUppercase{Conforme au programme officiel MINESEC}\par}
    \end{minipage}\hfill
    {\poplogo\fontsize{22}{22}\selectfont\color{popcream}OnBuch\color{poporange}+}
  \end{tcolorbox}%
}

% ============================================================
% Page de garde
% #1 titre · #2 matière · #3 niveau · #4 module · #5 n° leçon · #6 durée ·
% #7 description / objectifs (texte ou itemize)
% ============================================================
\newcommand{\pagegarde}[7]{%
  \thispagestyle{empty}
  \newgeometry{a4paper,margin=1.7cm,top=1.6cm,bottom=1.4cm}%
  \begin{tikzpicture}[remember picture,overlay]
    \fill[poporange!18] ([xshift=-3.2cm,yshift=-3.6cm]current page.north east) circle (4.6cm);
    \fill[poppurple!14] ([xshift=2.4cm,yshift=7.6cm]current page.south west) circle (3.8cm);
  \end{tikzpicture}%
  {\poplogo\fontsize{30}{30}\selectfont\color{popink}OnBuch\color{poporange}+}\hfill
  \raisebox{0.6ex}{\poppill{Cours signé \textbullet\ Premium}{popink}{popcream}}\par
  \vspace{1pt}\popsquiggle{poppurple}\par
  \vspace{8pt}
  \begin{tcolorbox}[enhanced,colback=poporange,colframe=popink,boxrule=2.8pt,arc=13pt,
      drop shadow=popink,left=18pt,right=18pt,top=12pt,bottom=13pt,after skip=10pt]
    \poppill{#2 \textbullet\ #3}{popink}{popcream}\hspace{5pt}\poppill{#4}{white}{popink}\par\vspace{8pt}
    {\popdisplay\fontsize{14}{14}\selectfont\color{popink}LEÇON #5}\par\vspace{4pt}
    {\raggedright\hyphenpenalty=10000\exhyphenpenalty=10000\popdisplay\fontsize{25}{27}\selectfont\color{popcream}\MakeUppercase{#1}\par}
    \vspace{8pt}
    {\popxbold\fontsize{9.5}{9.5}\selectfont\color{popink}\MakeUppercase{Durée conseillée : #6}}
  \end{tcolorbox}
  \begin{tcolorbox}[enhanced,colback=white,colframe=popink,boxrule=2.2pt,arc=10pt,
      drop shadow=popink,left=16pt,right=16pt,top=10pt,bottom=10pt,after skip=8pt]
    \poppill{Dans cette leçon}{poppurple}{white}\par\vspace{5pt}
    \popsemi\fontsize{9.3}{12.6}\selectfont\color{popink}#7
  \end{tcolorbox}
  \noindent\begin{minipage}[t]{0.487\linewidth}
    \begin{tcolorbox}[enhanced,colback=popgreen,colframe=popink,boxrule=2.2pt,arc=10pt,
        drop shadow=popink,left=11pt,right=11pt,top=10pt,bottom=10pt,height=3.1cm,valign=center]
      \tcbox[colback=white,colframe=popink,boxrule=1.3pt,arc=5pt,boxsep=0pt,nobeforeafter,
        left=3pt,right=3pt,top=3pt,bottom=3pt,valign=center]{\qrcode[height=1.9cm]{https://play.google.com/store/apps/details?id=com.luvvix.onbuch&hl=fr}}%
      \hspace{8pt}\begin{minipage}[c]{0.5\linewidth}
        {\popdisplay\fontsize{11}{12.5}\selectfont\color{white}\MakeUppercase{Télécharge OnBuch}}\par\vspace{4pt}
        {\raggedright\popsemi\fontsize{8.4}{11}\selectfont\color{white}Gratuit \textbullet\ Google Play\\Tuteur IA, annales, concours, résultats\par}
      \end{minipage}
    \end{tcolorbox}
  \end{minipage}\hfill
  \begin{minipage}[t]{0.487\linewidth}
    \begin{tcolorbox}[enhanced,colback=poppurple,colframe=popink,boxrule=2.2pt,arc=10pt,
        drop shadow=popink,left=11pt,right=11pt,top=10pt,bottom=10pt,height=3.1cm,valign=center]
      \tikz[baseline=-0.7ex]\node[circle,fill=white,draw=popink,line width=1.3pt,minimum size=1.6cm,inner sep=0]{\popdisplay\fontsize{24}{24}\selectfont\color{poppurple}+};%
      \hspace{8pt}\begin{minipage}[c]{0.6\linewidth}
        {\popdisplay\fontsize{11}{12.5}\selectfont\color{white}\MakeUppercase{Passe à OnBuch+}}\par\vspace{4pt}
        {\raggedright\popsemi\fontsize{8.4}{11}\selectfont\color{white}Tous les cours signés, Tuteur IA illimité, fiches PDF — onglet OnBuch+ de l'app\par}
      \end{minipage}
    \end{tcolorbox}
  \end{minipage}
  \vspace{8pt}
  \popcontenu
  \vfill
  \signatureonbuch
  \restoregeometry
  \newpage
}

% Rangée « Ce cours contient » (page de garde)
\newcommand{\poptile}[3]{% #1 couleur #2 gros texte #3 légende
  \begin{tcolorbox}[enhanced,colback=#1,colframe=popink,boxrule=2pt,arc=9pt,shadow={2.5pt}{-2.5pt}{0pt}{popink},
      left=6pt,right=6pt,top=9pt,bottom=9pt,halign=center,valign=center,height=1.75cm,width=\linewidth,nobeforeafter]
    {\popdisplay\fontsize{12.5}{13}\selectfont\color{white}\MakeUppercase{#2}}\par\vspace{3pt}
    {\centering\popsemi\fontsize{7.8}{9.5}\selectfont\color{white}#3\par}
  \end{tcolorbox}}
\newcommand{\popcontenu}{%
  \par\noindent{\popxboldls\fontsize{7.5}{7.5}\selectfont\color{popmuted}CE COURS SIGNÉ CONTIENT}\par\vspace{7pt}
  \noindent\begin{minipage}[t]{0.235\linewidth}\poptile{popblue}{Cours}{complet, illustré, pas à pas}\end{minipage}\hfill
  \begin{minipage}[t]{0.235\linewidth}\poptile{poporange}{Méthodes}{et exercices résolus}\end{minipage}\hfill
  \begin{minipage}[t]{0.235\linewidth}\poptile{poppink}{Exercices}{type examen + corrigés}\end{minipage}\hfill
  \begin{minipage}[t]{0.235\linewidth}\poptile{popgreen}{Fiche bilan}{l'essentiel à retenir}\end{minipage}\par}

% Sommaire
\newtcolorbox{sommaire}{enhanced,colback=white,colframe=popink,boxrule=2.2pt,arc=10pt,
  drop shadow=popink,left=18pt,right=18pt,top=13pt,bottom=13pt,before skip=6pt,after skip=20pt,
  before upper={\poppill{Sommaire}{poppurple}{white}\par\vspace{9pt}\popsemi\fontsize{10.6}{14.5}\selectfont\color{popink}}}

% Signature de fin de cours — poussée tout en bas de la dernière page
\newcommand{\findecours}{%
  \par\vfill\nopagebreak
  \begin{center}\popsquiggle{poporange}\end{center}\vspace{4pt}
  \signatureonbuch
}
\AtBeginDocument{\def\llbracket{\mathopen{[\mkern-3.5mu[}}\def\rrbracket{\mathclose{]\mkern-3.5mu]}}} % intervalles d'entiers (glyphe ⟦ absent de la police)
% Glyphes absents de Fira Math : redéfinis à partir de glyphes disponibles
\AtBeginDocument{%
  \def\setminus{\mathbin{\backslash}}\let\smallsetminus\setminus
  \def\vdots{\mathord{\vbox{\baselineskip3.2pt\lineskiplimit0pt\kern4pt\hbox{.}\hbox{.}\hbox{.}}}}%
  \def\ddots{\mathinner{\mkern1mu\raise6.5pt\hbox{.}\mkern2mu\raise3.6pt\hbox{.}\mkern2mu\raise0.7pt\hbox{.}\mkern1mu}}%
  \def\triangle{\mathord{\scalebox{0.9}{$\symup{\Delta}$}}}%
  \def\nabla{\mathord{\raisebox{\depth}{\scalebox{1}[-1]{$\symup{\Delta}$}}}}%
  \def\checkmark{\popcheck{popdark}}%
  \def\star{\mathbin{\ast}}\def\bigstar{\mathord{\ast}}%
  \def\diamond{\mathbin{\scalebox{0.6}{$\Diamond$}}}%
  \def\bigcup{\mathop{\mathchoice{\raisebox{-0.3em}{\scalebox{1.7}{$\cup$}}}{\scalebox{1.25}{$\cup$}}{\cup}{\cup}}\displaylimits}%
  \def\bigcap{\mathop{\mathchoice{\raisebox{-0.3em}{\scalebox{1.7}{$\cap$}}}{\scalebox{1.25}{$\cap$}}{\cap}{\cap}}\displaylimits}%
  \def\lesseqqgtr{\mathrel{\lessgtr}}\def\bigoplus{\mathop{\oplus}\displaylimits}%
}
\providecommand{\ssi}{si et seulement si} % abréviation fréquente
\AtBeginDocument{\providecommand{\wideparen}{\overparen}} % arc de cercle

%% FIGFIX-BEGIN v5 — correctifs globaux des figures (docs/onbuch-generation/tools/figfix)
\usepackage{adjustbox}\usepackage{etoolbox}\tcbuselibrary{raster}
% toute figure TikZ/pgfplots plus large que la ligne est réduite à la largeur disponible
\BeforeBeginEnvironment{tikzpicture}{\begin{adjustbox}{max width=\linewidth}}
\AfterEndEnvironment{tikzpicture}{\end{adjustbox}}
% légendes pgfplots lisibles par-dessus les courbes
\pgfplotsset{every axis/.append style={legend style={fill=white,fill opacity=0.92,text opacity=1,draw=black!15,inner sep=3pt}}}
% étiquettes d'arêtes (syntaxe edge["..."]) avec fond blanc
\tikzset{every edge quotes/.append style={fill=white,inner sep=1.2pt}}
%% FIGFIX-END
\begin{document}

\pagegarde{Organisation des données en informatique}{Informatique}{3ème}{Informatique – classe de 3ème}{N04}{3 séances (fiche de progression harmonisée)}{Cette leçon explique comment les données sont organisées de manière structurée (fichiers, dossiers, tableaux) et comment elles sont représentées et codifiées dans l'ordinateur grâce au système binaire et aux codes comme l'ASCII. Elle relie ces notions à des situations concrètes de la vie courante camerounaise.
\vspace{4pt}
\begin{multicols}{2}\small
\begin{itemize}
\item À la fin de cette leçon, je sais définir une donnée, une information et distinguer les principaux types de données (texte, nombre, image, son).
\item À la fin de cette leçon, je sais expliquer ce qu'est une organisation structurée des données (fichiers, dossiers, arborescence, tableau).
\item À la fin de cette leçon, je sais organiser des données de la vie courante dans une arborescence de dossiers et de fichiers bien nommés.
\item À la fin de cette leçon, je sais décrire le rôle de l'unité centrale, de la mémoire et des supports de stockage dans la représentation des données.
\item À la fin de cette leçon, je sais convertir un nombre entier du système décimal vers le binaire et inversement.
\item À la fin de cette leçon, je sais expliquer le principe de la codification des caractères (code ASCII) et coder/décoder un message court.
\end{itemize}\end{multicols}}

\begin{sommaire}
\begin{multicols}{2}
\begin{enumerate}
\item Données, informations et types de données
\item Organisation structurée des données
\item Représentation des données dans l'ordinateur
\item Codification des nombres : le binaire
\item Codification des caractères et des autres données
\item Synthèse et applications à la vie courante
\item Activité d'intégration
\item Méthodes et exercices résolus
\item Exercices
\item Corrigés des exercices
\item Fiche bilan
\end{enumerate}
\end{multicols}
\end{sommaire}

\begin{objectifs}
\begin{itemize}
\item À la fin de cette leçon, je sais définir une donnée, une information et distinguer les principaux types de données (texte, nombre, image, son).
\item À la fin de cette leçon, je sais expliquer ce qu'est une organisation structurée des données (fichiers, dossiers, arborescence, tableau).
\item À la fin de cette leçon, je sais organiser des données de la vie courante dans une arborescence de dossiers et de fichiers bien nommés.
\item À la fin de cette leçon, je sais décrire le rôle de l'unité centrale, de la mémoire et des supports de stockage dans la représentation des données.
\item À la fin de cette leçon, je sais convertir un nombre entier du système décimal vers le binaire et inversement.
\item À la fin de cette leçon, je sais expliquer le principe de la codification des caractères (code ASCII) et coder/décoder un message court.
\item À la fin de cette leçon, je sais citer les unités de mesure de la quantité de données (bit, octet, Ko, Mo, Go) et effectuer des conversions simples.
\item À la fin de cette leçon, je sais rédiger et justifier une démarche de résolution de problème concret utilisant l'organisation et la codification des données.
\end{itemize}
\end{objectifs}

\begin{prerequis}
\begin{itemize}
\item Notion de fichier et de dossier vue en classe de 4ème
\item Manipulation élémentaire du clavier et de la souris
\item Notion de système de numération décimale (mathématiques)
\item Vocabulaire de base : ordinateur, mémoire, clé USB, disque dur
\item Opérations élémentaires sur les puissances de 2 (calcul mental)
\end{itemize}
\end{prerequis}

\newpage

\coursec{Données, informations et types de données}

Au marché Mokolo de Yaoundé, la maman d'Amina tient un cahier où elle note chaque jour ses ventes de tomates : date, quantité, prix. À la fin du mois, elle passe des heures à tout additionner et se trompe souvent. Son neveu, élève en 3ème, lui montre qu'un téléphone ou un ordinateur peut ranger ces mêmes informations de façon ordonnée et faire les calculs en quelques secondes. Mais comment l'ordinateur range-t-il ces données ? Comment transforme-t-il les noms, les nombres et les images en suites de 0 et de 1 ? C'est ce que nous allons découvrir.

\courssub{De la donnée brute à l'information : définitions et distinction}

Avant de parler d'ordinateur, posons les bases. Dans la vie courante, nous sommes entourés de faits, de chiffres, de mots, d'images. Tant qu'ils sont isolés, sans lien ni traitement, ce ne sont que des \textbf{données brutes}. Dès qu'on les organise, qu'on les calcule, qu'on les met en perspective pour prendre une décision, elles deviennent de l'\textbf{information}.

\begin{definition}[Donnée et information]
Une \cle{donnée} (ou \cle{donnée brute}) est un élément élémentaire, non traité, qui représente un fait, une mesure, un symbole ou une observation brute (ex. : le chiffre \texttt{18}, le mot \texttt{« Yaoundé »}, un pixel rouge). Elle n'a pas de sens exploitable par elle-même.

Une \cle{information} est une donnée qui a été \textbf{traitées}, \textbf{organisée} et \textbf{contextualisée} pour avoir un sens précis et utile pour un destinataire (ex. : « La température à Yaoundé est de 18\,°C ce matin »). L'information permet de prendre une décision ou de comprendre une situation.
\end{definition}

\noindent
Le passage de la donnée à l'information, c'est le cœur du travail de l'informaticien : \textbf{le traitement de l'information}.

\begin{exemplebox}[Exemple camerounais : pluviométrie à Douala]
Imaginons une station météo à Douala. Chaque jour à 6\,h, un capteur relève la hauteur de pluie : \texttt{12,5}, \texttt{0}, \texttt{23,1}, \texttt{5,0}, \texttt{0}, \texttt{0}, \texttt{45,2} (en millimètres).
\begin{itemize}
    \item Ces sept nombres sont des \textbf{données brutes}. Seuls, ils ne disent pas s'il a beaucoup plu cette semaine.
    \item Le météorologue les \textbf{traite} : il fait la somme (85,8\,mm), la moyenne (12,3\,mm/jour), repère le maximum (45,2\,mm le dimanche). Il rédige : « Semaine très arrosée, cumuls supérieurs à la normale saisonnière, risque d'inondations dans les bas-fonds. »
    \item Ce bulletin est une \textbf{information} : la mairie peut décider d'envoyer des pompes, les habitants de ranger leurs affaires.
\end{itemize}
\end{exemplebox}

\begin{popfigure}
\begin{tikzpicture}[
    node distance=1.5cm and 2cm,
    box/.style={rounded rectangle, draw=popblue, thick, fill=popblueL, minimum width=3.5cm, minimum height=1.2cm, align=center, font=\small},
    arrow/.style={-Stealth, thick, draw=popdark}
]
\node[box, align=center] (donnees){Données brutes\\ (chiffres, symboles,\\ faits isolés)};
\node[box, right=of donnees, align=center] (traitement){Traitement\\ (tri, calcul, classement,\\ mise en contexte)};
\node[box, right=of traitement, align=center] (info){Information\\ (sens, décision,\\ connaissance)};
\draw[arrow] (donnees) -- node[above, font=\small\bfseries] {Saisie} (traitement);
\draw[arrow] (traitement) -- node[above, font=\small\bfseries] {Transformation} (info);
\end{tikzpicture}
\legende{Chaîne de transformation : de la donnée brute à l'information}
\end{popfigure}

\noindent
Pour bien fixer la distinction, comparons systématiquement données brutes et informations dans un tableau d'exemples locaux :

\begin{center}
\begin{tabularx}{\linewidth}{|>{\bfseries}l|X|X|}
\hline
\rowcolor{popblueL}
\textbf{Situation} & \textbf{Donnée brute (exemple)} & \textbf{Information obtenue} \\
\hline
Bulletin scolaire & Notes : 12, 14, 9, 16, 11 & « Moyenne 12,4 — Admis au BEPC, mention Assez Bien » \\
\hline
Gestion boutique & Ventes du jour : 5\,kg, 2\,kg, 8\,kg, 3\,kg & « Total 18\,kg vendus → Chiffre d'affaires 45\,000\,FCFA, stock restant 12\,kg » \\
\hline
Carnet de santé & Températures : 38,2 ; 38,5 ; 37,8 ; 37,0 & « Fièvre en baisse, traitement efficace, fin dans 2 jours » \\
\hline
Relevé téléphonique & Appels : 06xx… (12 min), 06yy… (3 min) & « Forfait dépassé de 15 min → Facture majorée de 1\,200\,FCFA » \\
\hline
Photo d'identité & Pixels RVB : (255,200,180), (254,199,179)… & « Photo conforme aux normes CNI : fond clair, visage visible » \\
\hline
\end{tabularx}
\end{center}

\begin{attention}[Erreur fréquente : confondre donnée brute et information]
Beaucoup d'élèves (et d'adultes !) disent : « J'ai une information : le numéro de téléphone de Paul est 6\,XX\,XX\,XX\,XX. »
\textbf{Faux.} Ce numéro, tant qu'il n'est pas utilisé pour appeler, situer géographiquement ou identifier un abonné, n'est qu'une \textbf{donnée brute} (une suite de chiffres). Il devient \textbf{information} quand on l'exploite : « Je compose le 6\,XX\,XX\,XX\,XX pour joindre Paul » ou « Ce numéro commence par 6, c'est un mobile MTN Cameroun ».
\emph{Astuce :} Demandez-vous toujours : « Est-ce que cela me permet d'agir ou de comprendre quelque chose maintenant ? » Si oui, c'est de l'information.
\end{attention}

\courssub{Les principaux types de données}

L'ordinateur ne « comprend » pas le sens des mots ni des images. Il ne manipule que des \textbf{nombres binaires} (des 0 et des 1). Pourtant, il gère votre texte, vos photos, vos musiques, vos vidéos. Comment ? En \textbf{codifiant} chaque type de donnée selon une convention précise. Au programme de 3ème, nous identifions cinq grands types de données élémentaires.

\begin{definition}[Types de données élémentaires]
\begin{enumerate}
    \item \cle{Texte} : suite de caractères (lettres, chiffres, signes de ponctuation). Ex. : un SMS, un mot de passe, un article de journal.
    \item \cle{Nombre} : valeurs numériques (entiers, décimaux, négatifs). Ex. : une note sur 20, un prix en FCFA, une température.
    \item \cle{Image} : ensemble de points (pixels) codés par leur couleur. Ex. : une photo de classe, un logo, une radiographie.
    \item \cle{Son} : signal audio numérisé (échantillons). Ex. : un message vocal WhatsApp, la sonnerie du téléphone, l'hymne national.
    \item \cle{Vidéo} : séquence d'images (frames) synchronisée avec une piste son. Ex. : un film Nollywood, un tutoriel YouTube, une visioconférence.
\end{enumerate}
Chaque type a sa propre \textbf{structure} et sa propre \textbf{codification} (que nous verrons en sections 4 et 5).
\end{definition}

\noindent
Reprenons nos exemples de la vie courante pour les classifier :

\begin{exemplebox}[Données du quotidien classées par type]
\begin{itemize}
    \item \textbf{Notes d'un élève} (12, 14, 9, 16) $\rightarrow$ \textbf{Nombre} (entiers ou décimaux).
    \item \textbf{Numéro de téléphone} (6\,XX\,XX\,XX\,XX) $\rightarrow$ \textbf{Texte} (suite de caractères chiffrés, on ne fait pas de calculs dessus).
    \item \textbf{Photo de famille} $\rightarrow$ \textbf{Image} (matrice de pixels).
    \item \textbf{Message vocal WhatsApp} (1 minute) $\rightarrow$ \textbf{Son} (fichier audio compressé, ex. Opus/OGG).
    \item \textbf{Cérémonie de remise des prix filmée} $\rightarrow$ \textbf{Vidéo} (suite d'images + son).
\end{itemize}
\emph{Remarque :} Un même objet peut contenir plusieurs types. Une page web = texte + images + vidéos + sons. Un fichier tableur = nombres + texte + formules (qui sont du texte exécutable).
\end{exemplebox}

\begin{exoresolu}[Calcul de la moyenne d'une classe : du brut à l'info]
\textbf{Énoncé.} Un professeur de 3ème a relevé les notes sur 20 de 6 élèves à un devoir d'informatique : 14 ; 9 ; 17 ; 11 ; 13 ; 16. Il veut obtenir l'information « Moyenne de la classe » pour le bulletin.
\textbf{Données brutes (entrée) :} Liste des 6 nombres.
\textbf{Traitement (algorithme) :}
\begin{lstlisting}[language=pascal]
Algorithme CalculMoyenne
Variables
    notes : Tableau[1..6] d'Entier
    somme, effectif : Entier
    moyenne : Reel
Debut
    notes <- [14, 9, 17, 11, 13, 16]
    somme <- 0
    Pour i de 1 a 6 faire
        somme <- somme + notes[i]
    FinPour
    effectif <- 6
    moyenne <- somme / effectif
    Afficher "Moyenne de la classe : ", moyenne
Fin
\end{lstlisting}
\textbf{Information (sortie) :} « La moyenne de la classe est de 13,33/20. »
\textbf{Interprétation :} Niveau global satisfaisant, pas de redoublement prévisible.
\tcblower
\textbf{Solution détaillée.}
On identifie clairement les étapes :
\begin{itemize}
    \item \textbf{Saisie :} Les 6 notes sont entrées au clavier ou scannées.
    \item \textbf{Stockage :} Elles sont placées en mémoire (variables ou tableau).
    \item \textbf{Calcul :} L'unité arithmétique et logique (UAL) du processeur effectue l'addition et la division.
    \item \textbf{Restitution :} Le résultat s'affiche à l'écran, s'imprime sur le bulletin, ou s'enregistre dans la base de données de l'établissement.
\end{itemize}
C'est exactement ce que fait un tableur (LibreOffice Calc, Excel) avec la formule \texttt{=MOYENNE(A1:A6)}.
\end{exoresolu}

\begin{savaistu}[Un message vocal WhatsApp d'une minute représente déjà des centaines de milliers de données binaires]
Un son est une onde analogique continue. Pour l'enregistrer, l'ordinateur l'\textbf{échantillonne} : il mesure l'amplitude de l'onde plusieurs milliers de fois par seconde.
\begin{itemize}
    \item Qualité « vocale » (bande étroite) : 8\,000 échantillons/seconde (8 kHz).
    \item Chaque échantillon est codé sur 16 bits (2 octets).
    \item Pour 1 minute (60 s) : $8\,000 \times 16 \times 60 = 7\,680\,000$ bits $\approx$ \textbf{960 kOctets} (presque 1 Mo) \textbf{sans compression}.
\end{itemize}
WhatsApp utilise le codec \textbf{Opus} qui compresse fortement (environ 12--24 kb/s). Votre message d'1 minute pèse donc \textbf{90 à 180 kOctets} seulement. Mais à la base, ce sont bien des \textbf{millions de 0 et de 1} qui voyagent sur la fibre CAMTEL ou la 4G MTN/Orange pour arriver chez votre correspondant !
\end{savaistu}

\courssub{Le cycle de traitement de l'information}

Tout système informatique (ordinateur, smartphone, tablette, serveur) suit le même cycle fondamental, hérité de l'architecture de Von Neumann. Comprendre ce cycle, c'est comprendre comment la maman d'Amina au marché Mokolo peut remplacer son cahier par une application simple.

\begin{methode}[Les quatre étapes du traitement de l'information]
\begin{enumerate}
    \item \textbf{Saisie (Entrée) :} Capture des données brutes via des périphériques d'entrée (clavier, souris, scanner, micro, capteur, caméra, écran tactile). \emph{Ex. : Taper les ventes du jour sur le clavier du téléphone.}
    \item \textbf{Stockage (Mémoire) :} Conservation des données et des programmes en mémoire vive (RAM) pendant le traitement, et en mémoire morte (disque dur, SSD, carte SD, cloud) pour la durée. \emph{Ex. : L'application enregistre chaque ligne dans un fichier local ou sur Google Drive.}
    \item \textbf{Traitement (Unité Centrale / Processeur) :} Exécution des instructions (calculs, tris, recherches, conversions) par le processeur (CPU). C'est l'étape « intelligente ». \emph{Ex. : Somme des quantités, multiplication par le prix unitaire, calcul du total mensuel.}
    \item \textbf{Restitution (Sortie) :} Présentation du résultat (information) à l'utilisateur via des périphériques de sortie (écran, imprimante, haut-parleur, port USB, réseau). \emph{Ex. : Affichage « Total mois : 45\,000 FCFA » + impression du récapitulatif.}
\end{enumerate}
\end{methode}

\begin{popfigure}
\begin{tikzpicture}[
    node distance=1.2cm and 1.5cm,
    io/.style={rectangle, draw=popgreen, thick, fill=popgreenL, rounded corners=4pt, minimum width=2.8cm, minimum height=1.1cm, align=center, font=\small},
    cpu/.style={rectangle, draw=poporange, thick, fill=poporangeL, rounded corners=4pt, minimum width=3cm, minimum height=1.8cm, align=center, font=\small},
    mem/.style={cylinder, draw=poppurple, thick, fill=poppurpleL, shape border rotate=90, aspect=0.3, minimum width=2.5cm, minimum height=1.5cm, align=center, font=\small},
    arrow/.style={-Stealth, thick, draw=popdark}
]
\node[io, align=center] (entree){Saisie\\ (Clavier, souris,\\ capteurs, micro...)};
\node[mem, right=of entree, align=center] (memoire){Stockage\\ (RAM, Disque,\\ Cloud)};
\node[cpu, right=of memoire, align=center] (processeur){Traitement\\ (Processeur CPU\\ UAL + Unité de contrôle)};
\node[io, right=of processeur, align=center] (sortie){Restitution\\ (Écran, imprimante,\\ haut-parleur, réseau...)};

\draw[arrow] (entree) -- node[above, font=\small] {Données brutes} (memoire);
\draw[arrow] (memoire) -- node[above, font=\small] {Données + Programmes} (processeur);
\draw[arrow] (processeur) -- node[above, font=\small] {Résultats} (memoire);
\draw[arrow] (memoire) -- node[above, font=\small] {Informations} (sortie);

% Boucle rétroaction mémoire-processeur (optionnel, mais réaliste)
\draw[arrow, popmuted, dashed] (processeur.150) |- ++(-1.5,1.2) -| (memoire.30);
\draw[arrow, popmuted, dashed] (memoire.330) |- ++(-1.5,-1.2) -| (processeur.210);
\end{tikzpicture}
\legende{Cycle fondamental : Saisie $\rightarrow$ Stockage $\rightarrow$ Traitement $\rightarrow$ Restitution (Architecture de Von Neumann simplifiée)}
\end{popfigure}

\noindent
Revenons à Amina et sa maman. Grâce à ce cycle :
\begin{itemize}
    \item \textbf{Saisie :} Maman dicte ou tape « 15 kg tomates à 500 FCFA/kg » sur son smartphone.
    \item \textbf{Stockage :} L'application range la ligne dans une base de données locale (SQLite) ou dans le cloud.
    \item \textbf{Traitement :} Le processeur calcule instantanément $15 \times 500 = 7\,500$ FCFA, met à jour le total journalier, le stock, le bénéfice estimé.
    \item \textbf{Restitution :} L'écran affiche un tableau clair, un graphique des ventes de la semaine, et peut envoyer un résumé par WhatsApp à Amina.
\end{itemize}
Fini les erreurs d'addition, les cahiers perdus, les heures de calcul. L'informatique, c'est d'abord \textbf{automatiser le traitement de l'information} pour gagner du temps et de la fiabilité.

\begin{experience}[TP guidé : Observer le cycle sur votre smartphone]
\textbf{Objectif :} Identifier les 4 étapes du cycle sur une action quotidienne.
\textbf{Matériel :} Un smartphone (ou ordinateur) avec une application de notes ou de calculatrice.
\textbf{Manipulation :}
\begin{enumerate}
    \item Ouvrez l'application \textbf{Calculatrice}.
    \item Tapez \texttt{1500 + 2500 + 3200} (Saisie : clavier tactile).
    \item Appuyez sur \texttt{=} (Traitement : le processeur additionne).
    \item Observez l'affichage \texttt{7200} (Restitution : écran).
    \item Fermez l'application, rouvrez-la : le résultat a disparu (Stockage volatile : RAM non persistante).
    \item Refaites le calcul, puis faites une capture d'écran (Stockage persistant : mémoire flash / galerie photos).
\end{enumerate}
\textbf{Observations :} Sans appui sur \texttt{=}, pas de traitement. Sans stockage persistant (capture, note), l'information est perdue à la fermeture.
\textbf{Interprétation :} Toute application informatique, du jeu vidéo au logiciel de paie de la fonction publique, repose sur ce même cycle. La différence n'est que dans la \textbf{complexité du traitement} et le \textbf{volume des données}.
\end{experience}

\begin{aretenir}[Résumé de la section]
\begin{itemize}
    \item Une \textbf{donnée} est un fait brut, sans sens direct (chiffre, mot, pixel, échantillon son).
    \item Une \textbf{information} est une donnée traitée, organisée, qui a du sens et aide à décider.
    \item Le passage de l'une à l'autre se fait par le \textbf{traitement de l'information} (Saisie $\rightarrow$ Stockage $\rightarrow$ Traitement $\rightarrow$ Restitution).
    \item Il existe 5 types de données de base : \textbf{Texte, Nombre, Image, Son, Vidéo}. Chacun a sa codification propre.
    \item \textbf{Piège :} Ne pas confondre le contenu brut (donnée) et son exploitation (information).
\end{itemize}
\end{aretenir}

\coursec{Organisation structurée des données}

\courssub{Pourquoi structurer les données~?}

Dans la section précédente, nous avons distingué la \cle{donnée} brute (un chiffre, un mot, un son) de l'\cle{information} (la donnée mise en contexte et exploitée). Nous avons aussi vu les \cle{types de données} (entiers, réels, booléens, chaînes de caractères). Mais une donnée isolée ne suffit pas~: dans la vie réelle, comme dans l'ordinateur, nous manipulons des \emph{collections} de données. Imagine un cybercafé à Yaoundé qui gère 500 élèves~: sans rangement, retrouver la fiche d'un seul élève prendrait des heures. C'est le problème fondamental de l'\cle{organisation structurée des données}~: \textbf{comment ranger, nommer et relier les données pour les retrouver vite, éviter les pertes et les doublons~?}

\begin{aretenir}[Le triptyque de l'organisation]
Une bonne organisation des données répond à trois impératifs~:
\begin{enumerate}
    \item \textbf{Accessibilité~:} retrouver l'information en quelques clics (pas en cherchant pendant 20 minutes).
    \item \textbf{Intégrité~:} éviter les pertes (fichier effacé par erreur) et les doublons (deux versions différentes de la même facture).
    \item \textbf{Évolutivité~:} pouvoir ajouter, supprimer ou modifier des données sans casser l'ensemble du système.
\end{enumerate}
\end{aretenir}

\courssub{Notion de fichier et de dossier}

\begin{definition}[Fichier]
Un \cle{fichier} est l'unité de base du stockage de l'information sur un support numérique (disque dur, clé USB, SSD). Il est identifié par~:
\begin{itemize}
    \item un \textbf{nom} (ex. \texttt{bulletin\_3eA}),
    \item une \textbf{extension} (suffixe après le point, ex. \texttt{.txt}, \texttt{.jpg}, \texttt{.mp3}, \texttt{.pdf}) qui indique le \textbf{format} des données et donc le logiciel capable de l'ouvrir,
    \item une \textbf{taille} (exprimée en octets : Ko, Mo, Go),
    \item une \textbf{date de création/modification}.
\end{itemize}
\end{definition}

\begin{savaistu}[L'extension, passeport du fichier]
L'extension n'est pas obligatoire pour que le fichier existe, mais elle est \emph{indispensable} pour l'utilisateur et le système d'exploitation. Sous Windows, masquer les extensions est une source d'erreurs fréquentes (ex. \texttt{photo.jpg.exe} qui passe pour une image alors que c'est un programme potentiellement dangereux). \textbf{Astuce~:} dans l'Explorateur de fichiers, onglet \textit{Affichage}, coche \textit{Extensions de noms de fichiers}.
\end{savaistu}

\begin{definition}[Dossier (ou Répertoire)]
Un \cle{dossier} (ou \cle{répertoire}) est un conteneur logique qui regroupe des fichiers et/ou d'autres dossiers (appelés \cle{sous-dossiers}). Il ne contient pas de données «~métier~» (texte, image, son), mais seulement des \textbf{références} vers d'autres objets. Un dossier a un nom, pas d'extension standard, et une taille qui correspond à la somme des tailles de son contenu.
\end{definition}

\courssub{L'arborescence et le chemin d'accès}

Pour organiser les dossiers entre eux, les systèmes de fichiers (NTFS, ext4, FAT32…) utilisent une structure en \textbf{arbre inversé}~: la \cle{racine} en haut, les branches (dossiers) et les feuilles (fichiers) en bas. C'est l'\cle{arborescence}.

\begin{popfigure}
\begin{tikzpicture}[
    grow'=down,
    level distance=1.6cm,
    edge from parent path={(\tikzparentnode.south) -- (\tikzchildnode.north)},
    every node/.style={draw, rounded corners, fill=popcream, thick, inner sep=3pt, font=\small\sffamily},
    level 1/.style={sibling distance=4.2cm, every node/.style={fill=popblueL, text=popdark}},
    level 2/.style={sibling distance=3.2cm, every node/.style={fill=popgreenL, text=popdark}},
    level 3/.style={sibling distance=3.4cm, every node/.style={fill=poporangeL, text=popdark}},
    root/.style={fill=poppurple, text=white, font=\bfseries\small\sffamily}
]
\node[root] {C:\textbackslash{} (Racine)}
    child {node {Cybercafe}
        child {node {Eleves}
            child {node {3eA\_liste.txt}}
            child {node {3eB\_notes.csv}}
        }
        child {node {Factures}
            child {node {janvier\_2026.xlsx}}
            child {node {fevrier\_2026.xlsx}}
        }
        child {node {Photos}
            child {node {inauguration.jpg}}
            child {node {equipe.png}}
        }
    };
\end{tikzpicture}
\legende{Arborescence type d'un poste «~Cybercafé~» : racine C:\textbackslash{}, dossiers principaux, sous-dossiers et fichiers (feuilles de l'arbre).}
\end{popfigure}

\begin{definition}[Chemin d'accès (Path)]
Le \cle{chemin d'accès} (ou \cle{path}) est l'adresse unique d'un fichier ou dossier dans l'arborescence. Il se lit de la racine vers la cible, en séparant les noms par un \textbf{séparateur}~:
\begin{itemize}
    \item \texttt{\textbackslash} (antislash) sous Windows (ex. \texttt{C:\textbackslash{}Ecole\textbackslash{}3eA\textbackslash{}notes.txt}),
    \item \texttt{/} (slash) sous Linux (ex. \texttt{/home/eleve/notes.txt}).
\end{itemize}
On distingue le \textbf{chemin absolu} (complet depuis la racine) du \textbf{chemin relatif} (depuis le dossier courant).
\end{definition}

\begin{exemplebox}[Lecture d'un chemin absolu Windows]
Prenons le chemin \texttt{C:\textbackslash{}Ecole\textbackslash{}3eA\textbackslash{}Devoirs\textbackslash{}DM01\_algo.txt}.
\begin{enumerate}
    \item \texttt{C:\textbackslash{}} : racine du disque système (lecteur C).
    \item \texttt{Ecole} : dossier principal à la racine.
    \item \texttt{3eA} : sous-dossier dans \texttt{Ecole}.
    \item \texttt{Devoirs} : sous-dossier dans \texttt{3eA}.
    \item \texttt{DM01\_algo.txt} : fichier cible (extension \texttt{.txt} $\rightarrow$ fichier texte).
\end{enumerate}
Si tu es \emph{dans} le dossier \texttt{3eA}, le chemin relatif vers le même fichier est simplement \texttt{Devoirs\textbackslash{}DM01\_algo.txt}.
\end{exemplebox}

\courssub{Règles de bon nommage}

Un nommage rigoureux évite les erreurs, les problèmes de compatibilité (Linux/Windows) et les pertes de temps.

\begin{propriete}[Règles d'or pour nommer fichiers et dossiers]
\begin{enumerate}
    \item \textbf{Pas d'espaces~:} utilise le \texttt{\_} (tiret bas ou \emph{underscore}) ou le \texttt{-} (tiret). \textit{Ex.~:} \texttt{liste\_eleves\_3eA.txt} (bon) vs \texttt{liste eleves 3eA.txt} (mauvais).
    \item \textbf{Caractères autorisés~:} lettres (sans accents de préférence), chiffres, \texttt{\_}, \texttt{-}. \textbf{Interdits sous Windows~:} \texttt{\textbackslash} \texttt{/} \texttt{:} \texttt{*} \texttt{?} \texttt{"} \texttt{<} \texttt{>} \texttt{|} (réservés au système).
    \item \textbf{Noms explicites et datés~:} \texttt{facture\_2026-01-15\_boutique\_ngono.pdf} est meilleur que \texttt{doc1.pdf}.
    \item \textbf{Extension cohérente~:} ne renomme pas \texttt{photo.jpg} en \texttt{photo.txt} pour «~l'ouvrir avec le Bloc-notes~»~: cela rend le fichier illisible par les logiciels.
    \item \textbf{Longueur raisonnable~:} vise moins de 50 caractères pour la lisibilité.
\end{enumerate}
\end{propriete}

\begin{attention}[Le piège du «~Bureau~» et de la racine]
\textbf{Erreur classique~:} enregistrer tous ses fichiers directement sur le \textbf{Bureau} (\texttt{Desktop}) ou à la racine \texttt{C:\textbackslash{}}.
\begin{itemize}
    \item \textbf{Conséquences~:} Bureau encombré $\rightarrow$ démarrage de l'ordinateur ralenti ; risque de suppression accidentelle ; impossibilité de retrouver un fichier parmi 300 icônes.
    \item \textbf{Bonne pratique~:} crée un dossier \texttt{MesTravaux\_3e} dans \texttt{Documents}, puis une arborescence par matière/projet \emph{avant} de créer le premier fichier.
\end{itemize}
\end{attention}

\courssub{Méthode : construire une arborescence claire en 4 étapes}

\begin{methode}[Créer son arborescence de travail]
\begin{enumerate}
    \item \textbf{Analyser les besoins~:} quels types de documents~? (Cours, Devoirs, Photos, Projets, Admin). Qui y accède~? (Moi seul, un groupe, le professeur).
    \item \textbf{Dessiner l'arbre sur papier~:} racine $\rightarrow$ 3 à 5 dossiers principaux maximum $\rightarrow$ sous-dossiers (maximum 2 à 3 niveaux de profondeur pour rester simple).
    \item \textbf{Appliquer le nommage~:} noms courts, sans espaces, sans accents, en minuscules (ex. \texttt{cours\_maths}, \texttt{projet\_site\_web}).
    \item \textbf{Créer l'arborescence \emph{avant} le premier fichier~:} clic droit $\rightarrow$ Nouveau $\rightarrow$ Dossier. Puis enregistre \emph{directement} dans le bon dossier (Fichier $\rightarrow$ Enregistrer sous $\rightarrow$ naviguer vers le dossier cible).
\end{enumerate}
\end{methode}

\courssub{Organisation en tableau : lignes et colonnes}

Au-delà des fichiers, les données \emph{à l'intérieur} d'un fichier (tableur, fichier CSV) suivent souvent une structure \textbf{tabulaire}. C'est la modélisation la plus intuitive~: un \textbf{tableau} organise les données en lignes et en colonnes.

\begin{definition}[Enregistrement et Champ]
Dans un tableau structuré~:
\begin{itemize}
    \item Une \textbf{ligne} représente un \cle{enregistrement} : une occurrence complète (ex. \textbf{une} vente, \textbf{un} élève, \textbf{un} produit).
    \item Une \textbf{colonne} représente un \cle{champ} : une propriété unique (ex. Date, Quantité, Prix). Chaque champ a un \textbf{nom} unique et un \textbf{type de données} (entier, réel, texte, date).
\end{itemize}
La première ligne contient généralement les \textbf{en-têtes} (noms des champs).
\end{definition}

\begin{popfigure}
\begin{center}
\begin{tabularx}{\linewidth}{|l|X|X|X|X|}
\hline
\rowcolor{popblueL}
\textbf{Date} & \textbf{Quantité (kg)} & \textbf{Prix unitaire (FCFA/kg)} & \textbf{Montant (FCFA)} & \textbf{Client} \\
\hline
15/01/2026 & 10 & 500 & 5\,000 & Mme Ngo \\
\hline
16/01/2026 & 5 & 550 & 2\,750 & M. Biya \\
\hline
17/01/2026 & 20 & 500 & 10\,000 & Restaurant Le Baobab \\
\hline
18/01/2026 & 3 & 600 & 1\,800 & Étudiant UY1 \\
\hline
\end{tabularx}
\end{center}
\legende{Tableau structuré «~Ventes de tomates~» : 4 enregistrements (lignes), 5 champs (colonnes). Le champ \textit{Montant} est calculé : Quantité $\times$ Prix unitaire.}
\end{popfigure}

\begin{exemplebox}[Analogie : le cahier de ventes du marché]
Mme Ngo, vendeuse de tomates au Marché Central de Douala, tient un cahier. Chaque page = un fichier. Chaque ligne = une vente (enregistrement). Colonnes = Date, Quantité, Prix/kg, Total, Client.
\begin{itemize}
    \item \textbf{Recherche~:} «~Combien de kg vendus le 16/01~?~» $\rightarrow$ parcourir les lignes en filtrant sur la colonne Date.
    \item \textbf{Calcul~:} «~Recette totale~?~» $\rightarrow$ somme de la colonne Montant : $5\,000 + 2\,750 + 10\,000 + 1\,800 = 19\,550$ FCFA.
    \item \textbf{Tri~:} «~Meilleurs clients~?~» $\rightarrow$ tri décroissant sur la colonne Montant.
\end{itemize}
C'est exactement ce que fait un tableur (LibreOffice Calc, Excel)~: \textbf{structurer pour interroger}.
\end{exemplebox}

\courssub{Analogie avec le classement d'une bibliothèque ou d'un magasin}

\begin{itemize}
    \item \textbf{Bibliothèque de l'école~:} rayons (dossiers principaux : \texttt{Romans}, \texttt{Sciences}, \texttt{Histoire}) $\rightarrow$ étagères (sous-dossiers : \texttt{Physique}, \texttt{Chimie}) $\rightarrow$ livres (fichiers). Chaque livre a une \textbf{cote} (chemin d'accès) : \texttt{Sciences/Physique/Mecanique\_Newton}. Sans cote, le livre est «~perdu~» dans la masse.
    \item \textbf{Magasin (boutique «~Tout pour l'École~»)~:} allées (dossiers) $\rightarrow$ rayons (sous-dossiers) $\rightarrow$ produits (fichiers). Le code-barres joue le rôle d'identifiant unique de chaque produit. L'inventaire consiste à interroger l'ensemble des données structurées.
\end{itemize}

\begin{experience}[TP guidé : créer l'arborescence de ton année de 3ème]
\textbf{Objectif~:} mettre en pratique la méthode des 4 étapes pour organiser tes documents scolaires.

\textbf{Matériel~:} poste Windows ou Linux, Explorateur de fichiers.

\textbf{Durée~:} 20 min.

\textbf{Manipulation~:}
\begin{enumerate}
    \item Ouvre l'Explorateur de fichiers (touche \texttt{Win+E} sous Windows). Va dans \texttt{Documents}.
    \item Crée le dossier racine : \texttt{3eme\_2025-2026} (respecte le nommage~!).
    \item À l'intérieur, crée ces 5 dossiers principaux~: \texttt{cours}, \texttt{devoirs}, \texttt{projets}, \texttt{admin}, \texttt{ressources}.
    \item Dans \texttt{cours}, crée un sous-dossier par matière : \texttt{maths}, \texttt{physique}, \texttt{svt}, \texttt{info}, \texttt{francais}, \texttt{anglais}.
    \item Dans \texttt{devoirs}, crée \texttt{ds\_1er\_trimestre}, \texttt{ds\_2eme\_trimestre}, \texttt{devoirs\_maison}.
    \item \textbf{Test~:} ouvre ton traitement de texte. Écris «~Test arborescence~». Enregistre sous \texttt{test\_arborescence.odt} \textbf{directement} dans \texttt{Documents\textbackslash{}3eme\_2025-2026\textbackslash{}cours\textbackslash{}info}.
    \item Vérifie dans l'Explorateur que le fichier est bien là.
\end{enumerate}

\textbf{Observations attendues~:}
\begin{itemize}
    \item L'arborescence est visible, logique, extensible.
    \item Le chemin d'accès du fichier de test est~: \texttt{C:\textbackslash{}Users\textbackslash{}TonNom\textbackslash{}Documents\textbackslash{}3eme\_2025-2026\textbackslash{}cours\textbackslash{}info\textbackslash{}test\_arborescence.odt}.
    \item Si tu cherches «~test~» dans la barre de recherche de l'Explorateur (en étant dans \texttt{3eme\_2025-2026}), tu le trouves instantanément.
\end{itemize}

\textbf{Interprétation~:} tu viens de construire une \textbf{structure de rangement physique} (sur le disque) qui reflète ta \textbf{organisation logique} (matières, types de documents). C'est la base de toute organisation des données en informatique.
\end{experience}

\courssub{Synthèse de la section}

\begin{aretenir}[Ce qu'il faut retenir de l'organisation structurée]
\begin{itemize}
    \item \textbf{Fichier~:} unité de stockage (nom + extension + taille). L'extension indique le programme capable de l'ouvrir.
    \item \textbf{Dossier~:} conteneur logique pour regrouper fichiers et dossiers. Pas d'extension.
    \item \textbf{Arborescence~:} structure en arbre (racine $\rightarrow$ branches $\rightarrow$ feuilles). Un chemin d'accès unique par objet.
    \item \textbf{Chemin d'accès~:} adresse absolue (depuis la racine) ou relative (depuis le dossier courant). Séparateur \texttt{\textbackslash} (Windows) ou \texttt{/} (Linux).
    \item \textbf{Nommage~:} sans espaces, sans accents, explicite, daté (AAAA-MM-JJ), extension correcte.
    \item \textbf{Tableau~:} lignes = enregistrements (objets), colonnes = champs (propriétés typées). Base du tableur.
    \item \textbf{Bonne pratique~:} créer l'arborescence \textbf{avant} de produire des fichiers. Ne jamais tout mettre sur le Bureau ou à la racine.
\end{itemize}
\end{aretenir}

\begin{exoresolu}[Analyse d'un chemin d'accès et correction de nommage]
\textbf{Énoncé~:} un élève a enregistré son devoir d'algorithmique sous le nom \texttt{Mon Devoir Algo Final V2.docx} sur son Bureau. Le chemin affiché est \texttt{C:\textbackslash{}Users\textbackslash{}Kamel\textbackslash{}Desktop\textbackslash{}Mon Devoir Algo Final V2.docx}.
\begin{enumerate}
    \item Identifie la racine, les dossiers, le fichier et l'extension.
    \item Cite 3 erreurs de nommage dans le nom du fichier.
    \item Propose un nom correct selon les règles vues en cours.
    \item Où aurait-il dû l'enregistrer idéalement (propose une arborescence cohérente)~?
\end{enumerate}
\tcblower
\textbf{Solution détaillée~:}
\begin{enumerate}
    \item \textbf{Racine~:} \texttt{C:\textbackslash{}}. \textbf{Dossiers~:} \texttt{Users}, \texttt{Kamel}, \texttt{Desktop}. \textbf{Fichier~:} \texttt{Mon Devoir Algo Final V2}. \textbf{Extension~:} \texttt{.docx} (document Word ou LibreOffice Writer).
    \item \textbf{Erreurs~:} (1) espaces multiples (posent problème dans les commandes et certains logiciels) ; (2) majuscules incohérentes (pas de convention) ; (3) «~V2~» non standard (mieux vaut une date ou un numéro de version propre) ; (4) «~Final~» trompeur (souvent ce n'est pas le dernier~!).
    \item \textbf{Nom correct~:} \texttt{devoir\_algo\_2026-01-20\_v02.docx}. Format de date clair (AAAA-MM-JJ), tirets bas, minuscules, version numérotée.
    \item \textbf{Arborescence idéale~:} \texttt{C:\textbackslash{}Users\textbackslash{}Kamel\textbackslash{}Documents\textbackslash{}3eme\_2025-2026\textbackslash{}devoirs\textbackslash{}devoirs\_maison\textbackslash{}devoir\_algo\_2026-01-20\_v02.docx}.
\end{enumerate}
\end{exoresolu}

\begin{exercice}[Entraînement personnel]
\begin{enumerate}
    \item Dessine l'arborescence (sous forme d'arbre) de ton téléphone portable : racine \texttt{Stockage\_interne} $\rightarrow$ dossiers \texttt{DCIM}, \texttt{Download}, \texttt{WhatsApp}, \texttt{Documents} $\rightarrow$ au moins 2 fichiers réalistes par dossier (ex. \texttt{IMG\_20260115\_1030.jpg}, \texttt{cv\_kamel.pdf}). Indique les extensions.
    \item Donne le chemin d'accès absolu (style Windows) du fichier \texttt{cv\_kamel.pdf} dans ton arborescence.
    \item Tu reçois un fichier nommé \texttt{facture\_mars\_2026.xlsx}. Tu veux le ranger. Décris précisément (clics et dossiers) où tu le places dans ton arborescence \texttt{3eme\_2025-2026} créée dans le TP, et pourquoi.
    \item Un camarade t'envoie un fichier nommé \texttt{photo\_vacances} (sans extension). Tu double-cliques : Windows demande «~Comment voulez-vous ouvrir ce fichier~?~». Explique pourquoi et comment résoudre le problème sans installer de logiciel.
\end{enumerate}
\end{exercice}

\begin{corrige}[Exercice — Organisation structurée]
\begin{enumerate}
    \item \textit{Arbre attendu~:} racine \texttt{Stockage\_interne} $\rightarrow$ \texttt{DCIM} (\texttt{IMG\_20260115\_1030.jpg}, \texttt{VID\_20260114\_1400.mp4}), \texttt{Download} (\texttt{cv\_kamel.pdf}, \texttt{cours\_algo.pdf}), \texttt{WhatsApp} (\texttt{IMG-20260110-WA0001.jpg}, \texttt{IMG-20260111-WA0002.jpg}), \texttt{Documents} (\texttt{lettre\_motivation.odt}, \texttt{notes\_physique.txt}).
    \item \texttt{Stockage\_interne\textbackslash{}Download\textbackslash{}cv\_kamel.pdf} (style Windows). Sous Android, on écrirait plutôt \texttt{/storage/emulated/0/Download/cv\_kamel.pdf} (style Linux).
    \item On le place dans \texttt{Documents\textbackslash{}3eme\_2025-2026\textbackslash{}admin\textbackslash{}facture\_mars\_2026.xlsx}. Le dossier \texttt{admin} regroupe les documents officiels (factures, inscriptions). Le nom du fichier est déjà conforme aux règles de nommage (pas d'espaces, date claire, extension correcte), il n'y a donc rien à renommer.
    \item L'extension manque $\rightarrow$ le système ne connaît pas le format du fichier ni le programme associé. Solution~: clic droit $\rightarrow$ Renommer $\rightarrow$ ajouter l'extension correcte, par exemple \texttt{photo\_vacances.jpg} (si c'est une photo). Puis double-cliquer : le fichier s'ouvre avec la visionneuse d'images.
\end{enumerate}
\end{corrige}

\coursec{Représentation des données dans l'ordinateur}

Dans la section précédente, nous avons appris à \emph{organiser} les données de manière structurée : un tableau pour les élèves d'une classe, un formulaire pour les clients d'une boutique. Mais une question fondamentale reste en suspens : une fois bien organisées, ces données, comment l'ordinateur les \emph{représente}-t-il physiquement ? Comment une machine faite de circuits électriques peut-elle mémoriser le prénom « Amina », la note 15/20 ou une photo de famille ? C'est ce que nous allons découvrir maintenant, en partant du plus petit élément possible jusqu'aux clés USB et disques durs que tu utilises chaque jour.

\courssub{L'ordinateur ne comprend que deux états}

\textbf{Une machine électrique}\par

Ouvre (en imagination !) le boîtier d'un ordinateur : tu y trouverais des millions de minuscules composants électroniques appelés \cle{transistors}. Un transistor fonctionne comme un interrupteur : soit le courant électrique \textbf{passe}, soit il \textbf{ne passe pas}. Il n'existe pas de troisième possibilité.

C'est une contrainte physique majeure : l'ordinateur ne peut distinguer que \textbf{deux états} :
\begin{itemize}
\item \textbf{absence de courant} (interrupteur ouvert, tension nulle) ;
\item \textbf{présence de courant} (interrupteur fermé, tension présente).
\end{itemize}

Toute l'ingéniosité de l'informatique consiste donc à représenter \emph{toutes} les données — nombres, textes, images, sons, vidéos — uniquement avec ces deux états. On leur associe par convention les symboles \cle{0} (absence de courant) et \cle{1} (présence de courant).

\begin{aretenir}[L'idée fondamentale]
Un ordinateur est une machine électrique qui ne connaît que deux états : courant présent ou courant absent. Toute donnée, quelle qu'elle soit, est donc codée à l'aide de deux symboles seulement : \textbf{0} et \textbf{1}. C'est ce qu'on appelle le codage \cle{binaire}.
\end{aretenir}

\begin{savaistu}[Pourquoi deux états et pas dix ?]
Nous comptons en base 10 parce que nous avons dix doigts. Pour une machine, fabriquer un circuit qui distingue dix niveaux de tension différents serait fragile : la moindre perturbation électrique provoquerait des erreurs. Avec seulement deux états très éloignés (« rien » ou « plein courant »), la machine est rapide et très fiable. C'est le choix fait dès les premiers ordinateurs, dans les années 1940, et il n'a jamais changé depuis.
\end{savaistu}

\courssub{Le bit et l'octet : les briques de base}

\textbf{Le bit : la plus petite unité de donnée}\par

\begin{definition}[Bit]
Un \cle{bit} (contraction de l'anglais \emph{binary digit}, « chiffre binaire ») est la \textbf{plus petite unité de donnée} qu'un ordinateur puisse manipuler. Un bit ne peut prendre que \textbf{deux valeurs : 0 ou 1}.
\end{definition}

Un bit seul ne dit pas grand-chose : c'est comme une lampe qui ne peut être qu'allumée ou éteinte. Mais en combinant plusieurs bits, on peut représenter des choses de plus en plus riches :
\begin{itemize}
\item avec 1 bit : 2 possibilités (0 ou 1) ;
\item avec 2 bits : 4 possibilités (00, 01, 10, 11) ;
\item avec 8 bits : $2^8 = 256$ possibilités !
\end{itemize}

\textbf{L'octet : un groupe de 8 bits}\par

\begin{definition}[Octet]
Un \cle{octet} (en anglais \emph{byte}) est un groupe de \textbf{8 bits}. C'est l'unité de base utilisée pour mesurer la quantité de données. Un octet peut prendre $2^8 = 256$ valeurs différentes (de 00000000 à 11111111).
\end{definition}

Pourquoi 8 bits ? Parce que c'est justement ce qu'il faut pour coder \textbf{un caractère} : une lettre, un chiffre, un signe de ponctuation. Avec 256 combinaisons possibles, on a largement de quoi représenter les 26 lettres en majuscules et minuscules, les 10 chiffres, les accents (é, è, à...) et les symboles.

\begin{exemplebox}[Un caractère $\approx$ 1 octet]
Le mot « CAMEROUN » contient 8 caractères. Il occupe donc environ \textbf{8 octets} en mémoire. La phrase « J'aime l'informatique ! » contient 24 caractères (espaces et apostrophe compris), soit environ 24 octets.
\end{exemplebox}

\begin{attention}[Erreur fréquente : confondre bit et octet]
Le \textbf{bit} et l'\textbf{octet} ne sont pas la même chose ! Retiens bien :
\[ 1 \text{ octet} = 8 \text{ bits} \]
Piège classique : un débit de connexion de 8 mégabits par seconde ne permet \emph{pas} de télécharger 8 mégaoctets par seconde, mais seulement 1 mégaoctet par seconde (car $8 \div 8 = 1$). Astuce : en abrégé, le bit s'écrit avec un « b » minuscule (Mb) et l'octet avec un « o » (Mo).
\end{attention}

\courssub{Les unités de mesure : Ko, Mo, Go, To}

Un octet, c'est très petit : un seul caractère. Une photo en contient des millions ! Pour éviter de manipuler des nombres gigantesques, on a défini des unités plus grandes.

\begin{propriete}[Relations entre les unités de stockage (admises)]
Les unités de mesure des données sont liées par des multiplications par 1024 :
\begin{align*}
1 \text{ Ko (kilooctet)} &= 1024 \text{ octets}\\
1 \text{ Mo (mégaoctet)} &= 1024 \text{ Ko}\\
1 \text{ Go (gigaoctet)} &= 1024 \text{ Mo}\\
1 \text{ To (téraoctet)} &= 1024 \text{ Go}
\end{align*}
\end{propriete}

Pourquoi 1024 et pas 1000 ? Parce que l'ordinateur travaille en binaire, et $1024 = 2^{10}$ est la puissance de 2 la plus proche de 1000. C'est un choix pratique des informaticiens.

\begin{popfigure}
\begin{center}
\begin{tikzpicture}[font=\small]
% Boîtes de la hiérarchie
\node[draw=popblue, fill=popblueL, rounded corners, minimum width=1.9cm, minimum height=1.1cm, align=center] (bit) at (0,0) {\textbf{bit}\\ 0 ou 1};
\node[draw=poporange, fill=poporangeL, rounded corners, minimum width=1.9cm, minimum height=1.1cm, align=center] (octet) at (3,0) {\textbf{octet}\\ 8 bits};
\node[draw=popgreen, fill=popgreenL, rounded corners, minimum width=1.9cm, minimum height=1.1cm, align=center] (ko) at (6,0) {\textbf{Ko}\\ 1024 octets};
\node[draw=poppurple, fill=poppurpleL, rounded corners, minimum width=1.9cm, minimum height=1.1cm, align=center] (mo) at (9,0) {\textbf{Mo}\\ 1024 Ko};
\node[draw=poppink, fill=poppinkL, rounded corners, minimum width=1.9cm, minimum height=1.1cm, align=center] (go) at (12,0) {\textbf{Go}\\ 1024 Mo};
\node[draw=popgold, fill=popgoldL, rounded corners, minimum width=1.9cm, minimum height=1.1cm, align=center] (to) at (15,0) {\textbf{To}\\ 1024 Go};
% Flèches
\draw[-{Stealth}, thick, popdark] (bit) -- (octet) node[midway, above, font=\footnotesize] {$\times 8$};
\draw[-{Stealth}, thick, popdark] (octet) -- (ko) node[midway, above, font=\footnotesize] {$\times 1024$};
\draw[-{Stealth}, thick, popdark] (ko) -- (mo) node[midway, above, font=\footnotesize] {$\times 1024$};
\draw[-{Stealth}, thick, popdark] (mo) -- (go) node[midway, above, font=\footnotesize] {$\times 1024$};
\draw[-{Stealth}, thick, popdark] (go) -- (to) node[midway, above, font=\footnotesize] {$\times 1024$};
% Flèches retour (division)
\draw[-{Stealth}, thick, popmuted, dashed] (octet.south) to[bend right=25] node[below, font=\footnotesize] {$\div 8$} (bit.south);
\draw[-{Stealth}, thick, popmuted, dashed] (ko.south) to[bend right=25] node[below, font=\footnotesize] {$\div 1024$} (octet.south);
\draw[-{Stealth}, thick, popmuted, dashed] (mo.south) to[bend right=25] node[below, font=\footnotesize] {$\div 1024$} (ko.south);
\draw[-{Stealth}, thick, popmuted, dashed] (go.south) to[bend right=25] node[below, font=\footnotesize] {$\div 1024$} (mo.south);
\draw[-{Stealth}, thick, popmuted, dashed] (to.south) to[bend right=25] node[below, font=\footnotesize] {$\div 1024$} (go.south);
\end{tikzpicture}
\end{center}
\legende{Hiérarchie des unités de données : on multiplie pour monter vers les grandes unités, on divise pour redescendre.}
\end{popfigure}

\courssub{Convertir entre les unités}

\begin{methode}[Convertir des unités de stockage]
Pour convertir une quantité de données d'une unité à une autre :
\begin{enumerate}
\item \textbf{Repère le sens de la conversion} sur le schéma de la hiérarchie (bit $\to$ octet $\to$ Ko $\to$ Mo $\to$ Go $\to$ To).
\item Si tu passes d'une \textbf{grande unité vers une petite unité} (ex. Go $\to$ Mo), tu \textbf{multiplies} par 1024 à chaque étape.
\item Si tu passes d'une \textbf{petite unité vers une grande unité} (ex. Mo $\to$ Go), tu \textbf{divises} par 1024 à chaque étape.
\item S'il y a plusieurs étapes (ex. Go $\to$ Ko), multiplie ou divise autant de fois qu'il y a d'étapes.
\end{enumerate}
\end{methode}

\begin{exoresolu}[Conversions d'unités de stockage]
\textbf{Énoncé.} (a) La clé USB de Manga a une capacité de 2 Go. Exprime cette capacité en Mo. (b) Un dossier de photos occupe 5120 Mo. Exprime cette taille en Go. (c) Combien d'octets y a-t-il dans 3 Ko ?
\tcblower
\textbf{Solution détaillée.}

(a) On passe d'une grande unité (Go) vers une petite unité (Mo) : on \textbf{multiplie} par 1024.
\[ 2 \text{ Go} = 2 \times 1024 \text{ Mo} = 2048 \text{ Mo} \]

(b) On passe d'une petite unité (Mo) vers une grande unité (Go) : on \textbf{divise} par 1024.
\[ 5120 \text{ Mo} = \frac{5120}{1024} \text{ Go} = 5 \text{ Go} \]

(c) On passe des Ko aux octets : on multiplie par 1024.
\[ 3 \text{ Ko} = 3 \times 1024 \text{ octets} = 3072 \text{ octets} \]

\textbf{Vérification mentale :} quand on descend vers une petite unité, le nombre doit \emph{augmenter} ; quand on monte vers une grande unité, il doit \emph{diminuer}. Si ton résultat va dans le mauvais sens, tu as inversé multiplication et division !
\end{exoresolu}

\courssub{Où sont stockées les données ?}

Les données codées en binaire doivent être \emph{conservées quelque part}. L'ordinateur utilise deux grandes familles de mémoires, aux rôles complémentaires.

\textbf{La mémoire vive (RAM) : la mémoire de travail}\par

La \cle{mémoire vive} ou \cle{RAM} (\emph{Random Access Memory}) est la mémoire dans laquelle l'ordinateur place les programmes et les données \textbf{pendant qu'il travaille}. Quand tu rédiges une lettre dans un traitement de texte, le texte en cours de saisie se trouve dans la RAM.

Sa caractéristique essentielle : elle est \cle{volatile}, c'est-à-dire \textbf{temporaire}. Dès qu'on éteint l'ordinateur (ou qu'il y a une coupure de courant — fréquentes chez nous !), tout ce qui se trouve dans la RAM est \textbf{perdu}. C'est pour cela qu'il faut \emph{enregistrer} régulièrement son travail.

\textbf{Les mémoires de masse : le stockage permanent}\par

Les \cle{mémoires de masse} conservent les données de façon \textbf{permanente}, même lorsque l'appareil est éteint. On y trouve :
\begin{itemize}
\item le \cle{disque dur} de l'ordinateur (souvent 500 Go ou 1 To) ;
\item la \cle{clé USB} (4 Go, 8 Go, 32 Go...) ;
\item la \cle{carte mémoire} des téléphones et appareils photo ;
\item le disque dur externe.
\end{itemize}

\begin{popfigure}
\begin{center}
\begin{tikzpicture}[font=\small]
% RAM
\node[draw=poporange, fill=poporangeL, rounded corners, minimum width=6.2cm, minimum height=3.4cm, align=center] (ram) at (0,0) {};
\node[font=\bfseries, popdark] at (0,1.35) {Mémoire vive (RAM)};
\draw[fill=popcream, draw=poporange, rounded corners] (-1.6,0.55) rectangle (1.6,1.0);
\foreach \x in {-1.2,-0.6,0,0.6,1.2} \draw[poporange, thick] (\x,0.55) -- (\x,0.35);
\node[align=center] at (0,-0.35) {Travail en cours\\ (rapide)};
\node[align=center, poppink, font=\bfseries] at (0,-1.15) {VOLATILE : s'efface\\ à l'extinction};
% Disque dur
\node[draw=popblue, fill=popblueL, rounded corners, minimum width=6.2cm, minimum height=3.4cm, align=center] (dd) at (8,0) {};
\node[font=\bfseries, popdark] at (8,1.35) {Mémoire de masse (disque dur, clé USB)};
\draw[fill=popcream, draw=popblue] (6.8,0.35) circle (0.45);
\draw[fill=popblue, draw=popblue] (6.8,0.35) circle (0.08);
\draw[fill=popcream, draw=popblue, rounded corners] (7.6,0.15) rectangle (9.4,0.6);
\node[align=center] at (8,-0.35) {Fichiers enregistrés\\ (grande capacité)};
\node[align=center, popgreen, font=\bfseries] at (8,-1.15) {PERMANENTE : conserve\\ les données hors tension};
% Flèche
\draw[-{Stealth}, very thick, popdark] (ram.east) -- (dd.west) node[midway, above, font=\footnotesize\bfseries] {enregistrer};
\draw[-{Stealth}, very thick, popdark] (dd.west) ++(0,-0.5) -- ++(-1.6,0) node[midway, below, font=\footnotesize\bfseries] {ouvrir};
\end{tikzpicture}
\end{center}
\legende{La RAM (volatile) contient le travail en cours ; la mémoire de masse (permanente) conserve les fichiers enregistrés.}
\end{popfigure}

\begin{aretenir}[RAM ou mémoire de masse ?]
\begin{center}
\begin{tabularx}{\linewidth}{|l|X|X|}
\hline
\rowcolor{popblueL} \textbf{Critère} & \textbf{RAM} & \textbf{Mémoire de masse} \\
\hline
Durée de conservation & Temporaire (volatile) & Permanente \\
\hline
Rôle & Travail en cours & Stockage des fichiers \\
\hline
Exemples & Barrette mémoire du PC & Disque dur, clé USB, carte mémoire \\
\hline
Capacité courante & 4 à 8 Go & 32 Go à 1 To et plus \\
\hline
\end{tabularx}
\end{center}
\end{aretenir}

\begin{attention}[Coupure de courant = travail perdu]
Tu rédiges depuis une heure un exposé au cybercafé et le courant se coupe : si tu n'as pas \textbf{enregistré} ton fichier sur le disque dur ou une clé USB, tout ce qui était dans la RAM est perdu. Prends l'habitude d'enregistrer ton travail toutes les quelques minutes (raccourci \texttt{Ctrl + S}) !
\end{attention}

\courssub{Ordres de grandeur et applications à la vie courante}

Pour te donner une intuition concrète des tailles de fichiers, voici quelques valeurs typiques :

\begin{center}
\begin{tabularx}{\linewidth}{|l|X|}
\hline
\rowcolor{popblueL} \textbf{Type de donnée} & \textbf{Taille approximative} \\
\hline
Une page de texte & quelques Ko \\
\hline
Une photo prise avec un téléphone & environ 4 Mo \\
\hline
Une chanson (fichier audio) & environ 4 à 5 Mo \\
\hline
Un film & environ 1 à 2 Go (700 Mo à 2 Go) \\
\hline
Un système d'exploitation complet & plusieurs Go \\
\hline
\end{tabularx}
\end{center}

\begin{exemplebox}[Combien de photos sur une clé USB de 8 Go ?]
Une clé USB de 8 Go peut-elle contenir les 2000 photos du mariage de ta cousine, chaque photo occupant environ 4 Mo ?

\textbf{Étape 1 :} convertir la capacité de la clé en Mo (grande unité $\to$ petite unité : on multiplie).
\[ 8 \text{ Go} = 8 \times 1024 = 8192 \text{ Mo} \]

\textbf{Étape 2 :} diviser la capacité par la taille d'une photo.
\[ \frac{8192 \text{ Mo}}{4 \text{ Mo/photo}} = 2048 \text{ photos} \]

\textbf{Conclusion :} la clé peut contenir environ \textbf{2000 photos} de 4 Mo (2048 exactement, en négligeant l'espace pris par le système de fichiers). Le mariage est sauvé !
\end{exemplebox}

\begin{experience}[Observer les tailles de fichiers et la mémoire d'un ordinateur]
\textbf{Objectif :} vérifier concrètement les notions d'octet, de Ko/Mo/Go et distinguer RAM et stockage.

\textbf{Matériel :} un ordinateur de la salle informatique (Windows), une clé USB contenant quelques fichiers (photos, musique, documents).

\textbf{Manipulation :}
\begin{enumerate}
\item Insère la clé USB, ouvre l'explorateur de fichiers et fais un clic droit sur la clé $\to$ \emph{Propriétés}. Note sa capacité totale et l'espace disponible (en Go).
\item Affiche le contenu de la clé en mode \emph{Détails} : la colonne \emph{Taille} indique la taille de chaque fichier en Ko ou Mo. Compare la taille d'un document texte et d'une photo.
\item Fais un clic droit sur \emph{Ce PC} $\to$ \emph{Propriétés} : repère la quantité de \emph{Mémoire installée (RAM)} et compare-la à la capacité du disque dur.
\end{enumerate}

\textbf{Observations attendues :} les fichiers texte sont minuscules (quelques Ko), les photos bien plus grosses (quelques Mo) ; la RAM (ex. 4 Go) est beaucoup plus petite que le disque dur (ex. 500 Go).

\textbf{Interprétation :} la RAM, petite et rapide, sert au travail momentané ; le disque dur, grand et permanent, sert à conserver tous les fichiers. Les tailles affichées confirment les relations $1 \text{ Ko} = 1024$ octets, $1 \text{ Mo} = 1024$ Ko, etc.
\end{experience}

\begin{exercice}[À toi de jouer]
\begin{enumerate}
\item Convertis : (a) 4 Go en Mo ; (b) 3072 Ko en Mo ; (c) 5 Ko en octets ; (d) 16 bits en octets.
\item La carte mémoire de l'appareil photo du club informatique a une capacité de 16 Go. Une photo occupe environ 4 Mo. Combien de photos la carte peut-elle contenir ?
\item Une clé USB de 4 Go contient déjà 3 films de 1 Go chacun. Peut-on encore y copier 200 chansons de 5 Mo ? Justifie ta réponse par un calcul.
\end{enumerate}
\end{exercice}

\begin{corrige}[Exercice « À toi de jouer »]
\begin{enumerate}
\item (a) $4 \times 1024 = 4096$ Mo. (b) $3072 \div 1024 = 3$ Mo. (c) $5 \times 1024 = 5120$ octets. (d) $16 \div 8 = 2$ octets.
\item $16 \text{ Go} = 16 \times 1024 = 16384$ Mo ; $16384 \div 4 = 4096$ photos.
\item Espace occupé : $3 \times 1 = 3$ Go, donc espace libre : $4 - 3 = 1$ Go $= 1024$ Mo. Les 200 chansons occupent $200 \times 5 = 1000$ Mo. Comme $1000 < 1024$, oui, on peut les copier (il restera environ 24 Mo).
\end{enumerate}
\end{corrige}

Nous savons désormais que toute donnée est une suite de 0 et de 1, mesurée en octets, Ko, Mo, Go ou To, et conservée soit temporairement dans la RAM, soit durablement sur une mémoire de masse. Reste à comprendre \emph{comment} ces 0 et ces 1 représentent concrètement un nombre : c'est l'objet de la prochaine section, consacrée au codage binaire des nombres.

\coursec{Codification des nombres : le binaire}

Dans la section précédente, nous avons vu que l'ordinateur ne manipule que des \textbf{signaux électriques} à l'intérieur de ses circuits. Pour représenter des nombres, du texte, des images ou des sons, il faut trouver une convention d'écriture qui ne repose que sur \textbf{deux états distincts}. C'est exactement le rôle du \textbf{système binaire} : c'est le « langage maternel » de la machine, celui que le processeur, la mémoire et les bus comprennent nativement. Avant d'apprendre à convertir des nombres, comprenons pourquoi ce choix s'impose physiquement.

\courssub{Pourquoi le binaire ? Deux états électriques seulement}

Un transistor, composant de base des puces électroniques, fonctionne comme un interrupteur microscopique : il laisse passer le courant (\textbf{état passant}) ou il le bloque (\textbf{état bloqué}). On associe conventionnellement :
\begin{itemize}
    \item l'état \textbf{passant} (présence de tension, ex. +5 V ou +3,3 V) au chiffre \textbf{1} ;
    \item l'état \textbf{bloqué} (absence de tension, 0 V) au chiffre \textbf{0}.
\end{itemize}

\begin{definition}[Système binaire (base 2)]
Le \cle{système binaire} est un système de numération de position qui utilise seulement \textbf{deux symboles} : \texttt{0} et \texttt{1}. Chaque chiffre s'appelle un \cle{bit} (contraction de \emph{binary digit}). La valeur d'un nombre binaire dépend de la \textbf{position} de chaque bit, exactement comme en base 10, mais les puissances sont des puissances de 2 au lieu de puissances de 10.
\end{definition}

\begin{savaistu}[Le bit, unité fondamentale]
Un \textbf{bit} est la plus petite unité d'information en informatique. Un groupe de 8 bits forme un \textbf{octet} (ou \emph{byte}), qui permet de coder 256 valeurs différentes ($2^8 = 256$). C'est l'unité de base pour mesurer la taille des fichiers (Ko, Mo, Go) et la capacité de la mémoire vive (RAM).
\end{savaistu}

\courssub{Rappel : la numération de position en base 10 (décimale)}

Avant d'aborder la base 2, rappelons le mécanisme que nous utilisons tous les jours. En base 10, nous avons 10 symboles (0 à 9). La valeur d'un chiffre dépend de sa position : unités, dizaines, centaines, milliers, etc. Ces positions correspondent aux puissances de 10 :
\[
\dots \quad 10^3 \quad 10^2 \quad 10^1 \quad 10^0
\]
Par exemple, le nombre \textbf{2 431} se décompose ainsi :
\[
2 \times 10^3 \;+\; 4 \times 10^2 \;+\; 3 \times 10^1 \;+\; 1 \times 10^0
= 2000 + 400 + 30 + 1 = 2431
\]

\courssub{La numération de position en base 2 (binaire)}

Le principe est \textbf{identique}, mais la base change : on utilise les puissances de 2.
\[
\dots \quad 2^5 \quad 2^4 \quad 2^3 \quad 2^2 \quad 2^1 \quad 2^0
\]
Soit les valeurs : \dots \quad 32 \quad 16 \quad 8 \quad 4 \quad 2 \quad 1.

\begin{popfigure}
\begin{center}
\begin{tikzpicture}[
    cell/.style={draw=popblue, thick, minimum width=1.8cm, minimum height=1cm, align=center, font=\sffamily\small},
    header/.style={cell, fill=popblue, text=white, font=\sffamily\bfseries\small},
    value/.style={cell, fill=popblueL, font=\ttfamily\small},
    pow/.style={cell, fill=poporangeL, font=\sffamily\small},
    sum/.style={cell, fill=popgreenL, font=\sffamily\small},
    arrow/.style={-Stealth, thick, popdark}
]
% En-têtes
\node[header] (h0) at (0,0) {Poids (Puissance de 2)};
\node[header] (h1) at (1.8,0) {$2^5=32$};
\node[header] (h2) at (3.6,0) {$2^4=16$};
\node[header] (h3) at (5.4,0) {$2^3=8$};
\node[header] (h4) at (7.2,0) {$2^2=4$};
\node[header] (h5) at (9.0,0) {$2^1=2$};
\node[header] (h6) at (10.8,0) {$2^0=1$};

% Ligne des bits de l'exemple 1101 (aligné à droite sur 4 bits -> 001101)
% On affiche 6 bits pour remplir le tableau : 0 0 1 1 0 1
\node[value] (b0) at (0,-1) {Bits du nombre};
\node[value] (b1) at (1.8,-1) {0};
\node[value] (b2) at (3.6,-1) {0};
\node[value] (b3) at (5.4,-1) {1};
\node[value] (b4) at (7.2,-1) {1};
\node[value] (b5) at (9.0,-1) {0};
\node[value] (b6) at (10.8,-1) {1};

% Ligne des produits (bit * poids)
\node[pow] (p0) at (0,-2) {Produit (Bit $\times$ Poids)};
\node[pow] (p1) at (1.8,-2) {$0 \times 32 = 0$};
\node[pow] (p2) at (3.6,-2) {$0 \times 16 = 0$};
\node[pow] (p3) at (5.4,-2) {$1 \times 8 = 8$};
\node[pow] (p4) at (7.2,-2) {$1 \times 4 = 4$};
\node[pow] (p5) at (9.0,-2) {$0 \times 2 = 0$};
\node[pow] (p6) at (10.8,-2) {$1 \times 1 = 1$};

% Ligne somme
\node[sum] (s0) at (0,-3) {Somme (Valeur décimale)};
\node[sum] (s1) at (1.8,-3) {};
\node[sum] (s2) at (3.6,-3) {};
\node[sum] (s3) at (5.4,-3) {};
\node[sum] (s4) at (7.2,-3) {};
\node[sum] (s5) at (9.0,-3) {};
\node[sum, minimum width=5cm] (s6) at (10.8,-3) {$8 + 4 + 0 + 1 = \mathbf{13}$};

% Flèches indicatives
\draw[arrow] (b3.south) -- ++(0,-0.3) -| (p3.north);
\draw[arrow] (b4.south) -- ++(0,-0.3) -| (p4.north);
\draw[arrow] (b6.south) -- ++(0,-0.3) -| (p6.north);

\node[align=center, font=\sffamily\small, text=popdark] at (5.4, -4.2) {
    \textbf{Exemple :} $1101_2 = 1\times2^3 + 1\times2^2 + 0\times2^1 + 1\times2^0 = 8+4+0+1 = 13_{10}$
};
\end{tikzpicture}
\end{center}
\legende{Décomposition du nombre binaire 1101 (écrit sur 6 bits : 001101) en somme de puissances de 2. Seuls les bits à 1 contribuent à la somme.}
\end{popfigure}

\courssub{Conversion binaire $\rightarrow$ décimal : la somme pondérée}

La méthode consiste à identifier les positions des bits à \texttt{1}, à écrire la puissance de 2 correspondante, et à faire la somme.

\begin{methode}[Convertir un nombre binaire en décimal]
\begin{enumerate}
    \item Écrire le nombre binaire en alignant le bit de poids fort (le plus à gauche) sur la plus grande puissance de 2 nécessaire.
    \item Pour chaque bit égal à \texttt{1}, noter la valeur de sa puissance de 2 ($2^{\text{position}}$, position commençant à 0 à droite).
    \item Additionner toutes ces valeurs. Le résultat est l'équivalent décimal.
\end{enumerate}
\end{methode}

\begin{exemplebox}[Conversion de $10110_2$ en décimal]
\begin{align*}
\text{Bits :}       &\quad 1 \quad 0 \quad 1 \quad 1 \quad 0 \\
\text{Puissances :} &\quad 2^4 \quad 2^3 \quad 2^2 \quad 2^1 \quad 2^0 \\
\text{Valeurs :}    &\quad 16 \quad  8 \quad  4 \quad  2 \quad  1 \\[4pt]
\text{Somme :}      &\quad 1\times16 \;+\; 0\times8 \;+\; 1\times4 \;+\; 1\times2 \;+\; 0\times1 \\
                    &= 16 + 4 + 2 = \mathbf{22}_{10}
\end{align*}
Donc $10110_2 = 22_{10}$.
\end{exemplebox}

\courssub{Conversion décimal $\rightarrow$ binaire : divisions successives par 2}

C'est l'opération inverse. On divise le nombre décimal par 2 successivement jusqu'à obtenir un quotient nul. Les \textbf{restes} (0 ou 1), lus \textbf{de bas en haut} (du dernier au premier), forment le nombre binaire.

\begin{methode}[Convertir un nombre décimal en binaire (divisions successives)]
\begin{enumerate}
    \item Diviser le nombre décimal par 2. Noter le \textbf{quotient} et le \textbf{reste} (0 ou 1).
    \item Recommencer avec le quotient obtenu, tant que le quotient est $\ge 1$.
    \item Lorsque le quotient devient 0, \textbf{arrêter}.
    \item Lire les restes \textbf{de bas en haut} (du dernier reste obtenu au premier). C'est l'écriture binaire.
\end{enumerate}
\end{methode}

\begin{popfigure}
\begin{center}
\begin{tikzpicture}[
    node distance=0.6cm and 1.2cm,
    box/.style={draw=popblue, thick, rounded corners, minimum width=2.5cm, minimum height=0.9cm, align=center, font=\sffamily\small},
    arr/.style={-Stealth, thick, popdark}
]
% Colonne Quotient
\node[box, fill=popblueL] (q0) {Quotient};
\node[box, below=of q0] (q1) {$13 \div 2 = 6$};
\node[box, below=of q1] (q2) {$6 \div 2 = 3$};
\node[box, below=of q2] (q3) {$3 \div 2 = 1$};
\node[box, below=of q3] (q4) {$1 \div 2 = 0$};

% Colonne Reste
\node[box, fill=poporangeL, right=of q0] (r0) {Reste};
\node[box, right=of q1] (r1) {$1$};
\node[box, right=of q2] (r2) {$0$};
\node[box, right=of q3] (r3) {$1$};
\node[box, right=of q4] (r4) {$1$};

% Flèches verticales
\draw[arr] (q1) -- (q2);
\draw[arr] (q2) -- (q3);
\draw[arr] (q3) -- (q4);
\draw[arr] (r1) -- (r2);
\draw[arr] (r2) -- (r3);
\draw[arr] (r3) -- (r4);

% Flèche de lecture (bas vers haut)
\draw[arr, popgreen, very thick] (r4.west) -- ++(-0.5,0) |- (r1.west) node[midway, left=0.3cm, popgreen, font=\bfseries\sffamily\small] {Lecture $\uparrow$};

% Résultat final
\node[align=center, font=\sffamily\small, text=popdark, below=0.8cm of q4] {
    \textbf{Résultat :} $13_{10} = 1101_2$ \quad (restes lus de bas en haut : 1, 1, 0, 1)
};
\end{tikzpicture}
\end{center}
\legende{Schéma en escalier des divisions successives de 13 par 2. La lecture des restes se fait de bas en haut (flèche verte).}
\end{popfigure}

\begin{attention}[Piège classique : le sens de lecture des restes]
C'est l'erreur \textbf{numéro 1} chez les débutants : lire les restes dans l'ordre où on les a trouvés (de haut en bas).
\begin{itemize}
    \item \textbf{Mauvais :} $13 \to$ restes $1, 0, 1, 1 \to$ binaire $1011_2$ (incorrect, cela vaut 11).
    \item \textbf{Bon :} On lit \textbf{de bas en haut} : le dernier reste (celui de la division $1 \div 2$) est le bit de poids fort (le plus à gauche). Donc $1, 1, 0, 1 \to 1101_2$.
\end{itemize}
\textbf{Astuce :} Le premier reste obtenu est toujours le bit de \textbf{poids faible} (le plus à droite, $2^0$). Le dernier reste est le bit de \textbf{poids fort}.
\end{attention}

\courssub{Table de correspondance : nombres de 0 à 15}

Il est utile de mémoriser les 16 premières valeurs (4 bits) pour gagner du temps lors des exercices ou de la lecture de dumps mémoire.

\begin{center}
\begin{tabularx}{\linewidth}{|c|X|c|X|}
\hline
\rowcolor{popblue} \textbf{Décimal} & \textbf{Binaire (4 bits)} & \textbf{Décimal} & \textbf{Binaire (4 bits)} \\
\hline
0 & 0000 & 8 & 1000 \\
1 & 0001 & 9 & 1001 \\
2 & 0010 & 10 & 1010 \\
3 & 0011 & 11 & 1011 \\
4 & 0100 & 12 & 1100 \\
5 & 0101 & 13 & 1101 \\
6 & 0110 & 14 & 1110 \\
7 & 0111 & 15 & 1111 \\
\hline
\end{tabularx}
\end{center}

\begin{aretenir}[Propriété : nombre de valeurs codables]
Avec $n$ bits, on peut coder exactement $2^n$ valeurs différentes (de $0$ à $2^n - 1$).
\begin{itemize}
    \item 4 bits $\to$ 16 valeurs ($0 \dots 15$)
    \item 8 bits (1 octet) $\to$ 256 valeurs ($0 \dots 255$)
    \item 16 bits $\to$ 65 536 valeurs
\end{itemize}
\end{aretenir}

\courssub{Exemple résolu complet : le nombre 13 sous toutes ses coutures}

\begin{exoresolu}[Codage du nombre 13]
\textbf{Énoncé :} Montrer par les deux méthodes que le nombre décimal 13 s'écrit $1101$ en binaire.

\textbf{Solution détaillée :}

\textbf{Méthode 1 : Puissances de 2 (Binaire $\to$ Décimal) — Vérification}
On cherche la décomposition de 13 en puissances de 2.
\begin{itemize}
    \item $2^3 = 8 \le 13$. On pose un 1 au rang $2^3$. Reste : $13 - 8 = 5$.
    \item $2^2 = 4 \le 5$. On pose un 1 au rang $2^2$. Reste : $5 - 4 = 1$.
    \item $2^1 = 2 > 1$. On pose un 0 au rang $2^1$.
    \item $2^0 = 1 \le 1$. On pose un 1 au rang $2^0$. Reste : $0$.
\end{itemize}
Bits obtenus (poids forts $\to$ poids faibles) : $1\;1\;0\;1$.
Vérification : $1\times8 + 1\times4 + 0\times2 + 1\times1 = 13$. \checkmark

\textbf{Méthode 2 : Divisions successives (Décimal $\to$ Binaire)}
\begin{center}
\begin{tabular}{r|l|l}
Division & Quotient & Reste \\ \hline
$13 \div 2$ & 6 & \textbf{1} (LSB - poids faible) \\
$6 \div 2$  & 3 & \textbf{0} \\
$3 \div 2$  & 1 & \textbf{1} \\
$1 \div 2$  & 0 & \textbf{1} (MSB - poids fort) \\
\end{tabular}
\end{center}
Lecture des restes \textbf{de bas en haut} : $\mathbf{1\,1\,0\,1}$.
Résultat : $13_{10} = 1101_2$.
\end{exoresolu}

\courssub{Complément : le système hexadécimal (base 16) — \emph{hors programme strict, culture informatique}}

Les informaticiens utilisent souvent l'\textbf{hexadécimal} comme raccourci du binaire. La base 16 utilise 16 symboles : \texttt{0-9} puis \texttt{A, B, C, D, E, F} (valeurs 10 à 15).
\begin{itemize}
    \item Un chiffre hexadécimal code \textbf{exactement 4 bits} (un \emph{quartet} ou \emph{nibble}).
    \item Conversion immédiate : on regroupe les bits par paquets de 4 (en partant de la droite) et on remplace chaque paquet par son équivalent hexa.
\end{itemize}

\begin{exemplebox}[Binaire $\leftrightarrow$ Hexadécimal]
$1101\,0110_2$ se groupe en $1101$ et $0110$.
\begin{itemize}
    \item $1101_2 = 13_{10} = \mathtt{D}_{16}$
    \item $0110_2 = 6_{10} = \mathtt{6}_{16}$
\end{itemize}
Donc $11010110_2 = \mathtt{D6}_{16}$.
C'est beaucoup plus compact à écrire et à lire pour un humain, tout en restant trivial à convertir vers/depuis le binaire pour la machine.
\end{exemplebox}

\courssub{Synthèse et applications à la vie courante}

\begin{itemize}
    \item \textbf{Adressage IP (IPv4) :} Une adresse comme \texttt{192.168.1.1} est stockée sur 32 bits (4 octets). $192 = 11000000_2$, $168 = 10101000_2$, etc. Les masques de sous-réseau (ex. \texttt{/24} ou \texttt{255.255.255.0}) sont des opérations sur les bits.
    \item \textbf{Couleurs web (RGB) :} \texttt{\#FF5733} signifie : Rouge=$FF_{16}=255$, Vert=$57_{16}=87$, Bleu=$33_{16}=51$. Chaque composante est codée sur 1 octet (8 bits).
    \item \textbf{Droits fichiers (Linux) :} \texttt{chmod 755} $\to$ $7=111_2$ (lecture+écriture+exécution pour le propriétaire), $5=101_2$ (lecture+exécution pour groupe/autres). C'est de la manipulation de bits (masques).
    \item \textbf{Capteurs et IoT :} Un capteur de température envoie une valeur sur 12 bits (0 à 4095). Le microcontrôleur lit les registres binaires bruts et applique une formule de conversion.
\end{itemize}

\begin{exercice}[Entraînement aux conversions]
\begin{enumerate}
    \item Convertir en décimal : $101_2$, $1111_2$, $10010_2$, $110011_2$.
    \item Convertir en binaire (méthode des divisions) : $5_{10}$, $18_{10}$, $31_{10}$, $63_{10}$.
    \item Compléter le tableau :
    \begin{center}
    \begin{tabular}{|c|c|c|}
    \hline
    Décimal & Binaire (8 bits) & Hexadécimal \\ \hline
    42      &                  &             \\ \hline
          & 10101010         &             \\ \hline
          &                  & \texttt{3F} \\ \hline
    \end{tabular}
    \end{center}
    \item \textbf{Contexte camerounais :} Un cybercafé à Yaoundé configure un routeur. L'adresse IP du serveur est \texttt{192.168.10.25}. Écrire les 4 octets de cette adresse en binaire (8 bits chacun).
\end{enumerate}
\end{exercice}

\begin{corrige}[Exercice n°1]
\begin{enumerate}
    \item $101_2 = 5$ ; $1111_2 = 15$ ; $10010_2 = 18$ ; $110011_2 = 51$.
    \item $5_{10} = 101_2$ ; $18_{10} = 10010_2$ ; $31_{10} = 11111_2$ ; $63_{10} = 111111_2$.
    \item
    \begin{center}
    \begin{tabular}{|c|c|c|}
    \hline
    Décimal & Binaire (8 bits) & Hexadécimal \\ \hline
    42      & 00101010         & \texttt{2A} \\ \hline
    170     & 10101010         & \texttt{AA} \\ \hline
    63      & 00111111         & \texttt{3F} \\ \hline
    \end{tabular}
    \end{center}
    \item $192 = 11000000_2$ ; $168 = 10101000_2$ ; $10 = 00001010_2$ ; $25 = 00011001_2$.
    Adresse complète : \texttt{11000000.10101000.00001010.00011001}.
\end{enumerate}
\end{corrige}

\begin{experience}[TP guidé : « Calculatrice binaire » en algorithmique]
\textbf{Objectif :} Écrire un algorithme qui convertit un entier positif saisi au clavier en sa représentation binaire (sans utiliser de fonction toute faite).

\textbf{Matériel :} Papier, crayon, puis éditeur de code (Python, C, ou pseudo-code).

\textbf{Manipulation :}
\begin{enumerate}
    \item \textbf{Variables :} \texttt{N} (entier, nombre à convertir), \texttt{R} (chaîne de caractères, résultat), \texttt{Q} (entier, quotient), \texttt{rest} (entier, reste).
    \item \textbf{Algorithme :}
    \begin{lstlisting}[language=Pascal]
Algorithme Conversion_Binaire
Variables
    N, Q, rest : Entier
    R : Chaine
Debut
    Afficher("Entier positif a convertir : ")
    Saisir(N)
    Q <- N
    R <- ""
    Si N = 0 Alors
        R <- "0"
    Sinon
        TantQue Q > 0 Faire
            rest <- Q mod 2      // Reste de la division par 2 (0 ou 1)
            R <- Caractere(rest) + R  // Concaténation À GAUCHE (lecture bas->haut)
            Q <- Q div 2         // Division entière
        FinTantQue
    FinSi
    Afficher("Binaire : ", R)
Fin
    \end{lstlisting}
    \item \textbf{Test :} Exécuter avec $N=13$, $N=0$, $N=255$. Vérifier que l'affichage correspond aux tables.
    \item \textbf{Observation :} Pourquoi concatène-t-on \texttt{Caractere(rest) + R} et non \texttt{R + Caractere(rest)} ?
    \item \textbf{Interprétation :} Cela correspond exactement à la lecture des restes « de bas en haut » : le premier reste calculé (poids faible) doit se retrouver à la \textbf{droite} de la chaîne finale.
\end{enumerate}
\end{experience}

\coursec{Codification des caractères et des autres données}

Dans la section précédente, nous avons appris à coder les \cle{nombres} en binaire : l'ordinateur ne comprend que des 0 et des 1, et tout nombre entier peut s'écrire comme une somme de puissances de 2. Mais un ordinateur ne manipule pas seulement des nombres. Quand tu tapes ton nom au clavier, quand tu rédiges une lettre dans un traitement de texte ou quand tu envoies un message à un camarade, l'ordinateur traite des \cle{caractères} : des lettres, des chiffres, des signes de ponctuation. Il traite aussi des images, des sons, des vidéos. Comment ces données, qui ne sont pas des nombres, peuvent-elles être représentées par des 0 et des 1 ? C'est toute la question de cette section, et la réponse est d'une élégance remarquable : \textbf{on associe à chaque donnée un nombre, puis on code ce nombre en binaire}.

\courssub{Le principe : donner un numéro à chaque caractère}

\textbf{Une situation de la vie courante.}\par

Imagine le directeur d'un collège de Yaoundé qui veut préparer la liste des élèves de 3\textsuperscript{ème} sans écrire leurs noms en entier, pour gagner du temps. Il décide d'un code secret : chaque élève reçoit un numéro (Aïcha = 1, Bertrand = 2, Chantal = 3\ldots). Désormais, au lieu d'écrire les noms, il écrit les numéros. Celui qui possède la \textbf{table de correspondance} peut retrouver les noms ; les autres n'y comprennent rien.

L'ordinateur fait exactement la même chose avec les caractères. Comme il ne sait manipuler que des nombres (eux-mêmes écrits en binaire), on a décidé, une fois pour toutes, d'associer un \textbf{numéro unique à chaque caractère} : la lettre A, la lettre a, le chiffre 0, le point, l'espace, etc. Écrire un texte dans la mémoire de l'ordinateur revient alors à écrire une suite de nombres, donc une suite de 0 et de 1.

\begin{definition}[Codification des caractères]
La \cle{codification des caractères} est l'opération qui consiste à associer à chaque caractère (lettre, chiffre, signe de ponctuation, symbole) un \textbf{nombre entier unique}, appelé son \textbf{code}, puis à représenter ce nombre en binaire pour qu'il soit manipulable par l'ordinateur. La table qui donne la correspondance entre les caractères et leurs codes s'appelle une \textbf{table de codage}.
\end{definition}

\begin{aretenir}[L'idée fondamentale]
Pour l'ordinateur, un texte n'est qu'une \textbf{suite de nombres}. Taper la lettre « A » au clavier, c'est en réalité enregistrer en mémoire le nombre 65, qui s'écrit 1000001 en binaire. Tout le reste de cette section n'est que l'application de cette idée.
\end{aretenir}

\courssub{Le code ASCII}

Pour que tous les ordinateurs du monde se comprennent, il fallait une table de codage \textbf{commune et universelle}. C'est ainsi qu'est né, aux États-Unis en 1963, le code ASCII, encore utilisé aujourd'hui.

\begin{definition}[Le code ASCII]
Le code \cle{ASCII} (sigle de \emph{American Standard Code for Information Interchange}, c'est-à-dire « Code américain normalisé pour l'échange d'information ») est une table de codage qui associe un nombre à \textbf{128 caractères} : les 26 lettres majuscules, les 26 lettres minuscules, les 10 chiffres, les signes de ponctuation, quelques symboles et des caractères de contrôle (comme le retour à la ligne). Les codes vont de 0 à 127, ce qui nécessite \textbf{7 bits} ($2^7 = 128$) ; on les range habituellement dans un octet (8 bits) en ajoutant un 0 de poids fort.
\end{definition}

\begin{propriete}[Codes ASCII à connaître]
Quelques repères utiles dans la table ASCII :
\begin{itemize}
\item les \textbf{chiffres} « 0 » à « 9 » ont les codes 48 à 57 ;
\item les \textbf{lettres majuscules} « A » à « Z » ont les codes 65 à 90 ;
\item les \textbf{lettres minuscules} « a » à « z » ont les codes 97 à 122 ;
\item l'\textbf{espace} a le code 32.
\end{itemize}
On remarque que les codes sont rangés dans l'ordre : B suit A, b suit a, « 1 » suit « 0 ». Pour passer d'une majuscule à la minuscule correspondante, on ajoute 32 (par exemple : $65 + 32 = 97$, donc A devient a).
\end{propriete}

\begin{exemplebox}[Trois codes à retenir par cœur]
\begin{itemize}
\item La lettre \textbf{A} (majuscule) a pour code ASCII \textbf{65}, soit 1000001 en binaire (sur 7 bits).
\item La lettre \textbf{a} (minuscule) a pour code ASCII \textbf{97}, soit 1100001 en binaire.
\item Le \textbf{chiffre 0} (le caractère « 0 ») a pour code ASCII \textbf{48}, soit 0110000 en binaire.
\end{itemize}
Attention : le code 48 représente le \emph{caractère} « 0 », pas le \emph{nombre} zéro !
\end{exemplebox}

\begin{popfigure}
\begin{center}
\renewcommand{\arraystretch}{1.2}
\begin{tabular}{|c|c|c|c|c|c|}
\hline
\rowcolor{popblueL} \textbf{Caractère} & \textbf{Code décimal} & \textbf{Code binaire (7 bits)} & \textbf{Caractère} & \textbf{Code décimal} & \textbf{Code binaire (7 bits)} \\
\hline
espace & 32 & 0100000 & M & 77 & 1001101 \\
\hline
0 & 48 & 0110000 & N & 78 & 1001110 \\
\hline
1 & 49 & 0110001 & O & 79 & 1001111 \\
\hline
3 & 51 & 0110011 & a & 97 & 1100001 \\
\hline
A & 65 & 1000001 & b & 98 & 1100010 \\
\hline
B & 66 & 1000010 & c & 99 & 1100011 \\
\hline
C & 67 & 1000011 & m & 109 & 1101101 \\
\hline
\end{tabular}
\end{center}
\legende{Extrait de la table ASCII : chaque caractère possède un code décimal, que l'ordinateur écrit ensuite en binaire sur 7 bits.}
\end{popfigure}

\courssub{Coder et décoder un mot avec la table ASCII}

\begin{methode}[Coder un mot en ASCII puis en binaire]
Pour coder un mot, par exemple « CAM » :
\begin{enumerate}
\item \textbf{Découper} le mot en caractères : C, A, M.
\item \textbf{Chercher} dans la table ASCII le code décimal de chaque caractère : C = 67, A = 65, M = 77.
\item \textbf{Convertir} chaque code en binaire (par divisions successives par 2, comme on l'a appris dans la section précédente) : 67 = 1000011, 65 = 1000001, 77 = 1001101.
\item \textbf{Écrire} la suite des codes binaires les uns à la suite des autres : c'est le mot tel que l'ordinateur le mémorise.
\end{enumerate}
Pour \textbf{décoder}, on fait le chemin inverse : on découpe la suite binaire en groupes de 7 bits, on convertit chaque groupe en décimal, puis on lit le caractère correspondant dans la table.
\end{methode}

\begin{exemplebox}[Le mot « CAM » codé en ASCII]
Le mot « CAM » se code, en décimal : \textbf{67 -- 65 -- 77}, puis en binaire :
\begin{center}
\textbf{1000011 \quad 1000001 \quad 1001101}
\end{center}
C'est cette suite de 21 bits (trois fois 7 bits) qui est réellement stockée dans la mémoire de l'ordinateur lorsque tu tapes « CAM ».
\end{exemplebox}

\begin{exoresolu}[Coder puis décoder]
\textbf{Énoncé.} (a) À l'aide de la table ASCII simplifiée ci-dessus, coder en décimal le mot « bac ». (b) Décoder la suite de codes : 66 -- 79 -- 78.
\tcblower
\textbf{Solution.}
(a) On découpe le mot : b, a, c. Dans la table : b = 98, a = 97, c = 99. Le mot « bac » se code donc \textbf{98 -- 97 -- 99}. En binaire : 1100010 -- 1100001 -- 1100011.

(b) On cherche chaque code dans la table : 66 correspond à B, 79 correspond à O, 78 correspond à N. La suite 66 -- 79 -- 78 représente donc le mot \textbf{« BON »}.
\end{exoresolu}

\begin{attention}[Ne pas confondre le chiffre « 3 » et le nombre 3]
C'est l'erreur la plus fréquente ! Le \textbf{nombre} 3 (une quantité, avec laquelle on calcule) se code en binaire : $3 = 11_2$. Mais le \textbf{caractère} « 3 » (le symbole que tu tapes au clavier dans un texte) a pour code ASCII \textbf{51}, soit 0110011 en binaire. Ce ne sont pas du tout les mêmes choses : dans « j'ai 3 stylos », le « 3 » est un caractère (code 51) ; dans un calcul $2 + 1 = 3$, le résultat est le nombre 3 (binaire 11). Retiens : \emph{tout ce qui est tapé au clavier est d'abord un caractère}.
\end{attention}

\courssub{Les limites de l'ASCII et l'Unicode}

Le code ASCII a été inventé aux États-Unis pour écrire l'anglais. Or l'anglais s'écrit sans aucun accent ! Avec seulement 128 caractères, impossible d'écrire correctement le \textbf{français} (é, è, à, ç), et encore moins les \textbf{langues camerounaises} qui utilisent des lettres spéciales, ni l'arabe, ni le chinois. Pour répondre à ce problème, on a d'abord créé des tables étendues (256 caractères sur 8 bits), puis une norme universelle : l'\cle{Unicode}.

\begin{savaistu}[L'Unicode : écrire toutes les langues du monde]
L'\textbf{Unicode} est une table de codage universelle qui attribue un numéro à plus de \textbf{140 000 caractères}, couvrant presque toutes les écritures du monde : alphabet latin, arabe, chinois, cyrillique\ldots{} Grâce à l'Unicode, on peut écrire sur ordinateur les lettres propres aux langues camerounaises, comme « \textbf{ə} » (présent en ewondo et en fulfuldé) ou « \textbf{ŋ} » (utilisé en ghomala' et en médumba). Ainsi, un élève de Bafoussam ou de Maroua peut écrire correctement son nom de famille ou un texte dans sa langue sur un ordinateur. Les premiers caractères de l'Unicode sont exactement ceux de l'ASCII : la lettre A garde le code 65, ce qui assure la compatibilité avec les anciens documents.
\end{savaistu}

\courssub{Codification des images et des sons : l'idée générale}

Le principe « associer un nombre à chaque élément » ne s'arrête pas au texte. Il s'applique à toutes les données.

\textbf{Les images.} Une image numérique est découpée en une grille de minuscules carrés appelés \cle{pixels} (contraction de \emph{picture elements}, « éléments d'image »). À chaque pixel, on associe un nombre qui indique sa couleur : par exemple 0 pour un pixel noir et 1 pour un pixel blanc dans une image en noir et blanc. Une image devient ainsi un grand tableau de nombres, donc une suite de 0 et de 1. Plus il y a de pixels, plus l'image est précise : c'est ce qu'on appelle la \textbf{résolution}. Pour les images en couleur, on utilise généralement trois nombres par pixel (rouge, vert, bleu -- codage RVB).

\begin{popfigure}
\begin{center}
\begin{tikzpicture}[scale=0.55]
% Grille de pixels formant un "C"
\foreach \x/\y in {1/4,2/4,3/4, 0/3, 0/2, 0/1, 1/0,2/0,3/0}{
  \fill[popblue] (\x,\y) rectangle (\x+1,\y+1);
}
\draw[popdark, thin] (0,0) grid[step=1] (5,5);
\node[popmuted] at (2.5,-0.7) {image agrandie : les pixels};
% Flèche
\draw[-{Stealth}, very thick, poporange] (5.8,2.5) -- (7.6,2.5);
% Codes binaires des lignes
\node[anchor=west, font=\ttfamily\small] at (8,4.5) {0 1 1 1 0};
\node[anchor=west, font=\ttfamily\small] at (8,3.5) {1 0 0 0 0};
\node[anchor=west, font=\ttfamily\small] at (8,2.5) {1 0 0 0 0};
\node[anchor=west, font=\ttfamily\small] at (8,1.5) {1 0 0 0 0};
\node[anchor=west, font=\ttfamily\small] at (8,0.5) {0 1 1 1 0};
\node[popmuted, anchor=west] at (8,-0.7) {1 = pixel color\'e, 0 = pixel blanc};
\end{tikzpicture}
\end{center}
\legende{Une image en noir et blanc, fortement agrandie : chaque pixel est codé par un bit (1 = coloré, 0 = blanc). L'image devient une suite de 0 et de 1.}
\end{popfigure}

\textbf{Les sons.} Un son est une vibration continue de l'air. Pour le stocker, on mesure l'intensité du son à intervalles réguliers (par exemple des milliers de fois par seconde) : c'est l'\cle{échantillonnage}. Chaque mesure est un nombre, que l'on code en binaire. Un morceau de musique devient ainsi une très longue suite de nombres. Pour écouter le son, l'ordinateur refait le chemin inverse : il transforme les nombres en vibrations du haut-parleur. La qualité dépend de la fréquence d'échantillonnage (nombre de mesures par seconde) et de la précision de chaque mesure (nombre de bits).

\begin{aretenir}[Un principe unique pour toutes les données]
Texte, image, son, vidéo : dans tous les cas, la démarche est la même.
\begin{enumerate}
\item Découper la donnée en éléments simples (caractères, pixels, échantillons sonores).
\item Associer un nombre à chaque élément.
\item Coder ces nombres en binaire.
\end{enumerate}
C'est pourquoi on dit que l'ordinateur est une machine \textbf{numérique} : tout, absolument tout, y est représenté par des nombres écrits en binaire.
\end{aretenir}

\begin{exercice}[À toi de jouer]
\begin{enumerate}
\item À l'aide de la table ASCII simplifiée de cette leçon, coder en décimal puis en binaire le mot « ON ».
\item Décoder la suite de codes ASCII : 67 -- 79 -- 78.
\item Expliquer, en deux phrases, pourquoi le code ASCII seul ne suffit pas pour écrire un texte en langue ewondo, et quelle solution existe.
\end{enumerate}
\end{exercice}

\begin{corrige}[Exercice]
\begin{enumerate}
\item O = 79 et N = 78. Le mot « ON » se code \textbf{79 -- 78} en décimal, soit \textbf{1001111 -- 1001110} en binaire (sur 7 bits).
\item 67 = C, 79 = O, 78 = N : la suite représente le mot \textbf{« CON »}.
\item L'ASCII ne contient que 128 caractères pensés pour l'anglais : il ne possède ni les accents du français ni les lettres spéciales des langues camerounaises comme « ə ». La solution est d'utiliser l'\textbf{Unicode}, table universelle qui attribue un code à plus de 140 000 caractères de toutes les écritures du monde.
\end{enumerate}
\end{corrige}

\coursec{Synthèse et applications à la vie courante}

Nous arrivons au terme de cette leçon consacrée à l'organisation et à la codification des données. Dans les sections précédentes, nous avons disséqué chaque maillon de la chaîne : nous savons distinguer une \cle{donnée} brute d'une \cle{information} traitée (section 1), nous maîtrisons l'organisation hiérarchique en \cle{fichiers}, \cle{dossiers} et \cle{arborescence} (section 2), nous comprenons que l'ordinateur ne manipule in fine que des \cle{bits} groupés en \cle{octets} (section 3), et nous avons appris à \cle{coder} les nombres entiers et les caractères (code ASCII) en binaire (sections 4 et 5).

Il est temps maintenant de \textbf{reassembler le puzzle} : comment ces notions s'enchaînent-elles dans la réalité d'un système informatique, du moment où tu appuies sur une touche du clavier jusqu'à l'affichage du résultat à l'écran ou au stockage sur une clé USB ? C'est l'objet de cette synthèse, illustrée par trois situations concrètes inspirées de ton quotidien au Cameroun.

\courssub{Bilan : de la donnée brute au fichier binaire stocké}

\begin{popfigure}
\begin{tikzpicture}[
    node distance=1.2cm and 1.5cm,
    box/.style={rectangle, draw=popdark, thick, rounded corners, fill=popcream, minimum width=3.2cm, minimum height=1.1cm, align=center, font=\small},
    arrow/.style={-Stealth, thick, popblue},
    start/.style={box, fill=popgreenL, draw=popgreen},
    end/.style={box, fill=poporangeL, draw=poporange},
    process/.style={box, fill=popblueL, draw=popblue},
    storage/.style={cylinder, draw=popdark, thick, shape border rotate=90, aspect=0.3, fill=poppinkL, minimum width=2.5cm, minimum height=1.8cm, align=center, font=\small}
]
% Nœuds
\node[start, align=center] (saisie){Saisie utilisateur\\(clavier, souris,\\micro)};
\node[process, right=of saisie, align=center] (codage){Codification\\source $\to$ binaire\\(ASCII, binaire\\pur)};
\node[process, right=of codage, align=center] (orga){Organisation\\logique\\Fichier + nom,\\extension, date};
\node[process, right=of orga, align=center] (fs){Système de fichiers\\Arborescence\\Dossiers};
\node[storage, right=of fs, align=center] (stockage){Support physique\\Disque dur,\\Clé USB,\\Carte SD};
\node[process, below=of fs, xshift=1.5cm, align=center] (traitement){Traitement / Lecture\\Processeur + RAM\\Décodage};
\node[end, below=of stockage, xshift=-1.5cm, align=center] (info){Information\\exploitable\\Écran, Imprimante,\\Haut-parleur};

% Flèches principales
\draw[arrow] (saisie) -- node[above, font=\tiny, popmuted] {Donnée brute} (codage);
\draw[arrow] (codage) -- node[above, font=\tiny, popmuted] {Suite de bits} (orga);
\draw[arrow] (orga) -- node[above, font=\tiny, popmuted] {Fichier nommé} (fs);
\draw[arrow] (fs) -- node[above, font=\tiny, popmuted] {Écriture} (stockage);

% Boucle lecture/traitement
\draw[arrow] (stockage) -- ++(0,-1.2) -| (traitement) node[pos=0.25, below, font=\tiny, popmuted] {Lecture};
\draw[arrow] (traitement) -- (info) node[above, font=\tiny, popmuted] {Décodage};
\end{tikzpicture}
\legende{La chaîne complète du traitement de l'information : de la saisie au stockage, puis de la lecture à l'exploitation.}
\end{popfigure}

\begin{aretenir}[La chaîne complète : Saisie $\to$ Codification $\to$ Organisation $\to$ Stockage $\to$ Traitement]
Chaque donnée manipulée par l'ordinateur suit un parcours en cinq étapes :
\begin{enumerate}
    \item \textbf{Saisie (Entrée)} : un périphérique (clavier, scanner, microphone) capture une information du monde réel et la convertit en signal numérique \textbf{brut}.
    \item \textbf{Codification (Encodage)} : le système traduit ce signal brut selon une \textbf{convention standardisée} (binaire pur pour un entier, ASCII pour un caractère). C'est l'étape \emph{Donnée $\to$ Suite de bits}.
    \item \textbf{Organisation logique (Fichier)} : la suite de bits est regroupée dans une \textbf{entité nommée} : le \cle{fichier}. On y adjoint des informations descriptives (nom, extension, date de création, taille). L'extension (\texttt{.txt}, \texttt{.jpg}, \texttt{.mp3}) indique \emph{comment} redécoder les bits.
    \item \textbf{Stockage physique} : le système de fichiers place le fichier dans une \textbf{arborescence de dossiers} et écrit ses octets sur le support (disque dur, clé USB, carte SD).
    \item \textbf{Traitement / Restitution (Sortie)} : pour être utile, le fichier est lu (chargé en mémoire vive), \textbf{décodé} (opération inverse de l'étape 2) par le processeur sous contrôle d'un logiciel, et présenté comme \cle{information} à l'utilisateur (texte à l'écran, son dans les enceintes, image imprimée).
\end{enumerate}
\emph{Règle d'or :} \textbf{Sans organisation (nom, dossier, extension), une suite de bits sur un disque est inutilisable. Sans codage standard, l'échange entre machines est impossible.}
\end{aretenir}

\courssub{Méthode générale : résoudre un problème d'organisation et de codification}

\begin{methode}[Démarche structurée pour un projet de gestion de données]
Face à une situation réelle (sauvegarder des photos, classer des factures, envoyer un message codé), suis ces \textbf{5 étapes} pour produire une solution rigoureuse et justifiée :

\begin{enumerate}
    \item \textbf{Analyser le besoin (Problème)} :
    \begin{itemize}
        \item Quelle est la \textbf{nature} des données ? (Texte, nombres, images, sons, mixte)
        \item Quel est le \textbf{volume} estimé ? (Nombre d'éléments $\times$ taille unitaire)
        \item Quelles sont les \textbf{contraintes} ? (Capacité du support, rapidité d'accès, partage, durée de conservation)
        \item Quel est l'\textbf{objectif final} ? (Archivage, consultation fréquente, envoi, impression)
    \end{itemize}

    \item \textbf{Choisir l'organisation logique (Structure)} :
    \begin{itemize}
        \item Définir une \textbf{règle de nommage} claire (ex : \texttt{AAAA-MM-JJ\_Sujet\_Type.ext}).
        \item Concevoir une \textbf{arborescence de dossiers} pertinente (par date, par projet, par type, par personne).
        \item Choisir les \textbf{formats de fichiers} (extensions) adaptés : \texttt{.pdf} pour les documents figés, \texttt{.jpg} pour les photos, \texttt{.txt} pour les textes simples.
    \end{itemize}

    \item \textbf{Déterminer la codification (Encodage)} :
    \begin{itemize}
        \item Pour du texte : code \cle{ASCII} (caractères courants) ou ses extensions pour les accents.
        \item Pour des nombres entiers : binaire pur (conversion vue en section 4).
        \item Vérifier la \textbf{taille résultante} (poids du fichier) après codage.
    \end{itemize}

    \item \textbf{Valider la faisabilité (Vérification quantitative)} :
    \begin{itemize}
        \item Calculer la \textbf{taille totale} nécessaire = somme des tailles de chaque fichier, plus une marge ($\approx 10\%$) pour le système de fichiers.
        \item Comparer à la \textbf{capacité réelle disponible} du support (attention : une clé de \texttt{16 Go} n'offre que $\approx 14,9$ Go utilisables, voir plus loin).
        \item Si \texttt{Taille totale > Capacité disponible} $\to$ revoir l'organisation (sélection, support plus grand).
    \end{itemize}

    \item \textbf{Exécuter, tester et documenter} :
    \begin{itemize}
        \item Effectuer la copie / l'envoi / la création.
        \item \textbf{Tester la lecture} (ouvrir 2 ou 3 fichiers au hasard) pour valider le décodage.
        \item Rédiger un \textbf{rapport court} : choix faits, volume final, support utilisé, date, problèmes rencontrés. C'est la traçabilité.
    \end{itemize}
\end{enumerate}
\end{methode}

\courssub{Application 1 : organiser les fichiers d'un cybercafé (factures, photos, devoirs)}

Imaginons \textbf{Paul}, gérant du cybercafé « \emph{Connexion Express} » à Yaoundé, quartier Mvog-Ada. Il doit gérer trois types de documents numériques sur le PC administratif (disque de 500 Go) :
\begin{itemize}
    \item \textbf{Factures clients} (PDF, environ 200 Ko chacune, environ 50 par mois).
    \item \textbf{Photos d'événements} (anniversaires, tournois de jeux vidéo) : JPEG, environ 3 Mo chacune, environ 200 par mois.
    \item \textbf{Devoirs d'élèves} (Word/PDF, environ 500 Ko, environ 100 par mois) imprimés sur place.
\end{itemize}

\begin{exemplebox}[Arborescence proposée pour le cybercafé]
Paul applique la \textbf{méthode} (étape 2). Il choisit une organisation \textbf{par année/mois/type}, la plus naturelle pour une gestion administrative chronologique.

\begin{lstlisting}
D:\CyberData\
  +-- 2025\
  |     +-- 01_Janvier\
  |     |     +-- Factures\
  |     |     |     +-- 2025-01-02_Facture_001_Client_Ngon.pdf
  |     |     |     +-- 2025-01-15_Facture_002_Client_Mbia.pdf
  |     |     +-- Photos_Evenements\
  |     |     |     +-- 2025-01-05_Tournoi_Jeux_Equipe_A.jpg
  |     |     |     +-- 2025-01-20_Anniv_Petite_Soeur.jpg
  |     |     +-- Devoirs_Eleves\
  |     |           +-- 2025-01-10_Devoir_Math_3eme_Kamga.docx
  |     |           +-- 2025-01-22_Devoir_Physique_Tle_Ondoa.pdf
  |     +-- 02_Fevrier\
  |           +-- ...
  +-- Archives_Anterieures\
        +-- 2024\
              +-- ...
\end{lstlisting}

\textbf{Justification des choix :}
\begin{itemize}
    \item \textbf{Nommage} : \texttt{AAAA-MM-JJ\_Objet\_ID\_Auteur.ext} $\to$ le tri alphabétique correspond au tri chronologique ; l'auteur ou le client est identifiable ; l'extension est claire.
    \item \textbf{Dossiers par mois} : évite d'avoir des milliers de fichiers dans un même dossier (ce qui ralentit l'explorateur et les recherches).
    \item \textbf{Séparation par type} : permet des sauvegardes différenciées (Factures = critiques $\to$ double sauvegarde ; Photos = volumineuses $\to$ disque externe ; Devoirs = temporaires $\to$ suppression après impression).
\end{itemize}
\end{exemplebox}

\begin{exoresolu}[Estimation mensuelle de l'espace disque]
\textbf{Énoncé :} Calculer l'espace disque mensuel nécessaire pour l'activité du cybercafé. Le disque de 500 Go est-il suffisant pour un an d'activité sans nettoyage ?

\tcblower

\textbf{Solution détaillée :}
\begin{enumerate}
    \item \textbf{Volume mensuel par catégorie :}
    \begin{align*}
        \text{Factures} &: 50 \times 200\,\text{Ko} = 10\,000\,\text{Ko} \approx 10\,\text{Mo} \\
        \text{Photos} &: 200 \times 3\,\text{Mo} = 600\,\text{Mo} \\
        \text{Devoirs} &: 100 \times 500\,\text{Ko} = 50\,000\,\text{Ko} \approx 50\,\text{Mo}
    \end{align*}
    \item \textbf{Total mensuel brut :} $10 + 600 + 50 = \mathbf{660\,\text{Mo}}$.
    \item \textbf{Avec marge système de fichiers (10\,\%) :} $660 \times 1{,}1 = \mathbf{726\,\text{Mo/mois}}$.
    \item \textbf{Total annuel estimé :} $726 \times 12 = 8\,712\,\text{Mo} \approx \mathbf{8{,}7\,\text{Go}}$.
    \item \textbf{Conclusion :} le disque de \texttt{500 Go} est \textbf{largement suffisant} pour plusieurs années. Le vrai problème n'est donc pas la capacité, mais l'\textbf{organisation} pour retrouver un fichier rapidement.
\end{enumerate}
\end{exoresolu}

\courssub{Application 2 : estimer l'espace pour sauvegarder les photos d'un mariage}

Situation classique : \textbf{Chantal} a photographié le mariage de sa cousine à Douala. Elle a \textbf{500 photos} de \textbf{25 Mo} chacune sur la carte SD de son appareil. Elle veut les sauvegarder sur sa \textbf{clé USB de 16 Go} avant de vider la carte.

\begin{exemplebox}[Calcul complet : 500 photos sur une clé de 16 Go]
\textbf{Données :}
\begin{itemize}
    \item Nombre de photos : $N = 500$.
    \item Taille unitaire : $T_u = 25\,\text{Mo}$.
    \item Capacité annoncée de la clé : $C = 16\,\text{Go}$.
\end{itemize}

\textbf{Étape 1 : Volume total des données}
\[
V = N \times T_u = 500 \times 25\,\text{Mo} = 12\,500\,\text{Mo} = 12{,}5\,\text{Go}.
\]

\textbf{Étape 2 : Capacité réellement utilisable de la clé}

Les fabricants comptent en puissances de 10 ($1\,\text{Go} = 10^9$ octets), mais l'ordinateur compte en puissances de 2 ($1\,\text{Gio} = 2^{30}$ octets $= 1\,073\,741\,824$ octets). La capacité vue par l'ordinateur est donc :
\[
C_{\text{réelle}} = \frac{16 \times 10^9}{2^{30}} \approx \frac{16\,000\,000\,000}{1\,073\,741\,824} \approx \mathbf{14{,}9\,\text{Gio}}.
\]
Le système de fichiers consomme encore une petite part ; on dispose d'environ $\mathbf{14{,}7\,\text{Gio}}$.

\textbf{Étape 3 : Comparaison (dans les mêmes unités)}

Convertissons le volume des données dans l'unité de l'ordinateur :
\[
V = \frac{12{,}5 \times 10^9}{2^{30}} \approx \mathbf{11{,}6\,\text{Gio}}.
\]
\textbf{Résultat :} $11{,}6\,\text{Gio} < 14{,}7\,\text{Gio}$ $\to$ \textbf{la copie est possible !} Il reste environ $3$ Gio de marge.

\textbf{Étape 4 : Et si Chantal compressait ses photos (5 Mo/photo) ?}
\[
V' = 500 \times 5\,\text{Mo} = 2\,500\,\text{Mo} = 2{,}5\,\text{Go} \approx 2{,}3\,\text{Gio}.
\]
Gain de place énorme, mais avec une \textbf{perte de qualité}. Le choix dépend de l'usage : impression grand format $\to$ garder les fichiers d'origine ; partage sur WhatsApp $\to$ la version compressée suffit.
\end{exemplebox}

\begin{attention}[Erreurs fréquentes : unités et limites des supports]
\begin{itemize}
    \item \textbf{Piège n°1 :} acheter une clé « 16 Go » en pensant pouvoir y mettre 16 Gio de données. \textbf{Faux.} On ne pourra y mettre que $\approx 14,9$ Gio. Règle pratique : \texttt{Capacité réelle $\approx$ Capacité annoncée $\times$ 0,93}.
    \item \textbf{Piège n°2 :} comparer directement $12\,500$ Mo à « 16 Go » sans conversion. Ici la copie tient quand même, mais avec une clé « 8 Go » ($\approx 7,45$ Gio réels), $11,6 > 7,45$ $\to$ \textbf{échec de la copie} en cours de route.
    \item \textbf{Piège n°3 :} oublier de laisser de la marge. \textbf{Bonne pratique :} toujours laisser \textbf{10 à 15\,\% d'espace libre} sur une clé USB ou une carte SD pour garantir son bon fonctionnement.
    \item \textbf{Piège n°4 :} oublier de \textbf{vérifier après la copie} : ouvrir quelques fichiers copiés avant d'effacer les originaux de la carte SD !
\end{itemize}
\end{attention}

\courssub{Application 3 : transmettre un message codé en binaire entre élèves}

Activité ludique mais rigoureuse : \textbf{Alice} et \textbf{Bob}, élèves en 3ème au Lycée de Biyem-Assi, veulent s'échanger un court message secret pendant l'heure d'informatique (sur papier, sans téléphone), en utilisant ce qu'ils ont appris sur le code ASCII.

\begin{experience}[TP guidé : échange d'un message binaire « main dans la main »]
\textbf{Objectif :} mettre en œuvre la chaîne complète : \emph{Texte $\to$ Codage ASCII (binaire) $\to$ Transmission (papier) $\to$ Décodage $\to$ Texte}.

\textbf{Matériel :} papier, crayon, table ASCII (extrait fourni ci-dessous).

\textbf{Manipulation :}
\begin{enumerate}
    \item \textbf{Accord préalable (protocole) :} Alice et Bob conviennent :
        \begin{itemize}
            \item Codage : \textbf{ASCII sur 8 bits par caractère} (on complète le code sur 7 bits par un 0 à gauche).
            \item Séparateur : \textbf{un espace} entre les octets (pour la lisibilité humaine).
            \item Fin de message : la séquence \texttt{00000000}.
        \end{itemize}
    \item \textbf{Écriture (Alice - Encodage) :} Alice choisit le mot \texttt{SALUT}.
        \begin{itemize}
            \item Elle consulte la table : \texttt{S} = 83, \texttt{A} = 65, \texttt{L} = 76, \texttt{U} = 85, \texttt{T} = 84.
            \item Elle convertit chaque décimal en binaire sur 8 bits (méthode des puissances de 2) :
            \begin{align*}
                83 &= 64+16+2+1 = \texttt{01010011} \\
                65 &= 64+1 = \texttt{01000001} \\
                76 &= 64+8+4 = \texttt{01001100} \\
                85 &= 64+16+4+1 = \texttt{01010101} \\
                84 &= 64+16+4 = \texttt{01010100}
            \end{align*}
            \item Elle écrit sur le papier : \texttt{01010011 01000001 01001100 01010101 01010100 00000000}.
        \end{itemize}
    \item \textbf{Transmission :} Alice passe le papier à Bob.
    \item \textbf{Lecture (Bob - Décodage) :} Bob groupe les bits par 8, convertit chaque octet en décimal, puis retrouve le caractère dans la table.
    \item \textbf{Vérification :} Bob lit « SALUT ». Succès !
\end{enumerate}

\textbf{Table de référence (extrait du code ASCII) :}
\begin{center}
\begin{tabularx}{\linewidth}{|c|c|X|c|c|X|}
\hline
\rowcolor{popblueL}
\textbf{Caractère} & \textbf{Décimal} & \textbf{Binaire (8 bits)} & \textbf{Caractère} & \textbf{Décimal} & \textbf{Binaire (8 bits)} \\ \hline
Espace & 32 & 00100000 & A & 65 & 01000001 \\ \hline
B & 66 & 01000010 & C & 67 & 01000011 \\ \hline
L & 76 & 01001100 & S & 83 & 01010011 \\ \hline
T & 84 & 01010100 & U & 85 & 01010101 \\ \hline
a & 97 & 01100001 & b & 98 & 01100010 \\ \hline
0 & 48 & 00110000 & 1 & 49 & 00110001 \\ \hline
2 & 50 & 00110010 & 5 & 53 & 00110101 \\ \hline
! & 33 & 00100001 & ? & 63 & 00111111 \\ \hline
\end{tabularx}
\end{center}

\textbf{Observations et interprétation :}
\begin{itemize}
    \item \textbf{Rigidité du protocole :} si Alice oublie un espace entre deux octets, Bob ne sait plus où couper la suite de bits : le message devient ambigu.
    \item \textbf{Coût :} 5 caractères + fin de message $\to$ 48 bits à écrire à la main. C'est long ! C'est pourquoi ce travail est confié aux machines.
    \item \textbf{Fragilité :} une erreur sur un seul bit change complètement la lettre : \texttt{01010011} (S) devient \texttt{01010010} (R). D'où l'importance de la concentration... et des contrôles !
\end{itemize}
\end{experience}

\courssub{Rédaction d'une démarche justifiée : du problème à la vérification}

Le programme insiste sur le \textbf{savoir-faire} : « Rédiger et justifier une démarche, une réponse ou une conclusion ». Voici le modèle type attendu au BEPC pour une question de synthèse sur l'organisation et la codification des données.

\begin{methode}[Structure type d'une réponse justifiée (épreuve écrite)]
Pour toute question du type « Proposez une organisation pour... » ou « Est-il possible de stocker... ? Justifiez. », structure ta réponse en \textbf{4 étapes distinctes} :

\begin{enumerate}
    \item \textbf{Problème et données (état des lieux)} : reprends les chiffres clés de l'énoncé.
    \begin{quote}
    \textit{« Le problème est de sauvegarder 500 photos de 25 Mo chacune sur une clé USB annoncée à 16 Go. Données : $N=500$, $T_u=25$ Mo, $C=16$ Go. »}
    \end{quote}

    \item \textbf{Choix d'organisation et de codage (stratégie)} : explique \textbf{pourquoi} tu choisis tel format, tel nommage, telle arborescence.
    \begin{quote}
    \textit{« Je choisis de conserver les fichiers d'origine (pas de perte de qualité) pour l'archivage. Les fichiers seront placés dans un dossier unique \texttt{2025-06-15\_Mariage} et nommés \texttt{Mariage\_001.jpg}, \texttt{Mariage\_002.jpg}, etc. »}
    \end{quote}

    \item \textbf{Calculs et vérification quantitative (preuve)} : montre les formules, les conversions d'unités, la comparaison.
    \begin{quote}
    \textit{« Volume total $V = 500 \times 25 = 12\,500$ Mo $= 12,5$ Go. Capacité réelle de la clé : $16 \times 0{,}93 \approx 14,9$ Gio. Conversion : $V \approx 11,6$ Gio. Comme $11,6 < 14,9$, la sauvegarde est \textbf{possible} avec environ 3 Gio de marge. »}
    \end{quote}

    \item \textbf{Conclusion et recommandations (synthèse)} : réponse finale claire + conseils.
    \begin{quote}
    \textit{« Conclusion : la clé de 16 Go convient. Recommandation : copier le dossier, vérifier l'ouverture de 3 photos au hasard, puis éjecter la clé proprement avant de la retirer. »}
    \end{quote}
\end{enumerate}
\end{methode}

\begin{exoresolu}[Sujet type BEPC : organisation des bulletins de notes]
\textbf{Énoncé :} L'administration du Lycée de Nkolbisson doit archiver les bulletins de notes de 600 élèves (3 trimestres, soit 3 bulletins par élève et par an) sur 5 ans. Chaque bulletin est un fichier PDF de 150 Ko. L'archive sera stockée sur un disque dur externe de 500 Go.
\begin{enumerate}
    \item Proposer une arborescence de dossiers et une convention de nommage.
    \item Calculer le volume total sur 5 ans. Vérifier la faisabilité.
    \item Donner deux avantages de cette organisation pour la recherche ultérieure.
\end{enumerate}

\tcblower

\textbf{Solution détaillée :}

\textbf{1. Organisation proposée :}
\begin{lstlisting}
D:\Archives_Bulletins\
  +-- 2024-2025\
  |     +-- 6eme\
  |     |     +-- 6eme_A\
  |     |     |     +-- 2024-2025_6eme_A_Abena_Paul_T1.pdf
  |     |     |     +-- ...
  |     |     +-- 6eme_B\
  |     +-- 5eme\
  |           +-- ...
  +-- 2025-2026\
  |     +-- ...
  +-- 2028-2029\
        +-- ...
\end{lstlisting}
\textbf{Nommage :} \texttt{AnneeScolaire\_Classe\_Nom\_Prenom\_Trimestre.pdf} (ex : \texttt{2024-2025\_3eme\_C\_Mvondo\_Jean\_T2.pdf}).

\textbf{Justification :} hiérarchie Année $\to$ Niveau $\to$ Classe $\to$ Fichiers. Elle permet de supprimer une année complète facilement, ou de retrouver tous les bulletins d'un élève en cherchant son nom.

\textbf{2. Calcul du volume :}
\begin{align*}
\text{Nombre de bulletins par an} &= 600 \times 3 = 1\,800 \\
\text{Volume par an} &= 1\,800 \times 150\,\text{Ko} = 270\,000\,\text{Ko} \approx 270\,\text{Mo} \\
\text{Volume sur 5 ans} &= 5 \times 270\,\text{Mo} = 1\,350\,\text{Mo} \approx 1{,}35\,\text{Go} \\
\text{Avec marge de 10\,\%} &\approx 1{,}5\,\text{Go}
\end{align*}
Le disque de 500 Go ($\approx 465$ Gio réels) est très largement suffisant : $1{,}5 \ll 465$. \textbf{L'archivage est faisable sans difficulté.}

\textbf{3. Avantages pour la recherche :}
\begin{enumerate}
    \item \textbf{Navigation intuitive :} on retrouve le bulletin de Jean Mvondo en 3 clics (Année $\to$ 3ème $\to$ 3ème C) sans outil de recherche.
    \item \textbf{Recherche par nom :} la règle de nommage régulière permet de retrouver rapidement, par exemple, tous les bulletins du trimestre T2 d'une classe donnée.
\end{enumerate}
\end{exoresolu}

\courssub{Le savais-tu ? : Le Mobile Money, une organisation de données géante dans ta poche}

\begin{savaistu}[MTN MoMo, Orange Money : la rigueur des données financières]
Chaque fois qu'un utilisateur fait un transfert \textbf{MTN MoMo} ou \textbf{Orange Money} depuis son téléphone, une chaîne de traitement de données très structurée s'exécute en quelques secondes :

\begin{enumerate}
    \item \textbf{Saisie structurée :} l'utilisateur tape le numéro du destinataire (9 chiffres), le montant (un nombre entier en FCFA) et son code secret. Le téléphone \textbf{vérifie le format} de chaque donnée avant l'envoi.
    \item \textbf{Codification :} la demande est convertie en une suite de bits et \textbf{protégée} (rendue illisible pour les autres) avant de voyager sur le réseau mobile.
    \item \textbf{Traitement par le serveur central :} l'ordinateur de l'opérateur effectue trois opérations indissociables : retirer le montant du compte de l'expéditeur, l'ajouter au compte du destinataire, et \textbf{enregistrer l'opération} dans un historique (date, heure, numéros, montant).
    \item \textbf{Organisation du stockage :} les données ne sont pas rangées dans de simples fichiers, mais dans des \textbf{tableaux structurés} (listes d'utilisateurs, de comptes, de transactions) conservés sur des serveurs sécurisés, avec des copies de sauvegarde.
    \item \textbf{Réponse (décodage) :} le serveur renvoie le résultat ; le téléphone le décode et affiche : « Transfert de 5 000 FCFA effectué. Nouveau solde : 12 300 FCFA. »
\end{enumerate}

\textbf{Le lien avec notre leçon :}
\begin{itemize}
    \item \textbf{Donnée brute} $\to$ les touches tapées sur le téléphone.
    \item \textbf{Codification} $\to$ la conversion en binaire et la protection du message.
    \item \textbf{Organisation} $\to$ des tableaux de données bien structurés et nommés.
    \item \textbf{Stockage} $\to$ des serveurs avec sauvegardes, dans des centres informatiques sécurisés.
    \item \textbf{Information} $\to$ le message de confirmation que l'utilisateur lit.
\end{itemize}
\textbf{Enjeu critique :} une erreur de codage (un montant mal converti), une mauvaise organisation (recherche trop lente) ou une panne de stockage sans sauvegarde (perte de l'historique) sont \textbf{inacceptables} quand il s'agit d'argent. C'est pourquoi les opérateurs de Mobile Money recrutent des informaticiens rigoureux, formés dès le collège à ces concepts de base !
\end{savaistu}

\courssub{Exercices de synthèse (préparation BEPC)}

\begin{exercice}[Exercice 1 : sauvegarde sélective]
Tu disposes d'une carte SD contenant :
\begin{itemize}
    \item 1\,200 photos JPEG (moyenne 4,5 Mo) ;
    \item 15 vidéos (moyenne 1,2 Go) ;
    \item 50 photos grand format (moyenne 28 Mo).
\end{itemize}
Tu veux tout copier sur un disque dur externe de 500 Go (déjà occupé à 320 Go).
\begin{enumerate}
    \item Calcule l'espace total occupé sur la carte SD (en Go).
    \item Calcule l'espace libre restant sur le disque dur (en Go).
    \item La copie complète est-elle possible ? Justifie par le calcul.
    \item Si l'espace libre n'était que de 6 Go, propose une stratégie pour sauvegarder le maximum de données prioritaires (Photos $>$ Vidéos $>$ Grand format).
\end{enumerate}
\end{exercice}

\begin{exercice}[Exercice 2 : codage et transmission d'un mot de passe]
Dans le cadre d'un projet de classe, vous devez transmettre le mot de passe Wi-Fi \texttt{Lycee2025!} à un camarade via un papier. Vous décidez de le coder en \textbf{binaire} à l'aide du code ASCII (1 octet par caractère).
\begin{enumerate}
    \item Donne le code décimal ASCII de chaque caractère : \texttt{L, y, c, e, e, 2, 0, 2, 5, !} (on donne : \texttt{y} = 121, \texttt{c} = 99, \texttt{e} = 101).
    \item Convertis chaque décimal en binaire sur 8 bits.
    \item Écris la trame complète à transmettre (octets séparés par un espace).
    \item Combien de bits au total ? Si l'on commet une erreur sur 1 bit sur les 80, quel est le taux d'erreur (en pourcentage) ?
\end{enumerate}
\end{exercice}

\begin{exercice}[Exercice 3 : organisation d'une médiathèque scolaire]
Le club informatique du lycée veut numériser 2\,000 fiches de lecture (texte + image de couverture). Chaque fiche = 1 page enregistrée en \texttt{.pdf} + 1 image JPEG (couverture).
\begin{itemize}
    \item Page PDF : environ 80 Ko.
    \item Couverture JPEG : environ 500 Ko.
\end{itemize}
\begin{enumerate}
    \item Propose une arborescence de dossiers et un nommage normalisé pour ces 2\,000 paires de fichiers.
    \item Estime le volume total en Mo puis en Go.
    \item Quelle clé USB minimale recommandes-tu (tailles standards : 8, 16, 32, 64 Go) ? Justifie avec une marge de 20\,\%.
\end{enumerate}
\end{exercice}

\begin{corrige}[Exercice 1]
\begin{enumerate}
    \item \textbf{Volume de la carte SD :}
    \begin{align*}
        \text{Photos} &= 1\,200 \times 4{,}5\,\text{Mo} = 5\,400\,\text{Mo} = 5{,}4\,\text{Go} \\
        \text{Vidéos} &= 15 \times 1{,}2\,\text{Go} = 18\,\text{Go} \\
        \text{Grand format} &= 50 \times 28\,\text{Mo} = 1\,400\,\text{Mo} = 1{,}4\,\text{Go} \\
        \text{Total} &= 5{,}4 + 18 + 1{,}4 = \mathbf{24{,}8\,\text{Go}}
    \end{align*}
    \item \textbf{Espace libre sur le disque :} $500 - 320 = \mathbf{180\,\text{Go}}$ (on néglige ici l'écart Go/Gio, identique pour les deux valeurs).
    \item \textbf{Faisabilité :} $24{,}8\,\text{Go} < 180\,\text{Go}$ $\to$ \textbf{oui, la copie complète est possible}, avec une large marge.
    \item \textbf{Stratégie avec 6 Go libres :} priorité 1 : les photos JPEG ($5{,}4$ Go) $\to$ elles tiennent dans les 6 Go. Priorité 2 : les vidéos (18 Go) $\to$ impossibles, on n'en copie aucune en entier. Priorité 3 : le grand format ($1{,}4$ Go) $\to$ ne tient plus ($5{,}4 + 1{,}4 = 6{,}8 > 6$). \textbf{Solution :} copier toutes les photos JPEG ($5{,}4$ Go) + environ 20 fichiers grand format ($0{,}56$ Go) pour rester sous 6 Go, et reporter les vidéos sur un autre support.
\end{enumerate}
\end{corrige}

\begin{corrige}[Exercice 2]
\begin{enumerate}
    \item \textbf{Codes décimaux ASCII :} L = 76, y = 121, c = 99, e = 101, e = 101, 2 = 50, 0 = 48, 2 = 50, 5 = 53, ! = 33.
    \item \textbf{Conversion en binaire sur 8 bits :}
    \begin{align*}
        76  &= 64+8+4 = \texttt{01001100} \\
        121 &= 64+32+16+8+1 = \texttt{01111001} \\
        99  &= 64+32+2+1 = \texttt{01100011} \\
        101 &= 64+32+4+1 = \texttt{01100101} \\
        101 &= \texttt{01100101} \\
        50  &= 32+16+2 = \texttt{00110010} \\
        48  &= 32+16 = \texttt{00110000} \\
        50  &= \texttt{00110010} \\
        53  &= 32+16+4+1 = \texttt{00110101} \\
        33  &= 32+1 = \texttt{00100001}
    \end{align*}
    \item \textbf{Trame complète :}

    \texttt{01001100 01111001 01100011 01100101 01100101 00110010 00110000 00110010 00110101 00100001}
    \item \textbf{Total :} $10 \times 8 = 80$ bits. Taux d'erreur : $\dfrac{1}{80} = 0{,}0125 = \mathbf{1{,}25\,\%}$.
\end{enumerate}
\end{corrige}

\begin{corrige}[Exercice 3]
\begin{enumerate}
    \item \textbf{Arborescence proposée :}
    \begin{lstlisting}
Mediatheque_Lycee\
  +-- Fiches_PDF\
  |     +-- 6eme\
  |     |     +-- 6eme_Abena_Paul_Fiche.pdf
  |     |     +-- ...
  |     +-- 5eme\
  |           +-- ...
  +-- Couvertures_JPG\
        +-- 6eme\
        |     +-- 6eme_Abena_Paul_Couv.jpg
        |     +-- ...
        +-- 5eme\
              +-- ...
    \end{lstlisting}
    \textbf{Nommage :} \texttt{Classe\_Nom\_Prenom\_Type.ext} (Type = Fiche ou Couv). La fiche et sa couverture sont liées par un nom de base commun.
    \item \textbf{Volume total :}
    \begin{align*}
        \text{Fiches} &= 2\,000 \times 80\,\text{Ko} = 160\,000\,\text{Ko} \approx 160\,\text{Mo} \\
        \text{Couvertures} &= 2\,000 \times 500\,\text{Ko} = 1\,000\,000\,\text{Ko} \approx 1\,000\,\text{Mo} \\
        \text{Total} &\approx 1\,160\,\text{Mo} \approx \mathbf{1{,}16\,\text{Go}}
    \end{align*}
    \item \textbf{Support :} avec une marge de 20\,\% : $1{,}16 \times 1{,}2 \approx 1{,}4$ Go nécessaires. Une clé de \textbf{8 Go} (environ 7,4 Go réels) suffit largement. \textbf{Recommandation :} une clé de \textbf{16 Go}, au prix très proche, offre une marge confortable pour les ajouts futurs de la médiathèque.
\end{enumerate}
\end{corrige}

\begin{aretenir}[Check-list finale : suis-je prêt pour le BEPC sur cette leçon ?]
\begin{itemize}
    \item[$\square$] Je sais distinguer \textbf{donnée}, \textbf{information} et \textbf{type de donnée} (entier, réel, caractère, chaîne de caractères).
    \item[$\square$] Je sais organiser une arborescence \textbf{logique, nommée et hiérarchique} adaptée à un contexte (cybercafé, école, usage personnel).
    \item[$\square$] Je connais les unités : \textbf{bit, octet, Ko, Mo, Go, To}, et je sais convertir de l'une à l'autre ($1\,\text{Ko} = 1\,024$ octets en binaire, $1\,\text{Mo} = 1\,024$ Ko, etc.).
    \item[$\square$] Je sais que la capacité réelle d'un support est inférieure à la capacité annoncée (règle pratique : $\times 0{,}93$).
    \item[$\square$] Je sais \textbf{coder un entier positif en binaire} (méthode des puissances de 2) et \textbf{décoder} un nombre binaire.
    \item[$\square$] Je connais le code \textbf{ASCII} : chaque caractère possède un numéro, codé sur 7 bits (128 caractères), que l'on écrit sur 8 bits en ajoutant un 0.
    \item[$\square$] Je sais qu'une \textbf{image} est un tableau de points (pixels) codés en binaire, et qu'un \textbf{son} est une suite de valeurs numériques.
    \item[$\square$] J'applique la \textbf{méthode en 5 étapes} (Analyser, Organiser, Coder, Vérifier, Exécuter/Tester) pour tout problème concret.
    \item[$\square$] Je sais \textbf{rédiger une démarche justifiée} : Problème $\to$ Choix $\to$ Calculs/Preuve $\to$ Conclusion/Recommandation.
\end{itemize}
\textbf{Bravo !} Tu as maintenant les bases solides pour comprendre comment l'ordinateur « pense » en binaire et comment nous, humains, organisons ce monde binaire en informations utiles. Les prochaines leçons te montreront comment \textbf{traiter} ces données automatiquement grâce aux algorithmes. Courage pour la suite !
\end{aretenir}

\newpage

\coursec{Activité d'intégration}

\coursec{Leçon 4 : Organisation des données en informatique}

\courssub{Objectifs de l'activité}

À l'issue de cette activité d'intégration, tu devras être capable de :

\begin{itemize}
\item distinguer \cle{donnée}, \cle{information} et \cle{type de données} dans une situation réelle de la vie courante ;
\item transformer des données brutes (écrites à la main, désordonnées) en un \cle{tableau structuré} dont les colonnes sont des champs bien définis ;
\item proposer une \cle{arborescence de dossiers et de fichiers} cohérente, avec des noms de fichiers corrects ;
\item \cle{coder un nombre entier en binaire} et \cle{coder un texte en ASCII} ;
\item calculer l'\cle{espace disque} nécessaire au stockage de fiches, en utilisant les unités de mesure de l'information (Ko, Mo) ;
\item rédiger et justifier une démarche, une réponse ou une conclusion, comme l'exige le programme.
\end{itemize}

\courssub{Situation}

Au marché Mokolo, l'un des plus grands marchés de Yaoundé, Madame Béatrice tient un étal de produits vivriers depuis quinze ans. Chaque soir, elle note ses ventes à la main dans un cahier, sans ordre précis. Voici un extrait fidèle de ses notes pour la semaine du lundi 6 au samedi 11 janvier 2025 :

\begin{exemplebox}[Le cahier de notes de Madame Béatrice (données brutes)]
\begin{itemize}
\item « Lundi 6/01 : 4 kg de tomates à 500 FCFA le kg ; 2 sacs d'oignons à 1\,000 FCFA le sac. »
\item « Mardi 7/01 : 6 kg de tomates à 500 FCFA ; 3 kg de piment à 800 FCFA le kg. »
\item « Mercredi 8/01 : 5 kg de tomates à 500 FCFA ; 1 sac d'oignons à 1\,000 FCFA ; 2 kg de piment à 800 FCFA. »
\item « Jeudi 9/01 : 3 kg de tomates à 500 FCFA ; 4 kg de piment à 800 FCFA. »
\item « Vendredi 10/01 : 7 kg de tomates à 500 FCFA ; 2 sacs d'oignons à 1\,000 FCFA. »
\item « Samedi 11/01 : 8 kg de tomates à 500 FCFA ; 3 sacs d'oignons à 1\,000 FCFA ; 1 kg de piment à 800 FCFA. »
\end{itemize}
\end{exemplebox}

Le fils de Madame Béatrice, élève en classe de 3ème, vient de recevoir un ordinateur offert par son oncle. Il propose à sa mère de créer un \emph{cahier numérique} pour remplacer le cahier papier. Mais il se pose plusieurs questions : comment ranger toutes ces ventes dans l'ordinateur ? Comment l'ordinateur va-t-il représenter les nombres et les noms de produits ? Combien de place faudra-t-il sur le disque dur ?

\textbf{Données numériques supplémentaires :}
\begin{itemize}
\item Madame Béatrice vend 3 produits différents : tomates, oignons, piment.
\item Elle estime créer \num{1500} fiches produits par an, chaque fiche occupant \num{2} Ko sur le disque.
\item Pour le codage ASCII, on donne : A = 65, C = 67, E = 69, I = 73, M = 77, O = 79, S = 83, T = 84 (codes décimaux des lettres majuscules utiles).
\end{itemize}

\courssub{Consignes}

Par groupes de 4 élèves, répondez aux questions suivantes en rédigeant soigneusement vos démarches. Chaque réponse doit être \textbf{justifiée}.

\begin{enumerate}
\item \textbf{Trier et structurer.} À partir du cahier de notes, construisez un tableau structuré des ventes de la semaine. Chaque ligne représentera une vente et chaque colonne un champ : Date, Produit, Quantité, Unité, Prix unitaire (FCFA), Montant (FCFA). Pour chaque champ, indiquez le \cle{type de données} choisi (texte, nombre entier, date) et justifiez votre choix.
\item \textbf{Calculer.} Ajoutez une colonne « Montant » et calculez le montant de chaque vente (Montant = Quantité $\times$ Prix unitaire), puis le montant total de la semaine.
\item \textbf{Ranger.} Proposez une arborescence de dossiers et de fichiers pour ranger les ventes de toute l'année, par mois et par produit. Donnez trois exemples de noms de fichiers corrects (sans espace, sans accent, sans caractère spécial interdit) et expliquez pourquoi ces règles de nommage sont importantes.
\item \textbf{Coder les nombres.} Calculez le nombre total de produits vendus dans la semaine (toutes quantités confondues), puis convertissez ce nombre en binaire en détaillant la méthode des divisions successives par 2.
\item \textbf{Coder les caractères.} Identifiez le produit le plus vendu de la semaine, puis codez son nom en ASCII (en majuscules) à l'aide des codes fournis. Donnez la suite des codes décimaux obtenue.
\item \textbf{Estimer l'espace disque.} Calculez l'espace disque nécessaire pour stocker les \num{1500} fiches de l'année, en Ko puis en Mo. Madame Béatrice dispose d'une clé USB de 4 Go : cet espace est-il suffisant ? Justifiez.
\item \textbf{Rédiger la conclusion.} Rédigez un court texte (10 à 15 lignes) adressé à Madame Béatrice, expliquant chaque choix effectué (structure du tableau, arborescence, codage) et les avantages du cahier numérique par rapport au cahier papier.
\end{enumerate}

\textbf{Restitution :} chaque groupe présente ses résultats à l'oral (5 minutes). La correction collective est construite au tableau avec l'enseignant.

\courssub{Ressources à mobiliser}

\begin{aretenir}[Boîte à outils de la leçon]
\begin{itemize}
\item \textbf{Types de données :} texte (chaîne de caractères), nombre entier, nombre décimal, date, booléen.
\item \textbf{Organisation structurée :} un tableau = des lignes (enregistrements) et des colonnes (champs) ; une arborescence = des dossiers qui contiennent des sous-dossiers et des fichiers.
\item \textbf{Noms de fichiers :} pas d'espace, pas d'accent, pas de caractères spéciaux (\texttt{\textbackslash\ / : * ? " < > |}) ; on préfère les tirets ou les tirets bas.
\item \textbf{Conversion décimal $\rightarrow$ binaire :} divisions successives par 2 ; on lit les restes \emph{de bas en haut}.
\item \textbf{Vérification :} un nombre binaire $(b_n \ldots b_1 b_0)_2$ vaut $\sum b_i \times 2^i$.
\item \textbf{Codage ASCII :} chaque caractère est représenté par un nombre (par exemple A = 65).
\item \textbf{Unités de stockage :} 1 Ko = \num{1024} octets ; 1 Mo = \num{1024} Ko ; 1 Go = \num{1024} Mo.
\end{itemize}
\end{aretenir}

\courssub{Démarche et corrigé commenté}

\begin{exoresolu}[Le cahier numérique du marché Mokolo — corrigé détaillé]
Énoncé : structurer les ventes de la semaine, proposer une arborescence, coder le total en binaire et le produit le plus vendu en ASCII, estimer l'espace disque, puis justifier chaque choix.
\tcblower
\textbf{Étape 1 — Tableau structuré des ventes.} On dépouille le cahier ligne par ligne. Chaque vente devient un enregistrement :

\begin{center}
\begin{tabular}{|l|l|c|l|c|c|}
\hline
\rowcolor{popblueL} \textbf{Date} & \textbf{Produit} & \textbf{Qté} & \textbf{Unité} & \textbf{P.U. (F)} & \textbf{Montant (F)} \\
\hline
06/01/2025 & Tomates & 4 & kg & 500 & \num{2000} \\
\hline
06/01/2025 & Oignons & 2 & sac & \num{1000} & \num{2000} \\
\hline
07/01/2025 & Tomates & 6 & kg & 500 & \num{3000} \\
\hline
07/01/2025 & Piment & 3 & kg & 800 & \num{2400} \\
\hline
08/01/2025 & Tomates & 5 & kg & 500 & \num{2500} \\
\hline
08/01/2025 & Oignons & 1 & sac & \num{1000} & \num{1000} \\
\hline
08/01/2025 & Piment & 2 & kg & 800 & \num{1600} \\
\hline
09/01/2025 & Tomates & 3 & kg & 500 & \num{1500} \\
\hline
09/01/2025 & Piment & 4 & kg & 800 & \num{3200} \\
\hline
10/01/2025 & Tomates & 7 & kg & 500 & \num{3500} \\
\hline
10/01/2025 & Oignons & 2 & sac & \num{1000} & \num{2000} \\
\hline
11/01/2025 & Tomates & 8 & kg & 500 & \num{4000} \\
\hline
11/01/2025 & Oignons & 3 & sac & \num{1000} & \num{3000} \\
\hline
11/01/2025 & Piment & 1 & kg & 800 & 800 \\
\hline
\end{tabular}
\end{center}

\textbf{Types de données justifiés :} Date : type \emph{date} (pour trier chronologiquement) ; Produit et Unité : type \emph{texte} ; Quantité et Prix unitaire : type \emph{nombre entier} (pour faire des calculs) ; Montant : nombre entier calculé.

\textbf{Étape 2 — Montant total de la semaine.}
\[ \num{2000}+\num{2000}+\num{3000}+\num{2400}+\num{2500}+\num{1000}+\num{1600}+\num{1500}+\num{3200}+\num{3500}+\num{2000}+\num{4000}+\num{3000}+800 = \num{32500} \text{ FCFA.} \]

\textbf{Étape 3 — Arborescence proposée.} On crée un dossier principal \texttt{VENTES\_2025}, contenant un sous-dossier par mois, lui-même contenant un fichier par produit :

\begin{lstlisting}
VENTES_2025/
|-- JANVIER/
|   |-- ventes_tomates_janvier.txt
|   |-- ventes_oignons_janvier.txt
|   |-- ventes_piment_janvier.txt
|-- FEVRIER/
|   |-- ventes_tomates_fevrier.txt
|   |-- ...
\end{lstlisting}

Justification : les noms sont en minuscules, sans espace ni accent, et décrivent le contenu ; on retrouve ainsi un fichier en quelques secondes, même après plusieurs mois.

\textbf{Étape 4 — Codage binaire du total.} Quantités vendues : tomates $4+6+5+3+7+8 = 33$ ; oignons $2+1+2+3 = 8$ ; piment $3+2+4+1 = 10$. Total : $33+8+10 = 51$ produits. Divisions successives par 2 :
\begin{align*}
51 &= 2 \times 25 + 1 & 25 &= 2 \times 12 + 1 & 12 &= 2 \times 6 + 0 \\
6 &= 2 \times 3 + 0 & 3 &= 2 \times 1 + 1 & 1 &= 2 \times 0 + 1
\end{align*}
On lit les restes de bas en haut : $51 = (110011)_2$. Vérification : $1\times 2^5 + 1\times 2^4 + 0\times 2^3 + 0\times 2^2 + 1\times 2^1 + 1\times 2^0 = 32+16+2+1 = 51$. Correct !

\textbf{Étape 5 — Codage ASCII du produit le plus vendu.} Le produit le plus vendu est la \textbf{tomate} (33 unités). En majuscules : TOMATES. Avec les codes fournis :
\begin{center}
\begin{tabular}{|c|c|c|c|c|c|c|}
\hline
\rowcolor{popblueL} T & O & M & A & T & E & S \\
\hline
84 & 79 & 77 & 65 & 84 & 69 & 83 \\
\hline
\end{tabular}
\end{center}
L'ordinateur stocke donc « TOMATES » comme la suite de nombres 84, 79, 77, 65, 84, 69, 83, chacun étant lui-même codé en binaire.

\textbf{Étape 6 — Espace disque.}
\[ \num{1500} \text{ fiches} \times 2 \text{ Ko} = \num{3000} \text{ Ko} = \frac{\num{3000}}{\num{1024}} \approx 2,93 \text{ Mo.} \]
La clé USB de 4 Go ($= \num{4096}$ Mo) est très largement suffisante : elle pourrait contenir plus de mille années d'archives à ce rythme !

\textbf{Étape 7 — Texte de justification (exemple attendu).} « Chère maman, nous avons rangé tes ventes dans un tableau où chaque colonne a un rôle précis : la date, le produit, la quantité, le prix et le montant. Ainsi, l'ordinateur peut calculer et trier automatiquement. Les fichiers sont rangés dans des dossiers par mois et par produit, avec des noms sans espace ni accent pour éviter les erreurs. Dans la mémoire de l'ordinateur, le nombre 51 s'écrit 110011 en binaire, et le mot TOMATES est remplacé par des nombres grâce au code ASCII. Toutes tes fiches de l'année n'occuperont que 3 Mo environ : ta clé USB suffit largement. Le cahier numérique évite les ratures, les pertes et les erreurs de calcul. »
\end{exoresolu}

\begin{attention}[Erreurs fréquentes à signaler lors de la correction]
\begin{itemize}
\item Lire les restes de la division \emph{de haut en bas} au lieu de bas en haut : on obtiendrait $(110011)_2$ lu à l'envers, c'est-à-dire une erreur si le nombre n'est pas un palindrome. Toujours vérifier par la somme des puissances de 2.
\item Confondre Ko et Mo : diviser par \num{1024}, pas par \num{1000}.
\item Utiliser des espaces ou des accents dans les noms de fichiers (\texttt{ventes piment janvier.txt} est à proscrire).
\item Oublier de justifier le type de chaque champ : c'est pourtant ce qui permet à l'ordinateur de traiter correctement les données.
\end{itemize}
\end{attention}

\courssub{Bilan de l'activité}

Cette activité a permis de mettre en évidence, à partir d'une situation authentique du marché Mokolo, la \textbf{chaîne complète de traitement des données} étudiée dans la leçon :

\begin{itemize}
\item des \textbf{données brutes} (notes manuscrites désordonnées) ne deviennent de véritables \textbf{informations} exploitables qu'après avoir été \textbf{structurées} en champs et en enregistrements, chaque champ recevant un \textbf{type de données} adapté ;
\item l'\textbf{organisation en arborescence} de dossiers et de fichiers, avec des règles de nommage strictes, garantit que l'on retrouve rapidement une information, même en grande quantité ;
\item à l'intérieur de l'ordinateur, \textbf{toute donnée est codée en binaire} : les nombres par conversion en base 2, les caractères grâce à des tables comme le code ASCII ;
\item enfin, le calcul de l'\textbf{espace de stockage} montre concrètement l'intérêt des unités de mesure de l'information (octet, Ko, Mo, Go) pour anticiper les besoins matériels.
\end{itemize}

Tu as ainsi mobilisé les deux savoir-faire officiels de la leçon : \emph{appliquer les notions de l'unité à une situation concrète de la vie courante} et \emph{rédiger et justifier une démarche}. Ces compétences te seront directement utiles pour organiser tes propres fichiers (photos, devoirs, musique) et pour comprendre comment les commerçants, les écoles et les administrations camerounaises gèrent leurs données au quotidien.

\newpage

\coursec{Méthodes et exercices résolus}

\coursec{Organisation des données en informatique}

\courssub{Méthodes et exercices résolus}

\begin{methode}[Différencier donnée, information et identifier le type de donnée]
    \textbf{Objectif :} Savoir distinguer la donnée brute de l'information traitée et classer une donnée selon sa nature (type).
    \begin{enumerate}
        \item \textbf{Identifier la source :} Observer le fait brut, non interprété (chiffre, mesure, symbole). C'est la \cle{donnée}.
        \item \textbf{Analyser le traitement :} Repérer l'opération qui donne du sens (contexte, structure, calcul). Le résultat est l'\cle{information}.
        \item \textbf{Classifier le type :} 
        \begin{itemize}
            \item \cle{Numérique} : Entier (ex: âge, nombre d'élèves) ou Réel (ex: note, température, prix).
            \item \cle{Textuelle (Chaîne de caractères)} : Suite de symboles (ex: nom, adresse, matricule).
            \item \cle{Booléenne} : Deux valeurs seulement (Vrai/Faux, Oui/Non, 1/0) (ex: admis/recalé, abonné/non-abonné).
            \item \cle{Date/Heure} : Repère temporel structuré.
        \end{itemize}
        \item \textbf{Justifier :} Écrire une phrase du type : « C'est une donnée de type \emph{entier} car elle représente un comptage exact sans virgule. »
    \end{enumerate}
\end{methode}

\begin{exoresolu}[Analyse d'une fiche d'inscription au BEPC]
    Un établissement scolaire enregistre pour chaque candidat : le \texttt{Matricule} (ex: \texttt{2025CE00123}), le \texttt{Nom}, la \texttt{Date de naissance}, le \texttt{Sexe} (M/F), la \texttt{Moyenne\_3ème} (ex: 12,50), et le \texttt{Statut\_Bourse} (Oui/Non).
    \begin{enumerate}
        \item Pour chaque rubrique, précisez s'il s'agit d'une donnée ou d'une information au moment de la saisie.
        \item Indiquez le type informatique le plus adapté pour stocker chaque rubrique.
        \item Le secrétariat calcule le \texttt{Pourcentage\_Reussite} de l'établissement. Est-ce une donnée ou une information ? Justifiez.
    \end{enumerate}
    \tcblower
    \textbf{Solution détaillée}
    \begin{enumerate}
        \item \textbf{Au moment de la saisie}, toutes les rubriques (\texttt{Matricule}, \texttt{Nom}, \texttt{Date de naissance}, \texttt{Sexe}, \texttt{Moyenne\_3ème}, \texttt{Statut\_Bourse}) sont des \textbf{données brutes}. Elles deviennent des \textbf{informations} seulement après traitement (ex: tri par moyenne, édition de la liste des boursiers, statistiques par sexe).
        \item \textbf{Types adaptés :}
        \begin{center}
        \begin{tabularx}{\linewidth}{|l|X|l|}
        \hline
        \rowcolor{popblueL}
        \textbf{Rubrique} & \textbf{Type informatique} & \textbf{Justification} \\ \hline
        Matricule & Chaîne de caractères (String) & Contient des lettres et des chiffres, pas de calcul mathématique. \\ \hline
        Nom & Chaîne de caractères (String) & Texte pur. \\ \hline
        Date de naissance & Date (DateTime) & Structure Jour/Mois/Année, permet calcul d'âge. \\ \hline
        Sexe & Chaîne (Char) ou Booléen & Valeurs figées « M » / « F ». Booléen possible (Homme=Vrai). \\ \hline
        Moyenne\_3ème & Réel (Float/Double) & Nombre à virgule (note sur 20). \\ \hline
        Statut\_Bourse & Booléen & Deux états exclusifs : Oui (Vrai) / Non (Faux). \\ \hline
        \end{tabularx}
        \end{center}
        \item Le \texttt{Pourcentage\_Reussite} est une \textbf{information}. Il résulte d'un \emph{traitement} (calcul : $\frac{\text{Nb admis}}{\text{Nb présents}} \times 100$) sur l'ensemble des données brutes (les notes de chaque élève). Il a un sens global pour la prise de décision (évaluation de l'établissement).
    \end{enumerate}
    \textbf{Conclusion :} La distinction se fait par la présence ou non d'un traitement donnant du sens ; le choix du type dépend de la nature des valeurs et des opérations prévues (calcul, tri, affichage).
\end{exoresolu}

\begin{methode}[Organiser des données de façon structurée (Fichier, Enregistrement, Champ)]
    \textbf{Objectif :} Modéliser un ensemble de données cohérent en identifiant les entités, les attributs (champs) et la clé primaire.
    \begin{enumerate}
        \item \textbf{Identifier l'entité :} Quel objet du monde réel gère-t-on ? (Ex: \emph{Élève}, \emph{Livre}, \emph{Produit}). Cela donne le nom du \cle{fichier} (ou table).
        \item \textbf{Lister les attributs (champs) :} Quelles caractéristiques faut-il retenir pour chaque occurrence ? Définir le \textbf{nom}, le \textbf{type} et la \textbf{taille max} pour chaque champ.
        \item \textbf{Choisir la clé primaire :} Repérer l'attribut (ou groupe d'attributs) \textbf{unique} et \textbf{non nul} pour chaque occurrence (ex: Matricule, ISBN, Code\_Produit). C'est l'identifiant unique (\cle{Clé primaire}).
        \item \textbf{Représenter la structure :} Sous forme de dictionnaire des données ou schéma : \texttt{NomFichier (CléPrimaire, Champ2, Champ3, ...)}.
        \item \textbf{Respecter l'atomicité :} Un champ ne doit contenir qu'une seule valeur élémentaire (pas de « Nom et Prénom » dans un seul champ, pas de liste de notes dans un champ).
    \end{enumerate}
\end{methode}

\begin{exoresolu}[Modélisation d'une bibliothèque scolaire]
    Le bibliothécaire du Lycée de Biyem-Assi veut informatiser la gestion des ouvrages. Pour chaque livre, on note : un code unique (ex: \texttt{LIV-001}), le titre, l'auteur, l'année d'édition, la catégorie (Roman, Sciences, Histoire), le nombre d'exemplaires disponibles.
    \begin{enumerate}
        \item Proposez une structure de fichier adaptée (nom, champs, types, clé primaire).
        \item Donnez un exemple d'enregistrement (une ligne) pour l'ouvrage « L'Enfant noir » de Camara Laye, édité en 1953, catégorie Roman, 3 exemplaires, code \texttt{LIV-045}.
        \item Pourquoi le « Titre » ne peut-il pas servir de clé primaire ?
    \end{enumerate}
    \tcblower
    \textbf{Solution détaillée}
    \begin{enumerate}
        \item \textbf{Structure du fichier :}
        \texttt{LIVRE (Code\_Livre, Titre, Auteur, Annee\_Edition, Categorie, Nb\_Exemplaires)}
        \begin{itemize}
            \item \texttt{Code\_Livre} : Chaîne (Clé Primaire) — Identifiant unique.
            \item \texttt{Titre} : Chaîne (ex: 100 caractères).
            \item \texttt{Auteur} : Chaîne (ex: 50 caractères).
            \item \texttt{Annee\_Edition} : Entier (4 chiffres).
            \item \texttt{Categorie} : Chaîne (ex: 20 caractères) ou Énumération.
            \item \texttt{Nb\_Exemplaires} : Entier (positif ou nul).
        \end{itemize}
        \item \textbf{Enregistrement exemple :}
        \texttt{LIV-045 | L'Enfant noir | Camara Laye | 1953 | Roman | 3}
        \item Le \texttt{Titre} n'est \textbf{pas une clé primaire} car :
        \begin{itemize}
            \item \textbf{Non-unicité :} Plusieurs livres différents peuvent avoir le même titre (ex: « Anthologie de la poésie camerounaise » tome 1 et tome 2, ou rééditions).
            \item \textbf{Instabilité :} Le titre peut changer (nouvelle édition, correction de coquille), une clé primaire ne doit jamais changer.
            \item \textbf{Non-atomicité potentielle :} Titre + Sous-titre peut être long et variable.
        \end{itemize}
        Seule \texttt{Code\_Livre} garantit l'unicité absolue et la stabilité.
    \end{enumerate}
    \textbf{Conclusion :} Une bonne organisation structurelle repose sur l'atomicité des champs et le choix d'une clé primaire technique, stable et unique.
\end{exoresolu}

\begin{methode}[Convertir un entier décimal positif en binaire (et inversement)]
    \textbf{Objectif :} Maîtriser la codification des nombres entiers en base 2 (langage machine).
    \begin{enumerate}
        \item \textbf{Décimal $\to$ Binaire (Division successive par 2) :}
        \begin{enumerate}
            \item Diviser le nombre $N$ par 2. Noter le \textbf{quotient} $Q$ et le \textbf{reste} $R$ (0 ou 1).
            \item Recommencer avec $Q$ tant que $Q \neq 0$.
            \item Le binaire se lit \textbf{de bas en haut} (du dernier reste au premier) : $(R_{dernier} ... R_1 R_0)_2$.
        \end{enumerate}
        \item \textbf{Binaire $\to$ Décimal (Puissances de 2) :}
        \begin{enumerate}
            \item Numéroter les bits de droite à gauche en partant de 0 (poids $2^0, 2^1, 2^2...$).
            \item Multiplier chaque bit par $2^{\text{position}}$.
            \item Faire la somme : $\sum \text{bit}_i \times 2^i$.
        \end{enumerate}
        \item \textbf{Vérification rapide :} Un nombre pair finit par 0 en binaire, impair par 1. $2^n$ s'écrit 1 suivi de $n$ zéros.
    \end{enumerate}
\end{methode}

\begin{exoresolu}[Codage binaire des effectifs d'une classe]
    Une classe de 3ème compte 47 élèves. Le proviseur veut connaître le codage binaire de cet effectif pour configurer un tableau d'affichage numérique à 8 segments (1 octet).
    \begin{enumerate}
        \item Convertir l'entier 47 en binaire sur 8 bits (1 octet).
        \item L'affichage reçoit le code binaire \texttt{00101101}. Quel effectif cela représente-t-il ?
        \item Combien d'élèves au maximum peut-on afficher avec 8 bits ? Justifiez.
    \end{enumerate}
    \tcblower
    \textbf{Solution détaillée}
    \begin{enumerate}
        \item \textbf{47 (décimal) $\to$ Binaire (Division par 2) :}
        \begin{center}
        \begin{tabular}{c|c|c}
        \hline
        \textbf{Division} & \textbf{Quotient} & \textbf{Reste} \\ \hline
        $47 / 2$ & 23 & \textbf{1} (LSB - bit de poids faible) \\
        $23 / 2$ & 11 & \textbf{1} \\
        $11 / 2$ & 5  & \textbf{1} \\
        $5 / 2$  & 2  & \textbf{1} \\
        $2 / 2$  & 1  & \textbf{0} \\
        $1 / 2$  & 0  & \textbf{1} (MSB - bit de poids fort) \\ \hline
        \end{tabular}
        \end{center}
        Lecture des restes de bas en haut : \textbf{101111}.
        Sur 8 bits (complété par des zéros à gauche) : \texttt{00101111}.
        \item \textbf{Binaire \texttt{00101101} $\to$ Décimal :}
        \begin{align*}
        \texttt{0}&\times2^7 + \texttt{0}\times2^6 + \texttt{1}\times2^5 + \texttt{0}\times2^4 + \texttt{1}\times2^3 + \texttt{1}\times2^2 + \texttt{0}\times2^1 + \texttt{1}\times2^0 \\
        &= 0 + 0 + 32 + 0 + 8 + 4 + 0 + 1 \\
        &= \textbf{45}
        \end{align*}
        L'effectif affiché est \textbf{45 élèves}.
        \item \textbf{Maximum sur 8 bits :} $2^8 - 1 = 256 - 1 = \textbf{255}$.
        Justification : Avec $n$ bits, on peut coder $2^n$ valeurs distinctes (de 0 à $2^n-1$). Pour $n=8$, valeurs de 0 à 255. 47 et 45 sont bien dans cet intervalle.
    \end{enumerate}
    \textbf{Conclusion :} La méthode des divisions successives (euclidiennes) est systématique pour l'encodage ; la somme des puissances de 2 pour le décodage. Attention à l'ordre de lecture des restes (LSB vers MSB).
\end{exoresolu}

\begin{methode}[Calculer la taille de stockage d'un fichier texte ou image]
    \textbf{Objectif :} Estimer l'espace mémoire nécessaire (octets, Kio, Mio, Gio) en fonction du codage utilisé.
    \begin{enumerate}
        \item \textbf{Texte (Codage caractères) :}
        \begin{itemize}
            \item \cle{ASCII / ISO-8859-1 / Windows-1252} : 1 octet (8 bits) par caractère $\to$ Taille = Nb\_caractères $\times$ 1 octet.
            \item \cle{UTF-8} : Variable. 1 octet pour ASCII standard (A-Z, 0-9), 2 octets pour accents (é, è, à), 3 octets pour sinogrammes/arabes, 4 pour émojis. \textbf{Approximation courante en 3ème :} 1 à 2 octets/caractère pour du français.
            \item \cle{UTF-16} : 2 octets fixes (parfois 4) par caractère.
        \end{itemize}
        \item \textbf{Image (Bitmap brut non compressé) :}
        \begin{itemize}
            \item Résolution : $L \times H$ pixels (Largeur $\times$ Hauteur).
            \item Profondeur (Couleurs) : Nb bits par pixel (ex: 1 bit = N\&B, 8 bits = 256 couleurs/niveaux de gris, 24 bits = Couleurs vraies / \cle{True Color}).
            \item Taille brute (bits) = $L \times H \times \text{Profondeur}$.
            \item Taille (Octets) = $\frac{\text{Taille (bits)}}{8}$.
        \end{itemize}
        \item \textbf{Conversion d'unités (Puissances de 2 - Norme CEI/IEC) :}
        \begin{align*}
        1 \text{ Kio (Kibioctet)} &= 2^{10} = 1\,024 \text{ octets} \\
        1 \text{ Mio (Mébioctet)} &= 2^{20} = 1\,048\,576 \text{ octets} \\
        1 \text{ Gio (Gibioctet)} &= 2^{30} \text{ octets}
        \end{align*}
        \emph{Note : Les constructeurs de disques durs utilisent souvent les puissances de 10 (Ko = 1000 o). Au BEPC, on utilise les puissances de 2 (Kio/Mio) pour la mémoire vive et les tailles de fichiers OS.}
    \end{enumerate}
\end{methode}

\begin{exoresolu}[Taille d'un fichier liste d'élèves et d'une photo de classe]
    \begin{enumerate}
        \item Un fichier texte \texttt{liste\_3eme.txt} contient 120 lignes. Chaque ligne : « Matricule Nom Prénom Moyenne » (en moyenne 35 caractères par ligne, accents compris). Le fichier est encodé en \textbf{UTF-8}. Estimer la taille en Kio.
        \item La photo de classe fait $3000 \times 2000$ pixels, en couleur vraie (24 bits/pixel), non compressée (format BMP brut). Calculer sa taille exacte en Mio.
    \end{enumerate}
    \tcblower
    \textbf{Solution détaillée}
    \begin{enumerate}
        \item \textbf{Fichier texte (UTF-8) :}
        \begin{itemize}
            \item Nombre total de caractères $\approx 120 \times 35 = 4\,200$ caractères.
            \item En UTF-8, les caractères ASCII (chiffres, lettres sans accent, espaces, ponctuation) = 1 octet. Les accents (é, è, ê, à, ô...) = 2 octets.
            \item Estimation prudente : $\approx 10\%$ d'accents $\to$ $420 \times 2 + 3780 \times 1 = 4\,620$ octets. Arrondissons à \textbf{4 700 octets}.
            \item Conversion en Kio : $\frac{4\,700}{1\,024} \approx \textbf{4,59 Kio}$.
        \end{itemize}
        \item \textbf{Image BMP (24 bits/pixel) :}
        \begin{align*}
        \text{Nb pixels} &= 3\,000 \times 2\,000 = 6\,000\,000 \text{ pixels} \\
        \text{Taille (bits)} &= 6\,000\,000 \times 24 = 144\,000\,000 \text{ bits} \\
        \text{Taille (octets)} &= \frac{144\,000\,000}{8} = 18\,000\,000 \text{ octets} \\
        \text{Taille (Mio)} &= \frac{18\,000\,000}{1\,048\,576} \approx \textbf{17,17 Mio}
        \end{align*}
        (Souvent arrondi à 17,2 Mio ou 18 Mo si puissance de 10).
    \end{enumerate}
    \textbf{Conclusion :} Le texte est léger (quelques Kio). L'image brute non compressée est lourde (dizaines de Mio), d'où l'intérêt de la compression (JPG, PNG) pour le stockage et le transfert (ex: envoi par WhatsApp ou email au Cameroun).
\end{exoresolu}

\begin{methode}[Comprendre le codage des caractères : ASCII, Unicode (UTF-8)]
    \textbf{Objectif :} Expliquer pourquoi l'ASCII est insuffisant et comment l'Unicode (UTF-8) résout le problème des langues du monde.
    \begin{enumerate}
        \item \textbf{ASCII (7 bits $\to$ 128 codes) :} Historiquement standard US. Codes 0-31 (contrôle), 32-127 (imprimables : A-Z, a-z, 0-9, symboles). \textbf{Problème :} Pas d'accents (é, ç, à), pas de caractères africains (ŋ, ɔ, Ɛ), pas d'arabe, chinois, émojis.
        \item \textbf{ISO-8859-1 (Latin-1, 8 bits = 256 codes) :} Étend l'ASCII pour l'Europe de l'Ouest (ajoute é, è, à, ç, ñ...). \textbf{Problème :} Incompatible entre zones géographiques (Latin-2 pour l'Est, Latin-6 pour l'Arabe...). Un fichier lu avec le mauvais encodage affiche des « \texttt{�} » (mojibake).
        \item \textbf{Unicode (Standard universel) :} Attribue un \textbf{point de code} unique (entier) à \emph{chaque} caractère de \emph{toutes} les écritures (U+0000 à U+10FFFF). Ex: 'A'=U+0041, 'é'=U+00E9, 'Ɛ'=U+0190, '😀'=U+1F600.
        \item \textbf{UTF-8 (Encodage Unicode dominant sur le Web/OS) :} Codage \textbf{à longueur variable} (1 à 4 octets).
        \begin{itemize}
            \item 1 octet (0xxxxxxx) : ASCII pur (compatibilité totale ascendante).
            \item 2 octets (110xxxxx 10xxxxxx) : Alphabets étendus (Latin, Grec, Cyrillique, Arabe, Hébreu, \textbf{Caractères africains spéciaux}).
            \item 3 octets : Sinogrammes (Chinois, Japonais, Coréen - BMP).
            \item 4 octets : Caractères rares, Émojis, Écritures historiques.
        \end{itemize}
        \item \textbf{Règle d'or :} Toujours déclarer/sauvegarder en \textbf{UTF-8} (éditeur de code, base de données, page HTML \texttt{<meta charset="UTF-8">}).
    \end{enumerate}
\end{methode}

\begin{exoresolu}[Diagnostic d'un problème d'affichage « Mojibake »]
    Un élève de 3ème au Lycée de Ngaoundéré ouvre un fichier \texttt{notes.csv} reçu par email. Il voit : \texttt{Nom : KOUAM\u{e} ; Note : 14/20} au lieu de \texttt{KOUAMÉ}.
    \begin{enumerate}
        \item Expliquez l'origine technique de ce défaut d'affichage.
        \item Le fichier a été créé sur un vieux système Windows (encodage par défaut Windows-1252 / ANSI) et lu sur un système moderne (Linux/Windows 10/11) qui attend de l'UTF-8. Donnez le code hexadécimal UTF-8 correct pour le caractère « É » (U+00C9).
        \item Proposez deux solutions concrètes pour l'élève afin de lire correctement son fichier.
    \end{enumerate}
    \tcblower
    \textbf{Solution détaillée}
    \begin{enumerate}
        \item \textbf{Origine :} C'est un classique \textbf{conflit d'encodage (Mojibake)}. Le fichier a été encodé en \textbf{Windows-1252} (1 octet par caractère étendu). Le caractère « É » y a pour code hexadécimal \texttt{C9} (1 octet). Le lecteur moderne (UTF-8) lit l'octet \texttt{C9} comme le début d'une séquence sur 2 octets (car \texttt{C9} = \texttt{11001001}$_2$ commence par \texttt{110}). Il attend un second octet de continuation (\texttt{10xxxxxx}) qui n'existe pas (ou est le caractère suivant « ; »). Ne trouvant pas de séquence valide, il affiche le caractère de remplacement \texttt{�} (U+FFFD) ou « \u{e} » selon la police.
        \item \textbf{Code UTF-8 de « É » (U+00C9) :}
        U+00C9 est dans l'intervalle U+0080..U+07FF $\to$ Codage sur \textbf{2 octets}.
        Binaire du point de code : \texttt{000 1100 1001} (11 bits utiles).
        Répartition UTF-8 : \texttt{110}xxxxx \texttt{10}xxxxxx
        $\to$ \texttt{110}00011 \texttt{10}001001
        $\to$ Hexadécimal : \texttt{C3 89}.
        \textbf{Vérification :} \texttt{C3} = \texttt{11000011} (début 2 octets), \texttt{89} = \texttt{10001001} (continuation). Correct.
        \item \textbf{Solutions pour l'élève :}
        \begin{enumerate}
            \item \textbf{Ouvrir avec un éditeur avancé (Notepad++, VS Code, Sublime Text) :} Menu \texttt{Encodage} $\to$ \texttt{Character Sets} $\to$ \texttt{Western European} $\to$ \texttt{Windows-1252}. Le texte s'affiche bien. Puis \texttt{Encodage} $\to$ \texttt{Convertir en UTF-8} et \textbf{Sauvegarder}.
            \item \textbf{En ligne de commande (Linux/Termux/WSL) :} \texttt{iconv -f WINDOWS-1252 -t UTF-8 notes.csv > notes\_utf8.csv}.
        \end{enumerate}
    \end{enumerate}
    \textbf{Conclusion :} L'UTF-8 est le standard universel. Tout fichier texte échangé (email, web, clé USB) doit être en UTF-8 pour éviter ce problème, surtout dans un pays multilingue comme le Cameroun (français, anglais, langues nationales).
\end{exoresolu}

\begin{methode}[Passer de l'algorithme au code C : Structure de base, Entrées/Sorties, Types]
    \textbf{Objectif :} Écrire un programme C complet, correct, compilable (GCC), respectant la syntaxe stricte pour résoudre un problème simple (calcul, conversion).
    \begin{enumerate}
        \item \textbf{Inclusions :} \texttt{\#include <stdio.h>} (indispensable pour \texttt{printf}, \texttt{scanf}).
        \item \textbf{Fonction principale :} \texttt{int main(void) \{ ... return 0; \}}.
        \item \textbf{Déclaration des variables :} Au début du bloc (C90) ou près de l'usage (C99). Types de base : \texttt{int} (entier), \texttt{float} / \texttt{double} (réel), \texttt{char} (caractère/octet).
        \item \textbf{Affichage (Sortie) :} \texttt{printf("Format", variables);} 
        \begin{itemize}
            \item \texttt{\%d} ou \texttt{\%i} pour \texttt{int}.
            \item \texttt{\%f} pour \texttt{float} (par défaut 6 décimales), \texttt{\%.2f} pour 2 décimales.
            \item \texttt{\%c} pour \texttt{char}, \texttt{\%s} pour chaîne.
            \item \texttt{\textbackslash n} pour retour à la ligne.
        \end{itemize}
        \item \textbf{Saisie (Entrée) :} \texttt{scanf("Format", \&variable);} 
        \begin{itemize}
            \item \textbf{Obligatoire :} Le \textbf{\&} (esperluette / adresse de) devant le nom de la variable (sauf pour les tableaux/chaînes).
            \item Le format doit correspondre exactement au type.
        \end{itemize}
        \item \textbf{Opérateurs arithmétiques :} \texttt{+ - * / \%} (modulo, reste entier). Attention : \texttt{/} sur deux \texttt{int} fait une \textbf{division entière} (tronquée). Forcer le réel : \texttt{(float)a / b} ou \texttt{a / 2.0}.
        \item \textbf{Compilation / Exécution (Terminal) :} \texttt{gcc -Wall -o prog prog.c} puis \texttt{./prog}.
    \end{enumerate}
\end{methode}

\begin{exoresolu}[Programme C : Conversion Binaire $\to$ Décimal]
    Écrire un programme C complet qui :
    \begin{itemize}
        \item Demande à l'utilisateur de saisir un \textbf{entier positif} représentant un nombre binaire (ex: \texttt{101111}).
        \item Vérifie que la saisie ne contient que des 0 et des 1 (optionnel mais recommandé).
        \item Calcule et affiche l'équivalent décimal.
    \end{itemize}
    \textit{Indication : Traiter le nombre saisi comme un entier (type \texttt{long long}) et extraire les chiffres par modulo 10 et division par 10.}
    \tcblower
    \textbf{Solution détaillée (Code C commenté)}
    \begin{lstlisting}[language=C]
#include <stdio.h>
#include <stdbool.h> // Pour le type bool, true, false

int main(void) {
    long long binaire; // Type assez grand pour stocker jusqu'à 18-19 chiffres binaires (64 bits)
    int decimal = 0;
    int poids = 1; // 2^0 = 1 au départ
    bool valide = true;

    // 1. Saisie
    printf("Entrez un nombre binaire (ex: 101111) : ");
    // On lit comme un entier. Si l'utilisateur tape 102, binaire vaudra 102.
    if (scanf("%lld", &binaire) != 1) {
        printf("Erreur de saisie.\n");
        return 1;
    }

    // 2. Validation et Conversion simultanées
    long long temp = binaire; // Copie pour ne pas perdre la valeur originale si besoin
    
    if (temp == 0) {
        decimal = 0; // Cas particulier : 0 binaire = 0 decimal
    } else {
        while (temp > 0) {
            int bit = temp % 10; // Extrait le chiffre de droite (poids faible)
            
            // Vérification : un chiffre binaire ne peut être que 0 ou 1
            if (bit != 0 && bit != 1) {
                valide = false;
                break; // Sort de la boucle immédiatement
            }
            
            // Calcul décimal : bit * 2^position
            // Ici 'poids' vaut 1, 2, 4, 8, 16... (puissances de 2 successives)
            decimal += bit * poids;
            
            // Préparation itération suivante
            poids *= 2;      // Poids suivant (x2)
            temp /= 10;      // Décale le nombre vers la droite (supprime le chiffre traité)
        }
    }

    // 3. Affichage résultat
    if (valide) {
        printf("Binaire saisi : %lld\n", binaire);
        printf("Equivalent decimal : %d\n", decimal);
    } else {
        printf("Erreur : '%lld' n'est pas un nombre binaire valide (que des 0 et 1).\n", binaire);
    }

    return 0; // Succès
}
    \end{lstlisting}
    \textbf{Explications clés :}
    \begin{itemize}
        \item \texttt{long long} et \texttt{\%lld} : Nécessaires car un binaire comme \texttt{111111111111111111} (18 uns) dépasse la capacité d'un \texttt{int} standard (souvent 32 bits max $\approx$ 2 milliards, mais $10^{18} > 2\times 10^{9}$). \texttt{long long} garantit 64 bits.
        \item \texttt{temp \% 10} : Récupère le dernier chiffre (le bit de poids faible courant).
        \item \texttt{poids *= 2} : Remplace le calcul de $2^i$ par une multiplication itérative plus efficace.
        \item \texttt{temp /= 10} : Décale le nombre pour traiter le bit suivant à la boucle d'après.
        \item Gestion du cas \texttt{0} explicite car la boucle \texttt{while(temp > 0)} ne s'exécuterait pas.
    \end{itemize}
    \textbf{Test mental :} Saisie \texttt{101111} (47).
    Boucle 1: bit=1, dec=1, poids=2, temp=10111.
    Boucle 2: bit=1, dec=3, poids=4, temp=1011.
    Boucle 3: bit=1, dec=7, poids=8, temp=101.
    Boucle 4: bit=1, dec=15, poids=16, temp=10.
    Boucle 5: bit=0, dec=15, poids=32, temp=1.
    Boucle 6: bit=1, dec=47, poids=64, temp=0. Fin. Affichage 47. \textbf{Correct.}
\end{exoresolu}

\begin{methode}[Résoudre un problème concret par l'algorithmique (Analyse, Algo, Codage, Test)]
    \textbf{Objectif :} Appliquer la démarche complète de résolution informatique à une situation de la vie courante (gestion de notes, stock, budget).
    \begin{enumerate}
        \item \textbf{Analyser la situation :} Identifier les \textbf{données d'entrée} (connues, saisies), les \textbf{résultats attendus} (sorties), et les \textbf{contraintes/règles de gestion} (formules, seuils, arrondis).
        \item \textbf{Concevoir l'algorithme (Pseudo-code) :} 
        \begin{itemize}
            \item Variables (Noms clairs, Types).
            \item Séquence logique : \texttt{Début} $\to$ \texttt{Saisir} $\to$ \texttt{Traiter (Si/Pour/TantQue)} $\to$ \texttt{Afficher} $\to$ \texttt{Fin}.
            \item Indentation rigoureuse (2 ou 4 espaces) pour voir les blocs.
        \end{itemize}
        \item \textbf{Coder en C (ou Python/JS si requis) :} Transcrire fidèlement l'algorithme. Respecter la syntaxe (\texttt{;}, \{\}, \texttt{\&} pour scanf).
        \item \textbf{Tester (Jeu de tests) :} 
        \begin{itemize}
            \item Cas \textbf{nominal} (valeurs normales).
            \item Cas \textbf{limites} (0, valeur max, seuil exact).
            \item Cas \textbf{erronés/invalides} (négatif, texte au lieu de nombre, division par 0).
        \end{itemize}
        \item \textbf{Documenter / Conclure :} Commenter le code (// but de la variable, // formule). Rédiger la phrase de réponse finale.
    \end{enumerate}
\end{methode}

\begin{exoresolu}[Gestion du stock d'un cybercafé (Algorithme + C)]
    Le gérant du cybercafé « @BiyemAssi » veut un programme pour gérer les recharges de temps. 
    Règles : 
    \begin{itemize}
        \item 1 heure = 300 FCFA.
        \item Si le client achète $\ge$ 5 heures, il bénéficie de 10\% de remise sur le total.
        \item Le programme demande le nombre d'heures souhaité (entier $> 0$), calcule et affiche le montant à payer (entier FCFA, arrondi au plus bas).
    \end{itemize}
    \begin{enumerate}
        \item Écrire l'algorithme en pseudo-code français structuré.
        \item Écrire le programme C correspondant.
        \item Donner les résultats pour 3h, 5h, 10h.
    \end{enumerate}
    \tcblower
    \textbf{Solution détaillée}
    \begin{enumerate}
        \item \textbf{Algorithme (Pseudo-code) :}
        \begin{lstlisting}
Algorithme CalculRechargeCyber
Variables
    nbHeures : Entier
    tarifHoraire, total, remise, aPayer : Reel (* ou Entier pour FCFA *)
Debut
    Afficher "Bienvenue au Cyber @BiyemAssi"
    Afficher "Tarif : 300 FCFA/h. Remise 10% si >= 5h."
    Repeter
        Afficher "Nombre d'heures desirees (entier > 0) : "
        Saisir nbHeures
    JusquA nbHeures > 0
    
    tarifHoraire <- 300
    total <- nbHeures * tarifHoraire
    
    Si nbHeures >= 5 Alors
        remise <- total * 0.10
    Sinon
        remise <- 0
    FinSi
    
    aPayer <- total - remise
    // Arrondi au plus bas (entier FCFA) : conversion implicite Entier ou floor()
    Afficher "Total brut : ", total, " FCFA"
    Afficher "Remise : ", remise, " FCFA"
    Afficher "Montant a payer : ", PartieEntiere(aPayer), " FCFA"
Fin
        \end{lstlisting}
        \item \textbf{Programme C :}
        \begin{lstlisting}[language=C]
#include <stdio.h>

int main(void) {
    int nbHeures;
    const int TARIF = 300; // Constante pour le tarif
    int total, remise, aPayer; // Tout en entiers (FCFA) pour éviter erreurs float

    printf("=== Cyber @BiyemAssi ===\n");
    printf("Tarif : %d FCFA/h. Remise 10%% si >= 5h.\n", TARIF);

    // Saisie sécurisée (boucle tant que invalide)
    do {
        printf("Nombre d'heures desirees (entier > 0) : ");
        if (scanf("%d", &nbHeures) != 1) { // Gestion saisie non-numerique
            printf("Erreur : saisir un nombre entier.\n");
            while(getchar() != '\n'); // Vider le buffer clavier
            nbHeures = 0; // Forcer reprise boucle
        }
    } while (nbHeures <= 0);

    total = nbHeures * TARIF;

    // Calcul remise en entier : (total * 10) / 100 = total / 10
    // Division entière = arrondi au plus bas automatique en C pour positifs
    if (nbHeures >= 5) {
        remise = total / 10; 
    } else {
        remise = 0;
    }

    aPayer = total - remise;

    printf("------------------------\n");
    printf("Total brut   : %d FCFA\n", total);
    printf("Remise (10%%) : %d FCFA\n", remise);
    printf("A PAYER     : %d FCFA\n", aPayer);

    return 0;
}
        \end{lstlisting}
        \item \textbf{Jeu de tests (Vérification manuelle) :}
        \begin{center}
        \begin{tabular}{|c|c|c|c|c|}
        \hline
        \rowcolor{popblueL}
        \textbf{Heures} & \textbf{Total Brut} & \textbf{Condition} & \textbf{Remise (Total/10)} & \textbf{À Payer} \\ \hline
        3 h & $3 \times 300 = 900$ & $< 5$ & 0 & \textbf{900 FCFA} \\ \hline
        5 h & $5 \times 300 = 1500$ & $\ge 5$ & $1500 / 10 = 150$ & \textbf{1350 FCFA} \\ \hline
        10 h & $10 \times 300 = 3000$ & $\ge 5$ & $3000 / 10 = 300$ & \textbf{2700 FCFA} \\ \hline
        \end{tabular}
        \end{center}
        Le programme C produit exactement ces résultats. L'utilisation d'entiers pour l'argent (FCFA) évite les erreurs d'arrondi du type \texttt{float} (ex: 1349.9999 au lieu de 1350).
    \end{enumerate}
    \textbf{Conclusion :} La démarche structurée (Analyse $\to$ Algo $\to$ Code $\to$ Test) garantit un programme fiable. L'usage de \texttt{const} pour les constantes métier (TARIF) et d'entiers pour la monnaie sont des bonnes pratiques professionnelles.
\end{exoresolu}

\begin{attention}[Les 5 erreurs qui coûtent des points au BEPC]
    \begin{enumerate}
        \item \textbf{Confondre Donnée et Information :} Dire « La note 12 est une information » sans préciser le contexte. \textbf{Correction :} « La note 12 est une \emph{donnée brute}. La \emph{décision} "Admis" prise à partir de la moyenne est une \emph{information}. »
        \item \textbf{Oublier le \texttt{\&} dans \texttt{scanf} :} Écrire \texttt{scanf("\%d", age);} au lieu de \texttt{scanf("\%d", \&age);}. \textbf{Conséquence :} Crash immédiat (Segmentation fault) ou écriture mémoire aléatoire. \textbf{Règle :} \texttt{\&} obligatoire pour \texttt{int, float, double, char} (variables scalaires).
        \item \textbf{Division entière non voulue :} Écrire \texttt{moyenne = (note1 + note2) / 2;} avec des \texttt{int}. Résultat : 12 au lieu de 12,5. \textbf{Correction :} \texttt{moyenne = (note1 + note2) / 2.0;} ou caster \texttt{(float)(note1+note2)/2}.
        \item \textbf{Mauvaise lecture des restes (Binaire) :} Lire les restes de haut en bas (premier reste = bit de poids fort). \textbf{Règle :} Le \textbf{premier reste} est le \textbf{bit de poids faible (LSB, bit 0)}. On lit \textbf{de bas en haut}.
        \item \textbf{Unités de stockage mélangées (Puissances 10 vs 2) :} Calculer $1 \text{ Mo} = 1\,000\,000 \text{ o}$ pour une taille de fichier RAM/Disque OS. \textbf{Au BEPC Informatique :} $1 \text{ Kio} = 1024 \text{ o}$, $1 \text{ Mio} = 1024 \text{ Kio}$. Utiliser les puissances de 2 ($2^{10}, 2^{20}$). Noter l'unité \textbf{Kio/Mio/Gio} (kibi, mebi, gibi) pour être précis.
    \end{enumerate}
\end{attention}

\newpage

\coursec{Exercices}

\courssub{Je vérifie mes connaissances}

\begin{exercice}[QCM : à toi de choisir]
Pour chaque question, une seule réponse est exacte. Recopie la lettre de la bonne réponse sur ta feuille.
\begin{enumerate}
\item Une \cle{donnée brute} est :
\begin{enumerate}
\item une information déjà interprétée ;
\item un élément brut non encore traité (chiffre, mot, image) ;
\item un programme de l'ordinateur.
\end{enumerate}
\item L'unité de base de l'information dans l'ordinateur est :
\begin{enumerate}
\item le caractère ; \item le bit ; \item le kilo-octet.
\end{enumerate}
\item Un octet est formé de :
\begin{enumerate}
\item 2 bits ; \item 8 bits ; \item 16 bits.
\end{enumerate}
\item Le code \cle{ASCII} sert à représenter :
\begin{enumerate}
\item les caractères (lettres, chiffres, symboles) ;
\item uniquement les nombres entiers ;
\item les images.
\end{enumerate}
\item $1$~Ko (kilo-octet) vaut environ :
\begin{enumerate}
\item 8 octets ; \item 100 octets ; \item 1024 octets.
\end{enumerate}
\end{enumerate}
\end{exercice}

\begin{exercice}[Vrai ou faux ? Justifie]
Réponds par Vrai ou Faux, puis justifie ta réponse en une ou deux phrases.
\begin{enumerate}
\item Une photo numérique est une information, jamais une donnée.
\item Dans l'ordinateur, toutes les données sont représentées à l'aide des deux chiffres 0 et 1.
\item Le nombre binaire $1011$ contient quatre bits.
\item Le code ASCII d'une lettre majuscule est le même que celui de la même lettre en minuscule.
\item Une clé USB de $1$~Go peut contenir plus de données qu'une disquette de $1$~Mo.
\end{enumerate}
\end{exercice}

\begin{exercice}[Définitions à compléter]
Recopie et complète les phrases suivantes avec les mots : \emph{bit}, \emph{information}, \emph{donnée}, \emph{octet}, \emph{arborescence}, \emph{binaire}.
\begin{enumerate}
\item Un élément brut comme la note $14$ est une \ldots{} ; la phrase « l'élève a la moyenne » est une \ldots{}.
\item Le \ldots{} est la plus petite unité d'information ; il ne peut prendre que les valeurs 0 ou 1.
\item Huit bits forment un \ldots{}.
\item Le système de numération qui n'utilise que les chiffres 0 et 1 est le système \ldots{}.
\item L'organisation des dossiers et fichiers sous forme d'arbre s'appelle une \ldots{}.
\end{enumerate}
\end{exercice}

\begin{exercice}[Calculs directs]
\begin{enumerate}
\item Convertis en décimal les nombres binaires suivants en détaillant les puissances de 2 : $101$ ; $1100$.
\item Convertis en binaire les nombres décimaux suivants par divisions successives par 2 : $10$ ; $19$.
\item Combien d'octets y a-t-il dans $3$~Ko ? dans $2$~Mo ?
\end{enumerate}
\end{exercice}

\courssub{Je m'entraîne}

\begin{exercice}[Données brutes ou informations ?]
Voici une liste d'éléments relevés dans le secrétariat d'un collège : une note de $12/20$ ; la moyenne générale d'un élève ; une photo d'identité ; une chanson enregistrée ; un numéro de téléphone ; la mention « Admis » sur un relevé.
\begin{enumerate}
\item Classe chaque élément dans la catégorie \emph{donnée brute} ou \emph{information}.
\item Justifie chacun de tes choix en une phrase.
\end{enumerate}
\end{exercice}

\begin{exercice}[Codage ASCII du mot « ECOLE »]
On donne l'extrait suivant de la table ASCII (codes décimaux) :
\begin{center}
\begin{tabular}{|c|c|c|c|c|c|c|c|c|}
\hline
\rowcolor{popblueL} Caractère & A & B & C & E & L & N & O & R \\
\hline
Code ASCII & 65 & 66 & 67 & 69 & 76 & 78 & 79 & 82 \\
\hline
\end{tabular}
\end{center}
\begin{enumerate}
\item Écris le code ASCII décimal de chaque lettre du mot « ECOLE ».
\item Convertis en binaire le code de la lettre « E » et celui de la lettre « O ».
\item À l'aide de la table, décode le message suivant : $66 - 79 - 78 - 74 - 79 - 85 - 82$ (le code de J est $74$ et celui de U est $85$).
\end{enumerate}
\end{exercice}

\begin{exercice}[La clé USB de Marlyse]
Marlyse veut copier sur sa clé USB de $1$~Go : $300$ photos de $2$~Mo chacune et $10$ chansons de $5$~Mo chacune. On rappelle que $1$~Go $= 1024$~Mo.
\begin{enumerate}
\item Calcule l'espace occupé par les photos, puis par les chansons.
\item Calcule l'espace total nécessaire en Mo.
\item La copie est-elle possible sur la clé ? Justifie ta réponse par un calcul et une conclusion rédigée.
\end{enumerate}
\end{exercice}

\begin{exercice}[L'arborescence du lycée]
Le lycée de Nkongsamba organise ses fichiers ainsi : un dossier principal « LYCEE » contient un dossier par niveau (« 2nde », « 1ere », « Tle ») ; chaque niveau contient un dossier par classe ; chaque classe contient un dossier par matière et un dossier « BULLETINS ».
\begin{enumerate}
\item Dessine l'arborescence complète pour les classes Tle C et Tle D, avec les matières Maths et Informatique.
\item Écris le chemin d'accès complet du fichier « bulletin1.pdf » de la classe de Tle C, en commençant par « LYCEE ».
\item Cite deux avantages d'une telle organisation structurée des données.
\end{enumerate}
\end{exercice}

\courssub{Je me prépare au Probatoire}

\begin{exercice}[Type BEPC : conversions binaire--décimal]
\textbf{Partie A.} Un professeur d'informatique écrit au tableau les nombres $25$ et $47$.
\begin{enumerate}
\item Convertis $25$ en binaire par la méthode des divisions successives par 2, en présentant toutes les étapes.
\item Convertis $47$ en binaire par la même méthode.
\item Vérifie tes résultats en reconvertissant chaque nombre binaire obtenu en décimal, à l'aide des puissances de 2.
\end{enumerate}
\textbf{Partie B.} Un élève a trouvé sur une affiche le nombre binaire $10110$.
\begin{enumerate}
\item Convertis $10110$ en décimal en détaillant la méthode des puissances de 2.
\item Combien de bits composent ce nombre ? Quel est le bit de poids fort ? le bit de poids faible ?
\item Rédige une conclusion expliquant pourquoi l'ordinateur utilise le binaire plutôt que le décimal.
\end{enumerate}
\end{exercice}

\begin{exercice}[Type BEPC : le cybercafé de Bafoussam]
M. Kamdem gère un cybercafé à Bafoussam. Il possède environ $2000$ fichiers clients : des documents Word, des photos scannées et des fichiers de musique, actuellement mélangés dans un seul dossier. Chaque fichier occupe en moyenne $3$~Mo.
\begin{enumerate}
\item Explique la différence entre une donnée et une information, en donnant un exemple de chaque tiré de la situation de M. Kamdem.
\item Propose une organisation structurée de ses fichiers sous forme d'arborescence (dossiers et sous-dossiers) et dessine-la.
\item Écris le chemin d'accès d'un fichier « contrat.pdf » rangé dans le sous-dossier « Documents » de l'année 2026.
\item Calcule l'espace disque total nécessaire en Mo, puis convertis-le en Go (on rappelle : $1$~Go $= 1024$~Mo).
\item M. Kamdem hésite entre un disque dur de $4$~Go et un disque de $10$~Go. Lequel conseilles-tu ? Rédige et justifie ta démarche complète.
\end{enumerate}
\end{exercice}

\begin{exercice}[Type BEPC : le bulletin de la 3e A]
Le principal d'un collège de Yaoundé veut numériser les bulletins de la classe de 3e A, qui compte $60$ élèves. Chaque bulletin scanné est une image de $500$~Ko. On donne l'extrait de table ASCII : A $= 65$, B $= 66$, C $= 67$, E $= 69$.
\begin{enumerate}
\item Le nom du collège commence par le mot « CEB ». Donne le code ASCII décimal de chaque lettre de ce mot, puis convertis le code de la lettre « B » en binaire.
\item Calcule l'espace total occupé par les $60$ bulletins en Ko, puis convertis le résultat en Mo (on rappelle : $1$~Mo $= 1024$~Ko).
\item Le principal veut graver ces bulletins sur un CD de $700$~Mo. Est-ce possible ? Justifie par un calcul.
\item Propose une arborescence permettant de ranger les bulletins par trimestre (T1, T2, T3) et écris le chemin d'accès du bulletin de l'élève « ABENA » pour le trimestre T1.
\item Explique en quelques lignes comment l'ordinateur représente une image scannée à l'aide de 0 et de 1.
\end{enumerate}
\end{exercice}

\newpage

\coursec{Corrigés des exercices}

\begin{corrige}[Exercice 1 — QCM : à toi de choisir]
\begin{enumerate}
\item \textbf{b.} Une donnée brute est un élément brut non encore traité (chiffre, mot, image). C'est le traitement qui la transforme en information.
\item \textbf{b.} L'unité de base de l'information est le \cle{bit} (0 ou 1).
\item \textbf{b.} Un octet est formé de $8$ bits.
\item \textbf{a.} Le code ASCII sert à représenter les caractères : lettres, chiffres et symboles.
\item \textbf{c.} $1$~Ko vaut environ $1024$ octets (et non $1000$).
\end{enumerate}
\end{corrige}

\begin{corrige}[Exercice 2 — Vrai ou faux ? Justifie]
\begin{enumerate}
\item \textbf{Faux.} Une photo numérique est d'abord une \emph{donnée} (un fichier d'image stocké dans l'ordinateur) ; elle devient une \emph{information} lorsqu'elle est interprétée, par exemple pour identifier un élève.
\item \textbf{Vrai.} Dans l'ordinateur, toutes les données (textes, images, sons, nombres) sont codées en \cle{binaire}, c'est-à-dire à l'aide des deux chiffres 0 et 1, car les circuits électroniques ne connaissent que deux états (courant qui passe ou ne passe pas).
\item \textbf{Vrai.} Le nombre $1011$ est écrit avec quatre chiffres binaires ; il contient donc $4$ bits.
\item \textbf{Faux.} En ASCII, chaque caractère a un code différent : par exemple « A » vaut $65$ et « a » vaut $97$. Majuscules et minuscules sont des caractères distincts.
\item \textbf{Vrai.} $1$~Go $= 1024$~Mo, donc $1$~Go $> 1$~Mo : la clé USB contient beaucoup plus de données que la disquette (environ $1024$ fois plus).
\end{enumerate}
\end{corrige}

\begin{corrige}[Exercice 3 — Définitions à compléter]
\begin{enumerate}
\item Un élément brut comme la note $14$ est une \emph{donnée} ; la phrase « l'élève a la moyenne » est une \emph{information}.
\item Le \emph{bit} est la plus petite unité d'information ; il ne peut prendre que les valeurs 0 ou 1.
\item Huit bits forment un \emph{octet}.
\item Le système de numération qui n'utilise que les chiffres 0 et 1 est le système \emph{binaire}.
\item L'organisation des dossiers et fichiers sous forme d'arbre s'appelle une \emph{arborescence}.
\end{enumerate}
\end{corrige}

\begin{corrige}[Exercice 4 — Calculs directs]
\begin{enumerate}
\item Conversion binaire $\rightarrow$ décimal :
\begin{align*}
101 &= 1\times 2^2 + 0\times 2^1 + 1\times 2^0 = 4 + 0 + 1 = \textbf{5} \\
1100 &= 1\times 2^3 + 1\times 2^2 + 0\times 2^1 + 0\times 2^0 = 8 + 4 + 0 + 0 = \textbf{12}
\end{align*}
\item Conversion décimal $\rightarrow$ binaire par divisions successives par 2 :
\begin{itemize}
\item Pour $10$ : $10 \div 2 = 5$ reste $0$ ; $5 \div 2 = 2$ reste $1$ ; $2 \div 2 = 1$ reste $0$ ; $1 \div 2 = 0$ reste $1$. On lit les restes de bas en haut : $10 = \textbf{1010}_2$.
\item Pour $19$ : $19 \div 2 = 9$ reste $1$ ; $9 \div 2 = 4$ reste $1$ ; $4 \div 2 = 2$ reste $0$ ; $2 \div 2 = 1$ reste $0$ ; $1 \div 2 = 0$ reste $1$. Lecture de bas en haut : $19 = \textbf{10011}_2$.
\end{itemize}
Vérification : $1010_2 = 8 + 2 = 10$ et $10011_2 = 16 + 2 + 1 = 19$. Correct.
\item Conversions d'unités :
\begin{itemize}
\item $3$~Ko $= 3 \times 1024 = \textbf{3072 octets}$ ;
\item $2$~Mo $= 2 \times 1024 \times 1024 = \textbf{$2\,097\,152$ octets}$.
\end{itemize}
\end{enumerate}
\end{corrige}

\begin{corrige}[Exercice 5 — Données brutes ou informations ?]
\begin{enumerate}
\item Classement :
\begin{center}
\begin{tabularx}{\linewidth}{|l|X|}
\hline
\rowcolor{popblueL} Élément & Catégorie \\
\hline
Une note de $12/20$ & Donnée brute \\
\hline
La moyenne générale d'un élève & Information \\
\hline
Une photo d'identité & Donnée brute \\
\hline
Une chanson enregistrée & Donnée brute \\
\hline
Un numéro de téléphone & Donnée brute \\
\hline
La mention « Admis » sur un relevé & Information \\
\hline
\end{tabularx}
\end{center}
\item Justifications :
\begin{itemize}
\item La note $12/20$ est un chiffre brut, non encore interprété : c'est une donnée.
\item La moyenne générale est le résultat d'un \emph{traitement} (calcul à partir de plusieurs notes) : c'est une information.
\item La photo d'identité est un fichier image brut : c'est une donnée.
\item La chanson enregistrée est un fichier son brut : c'est une donnée.
\item Le numéro de téléphone est une suite de chiffres brute : c'est une donnée.
\item La mention « Admis » résulte de l'interprétation de la moyenne (décision prise) : c'est une information.
\end{itemize}
\end{enumerate}
\end{corrige}

\begin{corrige}[Exercice 6 — Codage ASCII du mot « ECOLE »]
\begin{enumerate}
\item D'après la table : E $= 69$, C $= 67$, O $= 79$, L $= 76$, E $= 69$.\\
Le mot « ECOLE » se code donc : $69 - 67 - 79 - 76 - 69$.
\item Conversions en binaire (sur 7 bits, code ASCII standard) :
\begin{itemize}
\item « E » $= 69$ : $69 \div 2 = 34$ reste $1$ ; $34 \div 2 = 17$ reste $0$ ; $17 \div 2 = 8$ reste $1$ ; $8 \div 2 = 4$ reste $0$ ; $4 \div 2 = 2$ reste $0$ ; $2 \div 2 = 1$ reste $0$ ; $1 \div 2 = 0$ reste $1$. Donc $69 = \textbf{1000101}_2$.\\
Vérification : $64 + 4 + 1 = 69$. Correct.
\item « O » $= 79$ : $79 \div 2 = 39$ reste $1$ ; $39 \div 2 = 19$ reste $1$ ; $19 \div 2 = 9$ reste $1$ ; $9 \div 2 = 4$ reste $1$ ; $4 \div 2 = 2$ reste $0$ ; $2 \div 2 = 1$ reste $0$ ; $1 \div 2 = 0$ reste $1$. Donc $79 = \textbf{1001111}_2$.\\
Vérification : $64 + 8 + 4 + 2 + 1 = 79$. Correct.
\end{itemize}
\item Décodage : $66 = $ B ; $79 = $ O ; $78 = $ N ; $74 = $ J ; $79 = $ O ; $85 = $ U ; $82 = $ R.\\
Le message est : \textbf{BONJOUR}.
\end{enumerate}
\end{corrige}

\begin{corrige}[Exercice 7 — La clé USB de Marlyse]
\begin{enumerate}
\item Espace occupé :
\begin{itemize}
\item Photos : $300 \times 2 = \textbf{600~Mo}$ ;
\item Chansons : $10 \times 5 = \textbf{50~Mo}$.
\end{itemize}
\item Espace total : $600 + 50 = \textbf{650~Mo}$.
\item La clé a une capacité de $1$~Go $= 1024$~Mo. Or $650 < 1024$ : il restera même $1024 - 650 = 374$~Mo libres.\\
\textbf{Conclusion :} la copie est possible, car l'espace total nécessaire ($650$~Mo) est inférieur à la capacité de la clé ($1024$~Mo).
\end{enumerate}
\end{corrige}

\begin{corrige}[Exercice 8 — L'arborescence du lycée]
\begin{enumerate}
\item Arborescence demandée :
\begin{popfigure}
\begin{center}
\begin{tikzpicture}[
  grow=down,
  level 1/.style={sibling distance=4.5cm, level distance=1.1cm},
  level 2/.style={sibling distance=2.2cm, level distance=1.1cm},
  level 3/.style={sibling distance=1.6cm, level distance=1.1cm},
  every node/.style={draw=popblue, fill=popblueL, rounded corners=2pt, font=\small, inner sep=2pt},
  edge from parent/.style={draw=popdark, -}
]
\node {LYCEE}
  child { node {Tle}
    child { node {Tle C}
      child { node {Maths} }
      child { node {Info} }
      child { node {BULLETINS} }
    }
    child { node {Tle D}
      child { node {Maths} }
      child { node {Info} }
      child { node {BULLETINS} }
    }
  };
\end{tikzpicture}
\end{center}
\legende{Arborescence des fichiers du lycée (niveau Terminale, classes C et D).}
\end{popfigure}
\item Chemin d'accès complet : \texttt{LYCEE\textbackslash Tle\textbackslash Tle C\textbackslash BULLETINS\textbackslash bulletin1.pdf}
\item Deux avantages d'une organisation structurée :
\begin{itemize}
\item on \textbf{retrouve rapidement} un fichier grâce à son chemin d'accès, sans chercher parmi des milliers de fichiers mélangés ;
\item on \textbf{évite les confusions} entre fichiers (par exemple les bulletins de deux classes différentes) et on facilite les sauvegardes et le partage.
\end{itemize}
\end{enumerate}
\end{corrige}

\begin{corrige}[Exercice 9 — Type BEPC : conversions binaire--décimal]
\textbf{Partie A.}
\begin{enumerate}
\item Conversion de $25$ : $25 \div 2 = 12$ reste $1$ ; $12 \div 2 = 6$ reste $0$ ; $6 \div 2 = 3$ reste $0$ ; $3 \div 2 = 1$ reste $1$ ; $1 \div 2 = 0$ reste $1$.\\
Lecture des restes de bas en haut : $25 = \textbf{11001}_2$.
\item Conversion de $47$ : $47 \div 2 = 23$ reste $1$ ; $23 \div 2 = 11$ reste $1$ ; $11 \div 2 = 5$ reste $1$ ; $5 \div 2 = 2$ reste $1$ ; $2 \div 2 = 1$ reste $0$ ; $1 \div 2 = 0$ reste $1$.\\
Donc $47 = \textbf{101111}_2$.
\item Vérifications :
\begin{align*}
11001_2 &= 1\times 2^4 + 1\times 2^3 + 0\times 2^2 + 0\times 2^1 + 1\times 2^0 = 16 + 8 + 1 = 25 \quad \checkmark \\
101111_2 &= 2^5 + 2^3 + 2^2 + 2^1 + 2^0 = 32 + 8 + 4 + 2 + 1 = 47 \quad \checkmark
\end{align*}
\end{enumerate}
\textbf{Partie B.}
\begin{enumerate}
\item $10110_2 = 1\times 2^4 + 0\times 2^3 + 1\times 2^2 + 1\times 2^1 + 0\times 2^0 = 16 + 0 + 4 + 2 + 0 = \textbf{22}$.
\item Ce nombre est composé de \textbf{5 bits}. Le bit de poids fort (le plus à gauche, qui vaut $2^4 = 16$) est \textbf{1} ; le bit de poids faible (le plus à droite, qui vaut $2^0 = 1$) est \textbf{0}.
\item \textbf{Conclusion :} l'ordinateur utilise le binaire car ses circuits électroniques ne peuvent distinguer que deux états physiques stables : le courant passe (1) ou ne passe pas (0). Coder l'information avec deux symboles seulement rend la machine simple, rapide et fiable, alors que représenter dix chiffres exigerait de distinguer dix niveaux de tension différents, source d'erreurs.
\end{enumerate}
\textbf{Barème indicatif (10 points) :} A1 : 2 pts ; A2 : 2 pts ; A3 : 2 pts ; B1 : 1,5 pt ; B2 : 1 pt ; B3 : 1,5 pt.\\
\textbf{Points de vigilance :} présenter \emph{toutes} les divisions successives (quotient et reste) ; lire les restes \emph{de bas en haut} ; dans la méthode des puissances de 2, commencer par $2^0$ à droite ; citer explicitement les deux états électriques dans la conclusion.
\end{corrige}

\begin{corrige}[Exercice 10 — Type BEPC : le cybercafé de Bafoussam]
\begin{enumerate}
\item Une \emph{donnée} est un élément brut non traité : par exemple un fichier photo scanné d'un client. Une \emph{information} est une donnée interprétée ou traitée qui a un sens : par exemple la phrase « ce client doit encore payer $500$~F », obtenue à partir de ses factures.
\item Organisation proposée : un dossier principal « CYBERCAFE » contenant trois sous-dossiers « Documents », « Photos » et « Musique », chacun subdivisé par année.
\begin{popfigure}
\begin{center}
\begin{tikzpicture}[
  grow=down,
  level 1/.style={sibling distance=3.4cm, level distance=1.1cm},
  level 2/.style={sibling distance=1.8cm, level distance=1.1cm},
  every node/.style={draw=popgreen, fill=popgreenL, rounded corners=2pt, font=\small, inner sep=2pt},
  edge from parent/.style={draw=popdark, -}
]
\node {CYBERCAFE}
  child { node {Documents}
    child { node {2025} }
    child { node {2026} }
  }
  child { node {Photos}
    child { node {2025} }
    child { node {2026} }
  }
  child { node {Musique}
    child { node {2025} }
    child { node {2026} }
  };
\end{tikzpicture}
\end{center}
\legende{Arborescence proposée pour les fichiers du cybercafé.}
\end{popfigure}
\item Chemin d'accès : \texttt{CYBERCAFE\textbackslash Documents\textbackslash 2026\textbackslash contrat.pdf}
\item Espace total : $2000 \times 3 = \textbf{6000~Mo}$.\\
Conversion : $6000 \div 1024 \approx \textbf{5{,}86~Go}$.
\item \textbf{Démarche :} l'espace nécessaire est d'environ $5{,}86$~Go. Le disque de $4$~Go est insuffisant car $4 < 5{,}86$. Le disque de $10$~Go convient car $10 > 5{,}86$, et il laisse environ $10 - 5{,}86 = 4{,}14$~Go de libre pour les nouveaux fichiers.\\
\textbf{Conclusion :} je conseille à M. Kamdem le disque de $10$~Go, car il peut contenir tous ses fichiers actuels et offre une marge pour l'avenir.
\end{enumerate}
\textbf{Barème indicatif (10 points) :} 1 : 2 pts ; 2 : 2 pts ; 3 : 1 pt ; 4 : 2 pts ; 5 : 3 pts.\\
\textbf{Points de vigilance :} distinguer clairement donnée/information avec des exemples \emph{tirés de la situation} ; l'arborescence doit avoir un dossier racine et des sous-dossiers ; le chemin commence par le dossier principal ; la conversion Mo $\rightarrow$ Go se fait par \emph{division} par $1024$ ; la réponse à la question 5 doit comparer les deux disques et conclure par une phrase rédigée.
\end{corrige}

\begin{corrige}[Exercice 11 — Type BEPC : le bulletin de la 3e A]
\begin{enumerate}
\item Codes ASCII de « CEB » : C $= 67$, E $= 69$, B $= 66$.\\
Conversion de « B » $= 66$ en binaire : $66 \div 2 = 33$ reste $0$ ; $33 \div 2 = 16$ reste $1$ ; $16 \div 2 = 8$ reste $0$ ; $8 \div 2 = 4$ reste $0$ ; $4 \div 2 = 2$ reste $0$ ; $2 \div 2 = 1$ reste $0$ ; $1 \div 2 = 0$ reste $1$.\\
Donc $66 = \textbf{1000010}_2$. Vérification : $64 + 2 = 66$. Correct.
\item Espace total : $60 \times 500 = \textbf{30\,000~Ko}$.\\
Conversion : $30\,000 \div 1024 \approx \textbf{29{,}3~Mo}$.
\item Le CD a une capacité de $700$~Mo. Or $29{,}3 < 700$ : la gravure est \textbf{possible}, et il restera environ $700 - 29{,}3 = 670{,}7$~Mo libres.
\item Arborescence : un dossier « BULLETINS\_3eA » contenant trois sous-dossiers « T1 », « T2 » et « T3 ».\\
Chemin d'accès : \texttt{BULLETINS\_3eA\textbackslash T1\textbackslash ABENA.pdf}
\item Une image scannée est découpée en petits points appelés \cle{pixels}. À chaque pixel, l'ordinateur associe un nombre qui représente sa couleur (ou son niveau de gris). Chacun de ces nombres est ensuite converti en binaire, c'est-à-dire en une suite de 0 et de 1. L'image entière est donc stockée comme une longue suite de bits, que l'ordinateur peut reconstituer à l'écran.
\end{enumerate}
\textbf{Barème indicatif (10 points) :} 1 : 2 pts ; 2 : 2 pts ; 3 : 1,5 pt ; 4 : 2 pts ; 5 : 2,5 pts.\\
\textbf{Points de vigilance :} pour le codage ASCII, donner les trois codes \emph{et} la conversion binaire complète de « B » ; la conversion Ko $\rightarrow$ Mo est une \emph{division} par $1024$ ; la question 3 exige une comparaison chiffrée et une conclusion rédigée ; pour la question 5, les mots-clés attendus sont \emph{pixels}, \emph{couleur}, \emph{nombre}, \emph{binaire}.
\end{corrige}

\newpage

\coursec{Fiche bilan}

\begin{popfigure}
\centering
\begin{tikzpicture}[
  mindmap,
  grow cyclic,
  every node/.style={concept, execute at begin node=\hskip0pt},
  concept color=popblue,
  level 1/.append style={level distance=4.5cm, sibling angle=360/7},
  level 2/.append style={level distance=3cm, sibling angle=45},
  texte court/.style={text width=2.8cm, align=center, font=\small\bfseries},
  branch1/.style={concept color=poporange, texte court},
  branch2/.style={concept color=popgreen, texte court},
  branch3/.style={concept color=poppurple, texte court},
  branch4/.style={concept color=poppink, texte court},
  branch5/.style={concept color=popgold, texte court},
  branch6/.style={concept color=popblue, texte court},
  branch7/.style={concept color=popmuted, texte court},
]
\node[concept, text width=3.2cm, align=center, font=\bfseries\large] {Organisation\\ des données}
  child[branch1] { node[align=center]{Données vs\\ Information} }
  child[branch2] { node[align=center]{Organisation\\ structurée} }
  child[branch3] { node[align=center]{Représentation\\ interne (bit, octet)} }
  child[branch4] { node[align=center]{Codage nombres\\ (Binaire, Hexa)} }
  child[branch5] { node[align=center]{Codage caractères\\ (ASCII, Unicode)} }
  child[branch6] { node[align=center]{Codage images/son\\ (Pixels, Échantillons)} }
  child[branch7] { node[align=center]{Unités de mesure\\ (Ko, Mo, Go, To)} };
\end{tikzpicture}
\legende{Carte mentale : Organisation des données en informatique (Leçon N04)}
\end{popfigure}

\begin{bilan}

\par\medskip\textbf{Définitions clés}\par

\begin{itemize}
  \item \cle{Donnée} : Représentation brute d'un fait, d'une mesure, d'un symbole (ex. : \texttt{23}, \texttt{A}, \texttt{true}). Sans contexte, elle n'a pas de sens.
  \item \cle{Information} : Donnée traitée, contextualisée, qui a un sens pour le destinataire et permet une décision (ex. : « La température est de \SI{23}{\celsius} »).
  \item \cle{Type de données} : Classification qui détermine les valeurs possibles et les opérations autorisées. Principaux types élémentaires : \cle{Entier}, \cle{Réel}, \cle{Booléen}, \cle{Caractère}, \cle{Chaîne de caractères}.
  \item \cle{Structure de données} : Organisation logique des données en mémoire (tableau, liste, pile, fichier, arbre, table de base de données).
  \item \cle{Bit} (Binary Digit) : Plus petite unité d'information, vaut \texttt{0} ou \texttt{1}.
  \item \cle{Octet} (Byte) : Groupe de \textbf{8 bits}. Unité de base de l'adressage mémoire. $1\ \text{octet} = 8\ \text{bits}$.
  \item \cle{Codage} : Convention de représentation d'une information (nombre, texte, image, son) par une suite de bits.
\end{itemize}

\par\medskip\textbf{Organisation structurée des données (Niveaux)}\par

\begin{center}
\begin{tabularx}{\linewidth}{|l|X|l|}
\hline
\rowcolor{popblueL}
\textbf{Niveau} & \textbf{Description} & \textbf{Exemple} \\
\hline
Bit & Unité élémentaire & \texttt{0}, \texttt{1} \\
\hline
Octet & 8 bits & \texttt{01000001} (code ASCII 'A') \\
\hline
Champ / Zone & Donnée élémentaire identifiée (nom + type + valeur) & \texttt{Nom : "KAMGA"} \\
\hline
Enregistrement / Tuple & Ensemble de champs liés à un même objet & Fiche élève : \{\texttt{Matricule}, \texttt{Nom}, \texttt{Classe}\} \\
\hline
Fichier / Table & Ensemble d'enregistrements de même structure & Fichier \texttt{ELEVES.csv} \\
\hline
Base de données & Ensemble de fichiers liés, gérés par un SGBD & Base \texttt{CollegeDB} \\
\hline
\end{tabularx}
\end{center}

\par\medskip\textbf{Représentation des nombres : Conversions à maîtriser}\par

\begin{center}
\begin{tabularx}{\linewidth}{|l|X|}
\hline
\rowcolor{popblueL}
\textbf{Conversion} & \textbf{Méthode express} \\
\hline
Décimal $\to$ Binaire & Divisions successives par 2, lire les restes de bas en haut. \\
\hline
Binaire $\to$ Décimal & Somme des puissances de 2 : $\sum b_i \times 2^i$ (poids de droite à gauche : $1, 2, 4, 8\dots$). \\
\hline
Binaire $\leftrightarrow$ Hexadécimal & Groupes de 4 bits (nibble) $\leftrightarrow$ 1 chiffre hexa ($0\dots9, A\dots F$). \\
\hline
Décimal $\to$ Hexadécimal & Via le binaire ou divisions successives par 16. \\
\hline
Entiers signés & \cle{Complément à deux} : Inverser les bits + ajouter 1. Le bit de poids fort (MSB) indique le signe ($0=+$, $1=-$). \\
\hline
Réels (approche) & Notation scientifique binaire (standard IEEE 754) : Signe $\times$ Mantisse $\times 2^{\text{Exposant}}$. \\
\hline
\end{tabularx}
\end{center}

\par\medskip\textbf{Codage des caractères}\par

\begin{itemize}
  \item \cle{ASCII} (7 bits, 128 codes) : Contrôles (0-31), Chiffres (48-57), Majuscules (65-90), Minuscules (97-122). Étendu à 8 bits (256 codes) pour accents (ISO-8859-1).
  \item \cle{Unicode / UTF-8} : Standard universel. UTF-8 : codage à longueur variable (1 à 4 octets). Compatible ASCII sur 1 octet. \textbf{Indispensable pour le web et le multilinguisme (dont langues camerounaises)}.
\end{itemize}

\par\medskip\textbf{Unités de mesure (Puissances de 2)}\par

\[
\begin{array}{rcl}
1\ \text{Kio (Kibioctet)} &=& 2^{10}\ \text{octets} = 1\,024\ \text{octets} \\
1\ \text{Mio (Mébioctet)} &=& 2^{20}\ \text{octets} = 1\,024\ \text{Kio} \\
1\ \text{Gio (Gibioctet)} &=& 2^{30}\ \text{octets} = 1\,024\ \text{Mio} \\
1\ \text{Tio (Tébioctet)} &=& 2^{40}\ \text{octets} = 1\,024\ \text{Gio}
\end{array}
\]


\par\medskip\textbf{Méthodes express}\par

\end{bilan}

\begin{methode}[Convertir un entier décimal positif en binaire]
\begin{enumerate}
  \item Diviser le nombre par 2.
  \item Noter le reste (0 ou 1).
  \item Recommencer avec le quotient jusqu'à obtenir 0.
  \item Le binaire se lit des restes du \textbf{dernier au premier} (bas vers haut).
\end{enumerate}
\end{methode}

\begin{methode}[Convertir un binaire en décimal]
\begin{enumerate}
  \item Écrire les poids des bits (puissances de 2) au-dessus de chaque bit, en partant de $2^0=1$ à droite.
  \item Additionner les poids correspondant aux bits à \texttt{1}.
\end{enumerate}
\end{methode}

\begin{methode}[Calculer la taille d'un fichier image (bitmap non compressé)]
\begin{enumerate}
  \item Résolution = Largeur $\times$ Hauteur (en pixels).
  \item Profondeur = Nombre de bits par pixel (ex. 24 bits = 3 octets pour RVB Vrai Couleur).
  \item Taille (bits) = Résolution $\times$ Profondeur.
  \item Taille (octets) = Taille (bits) $/ 8$.
\end{enumerate}
\end{methode}


\begin{aretenir}[Checklist avant le Bac]
\begin{itemize}
  \item $\square$ Différencier \textbf{donnée} (brute) et \textbf{information} (contextualisée).
  \item $\square$ Citer les 5 types de données élémentaires et donner un exemple de valeur pour chacun.
  \item $\square$ Expliquer la hiérarchie : Bit $\to$ Octet $\to$ Champ $\to$ Enregistrement $\to$ Fichier $\to$ Base de données.
  \item $\square$ Convertir un entier \textbf{décimal $\leftrightarrow$ binaire} (ex. $42_{10} = 101010_2$).
  \item $\square$ Convertir \textbf{binaire $\leftrightarrow$ hexadécimal} par groupes de 4 bits (ex. $1010_2 = A_{16}$).
  \item $\square$ Représenter un entier \textbf{négatif en complément à deux} sur 8 bits (ex. $-5$).
  \item $\square$ Expliquer la différence entre \textbf{ASCII} (7/8 bits, limité) et \textbf{UTF-8} (variable, universel).
  \item $\square$ Calculer la taille en octets d'un fichier texte (1 octet/caractère ASCII, variable UTF-8) ou image brute (Largeur $\times$ Hauteur $\times$ Profondeur/8).
  \item $\square$ Convertir des unités de stockage : \textbf{Kio, Mio, Gio, Tio} (puissances de 1024) vs \textbf{ko, Mo, Go} (puissances de 1000).
  \item $\square$ Identifier le codage adapté à une situation : nombre entier (binaire/2e complément), réel (IEEE 754), texte (UTF-8), image (RVB + compression), son (échantillonnage).
  \item $\square$ Justifier l'intérêt de l'organisation structurée (accès rapide, intégrité, partage, sécurité).
  \item $\square$ Résoudre un problème concret : « Combien d'octets pour stocker 500 noms de 20 caractères en UTF-8 ? » (Réponse : $\approx 500 \times 20 \times 1 = 10\,000$ octets si ASCII, plus si accents).
\end{itemize}
\end{aretenir}

\begin{attention}[Notation commerciale]
Les constructeurs de disques durs utilisent souvent les puissances de 10 (ko = $10^3$, Mo = $10^6$). L'OS affiche en puissances de 2 (Kio, Mio). D'où la différence affichée (ex. disque \SI{500}{Go} $\approx$ \SI{465}{Gio}).
\end{attention}


\findecours

\end{document}
